% Options for packages loaded elsewhere
\PassOptionsToPackage{unicode}{hyperref}
\PassOptionsToPackage{hyphens}{url}
\PassOptionsToPackage{dvipsnames,svgnames,x11names}{xcolor}
\documentclass[
  11pt,
]{article}
\usepackage{xcolor}
\usepackage[margin=1in]{geometry}
\usepackage{amsmath,amssymb}
\setcounter{secnumdepth}{-\maxdimen} % remove section numbering
\usepackage{iftex}
\ifPDFTeX
  \usepackage[T1]{fontenc}
  \usepackage[utf8]{inputenc}
  \usepackage{textcomp} % provide euro and other symbols
\else % if luatex or xetex
  \usepackage{unicode-math} % this also loads fontspec
  \defaultfontfeatures{Scale=MatchLowercase}
  \defaultfontfeatures[\rmfamily]{Ligatures=TeX,Scale=1}
\fi
\usepackage{lmodern}
\ifPDFTeX\else
  % xetex/luatex font selection
\fi
% Use upquote if available, for straight quotes in verbatim environments
\IfFileExists{upquote.sty}{\usepackage{upquote}}{}
\IfFileExists{microtype.sty}{% use microtype if available
  \usepackage[]{microtype}
  \UseMicrotypeSet[protrusion]{basicmath} % disable protrusion for tt fonts
}{}
\makeatletter
\@ifundefined{KOMAClassName}{% if non-KOMA class
  \IfFileExists{parskip.sty}{%
    \usepackage{parskip}
  }{% else
    \setlength{\parindent}{0pt}
    \setlength{\parskip}{6pt plus 2pt minus 1pt}}
}{% if KOMA class
  \KOMAoptions{parskip=half}}
\makeatother
\usepackage{longtable,booktabs,array}
\newcounter{none} % for unnumbered tables
\usepackage{calc} % for calculating minipage widths
% Correct order of tables after \paragraph or \subparagraph
\usepackage{etoolbox}
\makeatletter
\patchcmd\longtable{\par}{\if@noskipsec\mbox{}\fi\par}{}{}
\makeatother
% Allow footnotes in longtable head/foot
\IfFileExists{footnotehyper.sty}{\usepackage{footnotehyper}}{\usepackage{footnote}}
\makesavenoteenv{longtable}
\setlength{\emergencystretch}{3em} % prevent overfull lines
\providecommand{\tightlist}{%
  \setlength{\itemsep}{0pt}\setlength{\parskip}{0pt}}
\usepackage{bookmark}
\IfFileExists{xurl.sty}{\usepackage{xurl}}{} % add URL line breaks if available
\urlstyle{same}
\hypersetup{
  colorlinks=true,
  linkcolor={Maroon},
  filecolor={Maroon},
  citecolor={Blue},
  urlcolor={Blue},
  pdfcreator={LaTeX via pandoc}}

\usepackage{amssymb}
\usepackage{newunicodechar}
\newunicodechar{°}{\ensuremath{^\circ}}
\newunicodechar{±}{\ensuremath{\pm}}
\newunicodechar{²}{\ensuremath{^{2}}}
\newunicodechar{·}{\ensuremath{\cdot}}
\newunicodechar{¼}{\ensuremath{\frac{1}{4}}}
\newunicodechar{½}{\ensuremath{\frac{1}{2}}}
\newunicodechar{×}{\ensuremath{\times}}
\newunicodechar{÷}{\ensuremath{\div}}
\newunicodechar{ħ}{\ensuremath{\hbar}}
\newunicodechar{̄}{}
\newunicodechar{Δ}{\ensuremath{\Delta}}
\newunicodechar{Λ}{\ensuremath{\Lambda}}
\newunicodechar{Σ}{\ensuremath{\Sigma}}
\newunicodechar{Φ}{\ensuremath{\Phi}}
\newunicodechar{Ω}{\ensuremath{\Omega}}
\newunicodechar{α}{\ensuremath{\alpha}}
\newunicodechar{β}{\ensuremath{\beta}}
\newunicodechar{γ}{\ensuremath{\gamma}}
\newunicodechar{δ}{\ensuremath{\delta}}
\newunicodechar{η}{\ensuremath{\eta}}
\newunicodechar{θ}{\ensuremath{\theta}}
\newunicodechar{λ}{\ensuremath{\lambda}}
\newunicodechar{μ}{\ensuremath{\mu}}
\newunicodechar{ν}{\ensuremath{\nu}}
\newunicodechar{ξ}{\ensuremath{\xi}}
\newunicodechar{π}{\ensuremath{\pi}}
\newunicodechar{ρ}{\ensuremath{\rho}}
\newunicodechar{σ}{\ensuremath{\sigma}}
\newunicodechar{τ}{\ensuremath{\tau}}
\newunicodechar{χ}{\ensuremath{\chi}}
\newunicodechar{ψ}{\ensuremath{\psi}}
\newunicodechar{ᵘ}{\ensuremath{^{\mu}}}
\newunicodechar{ᵤ}{\ensuremath{_{\mu}}}
\newunicodechar{″}{\ensuremath{^{\prime\prime}}}
\newunicodechar{ℂ}{\ensuremath{\mathbb{C}}}
\newunicodechar{ℍ}{\ensuremath{\mathbb{H}}}
\newunicodechar{ℓ}{\ensuremath{\ell}}
\newunicodechar{ℝ}{\ensuremath{\mathbb{R}}}
\newunicodechar{ℤ}{\ensuremath{\mathbb{Z}}}
\newunicodechar{→}{\ensuremath{\to}}
\newunicodechar{∂}{\ensuremath{\partial}}
\newunicodechar{−}{\ensuremath{-}}
\newunicodechar{∓}{\ensuremath{\mp}}
\newunicodechar{√}{\ensuremath{\surd}}
\newunicodechar{∝}{\ensuremath{\propto}}
\newunicodechar{∫}{\ensuremath{\int}}
\newunicodechar{≈}{\ensuremath{\approx}}
\newunicodechar{≥}{\ensuremath{\ge}}
\newunicodechar{≲}{\ensuremath{\lesssim}}
\newunicodechar{≳}{\ensuremath{\gtrsim}}
\newunicodechar{𝕆}{\ensuremath{\mathbb{O}}}
\newunicodechar{⊕}{\ensuremath{\oplus}}
\newunicodechar{⁻}{\ensuremath{^{-}}}
\newunicodechar{¹}{\ensuremath{^{1}}}
\newunicodechar{⁴}{\ensuremath{^{4}}}
\newunicodechar{₀}{\ensuremath{_{0}}}
\newunicodechar{₁}{\ensuremath{_{1}}}
\newunicodechar{₂}{\ensuremath{_{2}}}
\newunicodechar{₃}{\ensuremath{_{3}}}
\newunicodechar{₄}{\ensuremath{_{4}}}
\newunicodechar{₅}{\ensuremath{_{5}}}
\newunicodechar{₆}{\ensuremath{_{6}}}
\newunicodechar{ₖ}{\ensuremath{_{k}}}
\newunicodechar{̂}{}
\newunicodechar{φ}{\ensuremath{\phi}}
\newunicodechar{∮}{\ensuremath{\oint}}
\newunicodechar{‖}{\ensuremath{\Vert}}
\newunicodechar{≡}{\ensuremath{\equiv}}
\newunicodechar{≢}{\ensuremath{\not\equiv}}
\newunicodechar{∎}{\ensuremath{\blacksquare}}
\newunicodechar{𝒯}{\ensuremath{\mathcal{T}}}
\newunicodechar{⁸}{\ensuremath{^{8}}}
\newunicodechar{³}{\ensuremath{^{3}}}
\newunicodechar{⁰}{\ensuremath{^{0}}}
\newunicodechar{≤}{\ensuremath{\le}}
\newunicodechar{ε}{\ensuremath{\epsilon}}
\newunicodechar{Γ}{\ensuremath{\Gamma}}
\newunicodechar{∗}{\ensuremath{*}}
\newunicodechar{′}{\ensuremath{^{\prime}}}
\newunicodechar{ᵀ}{\ensuremath{^{T}}}
\newunicodechar{₊}{\ensuremath{_{+}}}
\newunicodechar{₋}{\ensuremath{_{-}}}
\newunicodechar{⁵}{\ensuremath{^{5}}}
\newunicodechar{⁹}{\ensuremath{^{9}}}
\newunicodechar{⁶}{\ensuremath{^{6}}}
\newunicodechar{⁷}{\ensuremath{^{7}}}
\author{}
\date{}

\begin{document}
\hypertarget{dynamic-relational-transitioning-geometry-a-relational-entanglement-based-account-of-known-physics-with-testable-predictions}{%
\section{Dynamic Relational Transitioning Geometry: a relational, entanglement-based account of known physics, with testable predictions}\label{dynamic-relational-transitioning-geometry-a-relational-entanglement-based-account-of-known-physics-with-testable-predictions}}

\textbf{Matthew Gelfman} (ORCID 0009-0000-6402-2258)\\
\textbf{Melissa Behrends-Gelfman} (ORCID 0009-0001-6281-1947)

Version of 27 November 2026 (v1.4.2: Appendix B added, which defines the particle side on a lattice by extending the construction of Craig (2026), with companion paper F; no result, rule, prediction or count changed. v1.4.1: v1.4.0 with a fuller statement of how AI assistants were used, Acknowledgements; no result changed. v1.4.0: the framework restated as five rules and one bet, read as one matrix of relationships with three roles, Section 2; a short start phase driven by the stiffness of space kept as the alternative start, to be decided by the tilt of the first ripples, Sections 5 and 9; the exact right-angled form of quark mixing added as a dated lead and a charge rule for the far-side weight tested, Section 10; a record of the tests run after 18 November that failed, Appendix A.7, and where each open problem now stands, Appendix A.8. v1.3.3 recalibrated v1.3.2 to the Particle Data Group's 2026 values and JUNO's 2026 splittings, with a full audit of tests set after an earlier attempt failed, Appendix A.6). Built on the work of the physicists credited at the end; this framework is a perspective that aims to show they were each right.

trpgphysics.org \ensuremath{\cdot} November 2026 \ensuremath{\cdot} DOI \href{https://doi.org/10.5281/zenodo.24330763}{10.5281/zenodo.24330763} (all versions)

\begin{quote}
\textbf{Status note, 27 November 2026 (v1.4.2).} Appendix B is new. Craig (2026) has proposed an exact, gauge-invariant lattice definition of the Standard Model's one-handed fermions. With two short lemmas it covers everything on the framework's main line: the far-side states (which come in both handednesses), the hidden partners, the Higgs field, the mass links, the Majorana mass and the majoron. The lemmas are checked numerically. The majoron's symmetry, B \ensuremath{-} L, is exact on the lattice, and baryon number can change only in steps of three, so the lattice agrees with the statement that protons do not decay. Only the main line is covered; the two recorded alternatives are not. All of this is conditional on Craig's theorem, which is a preprint. We rechecked every finite, computable step of his proof independently: 18 of 18 checks pass and no error was found, while his analytic estimates and formal proof files were not checked (B.7). The recheck also showed that his theorem already covers the full 27, so Lemma A is a second, explicit route for the far-side pair rather than a required step. Reflection positivity and the zero-spacing limit stay open, as they are for the Standard Model. No number, rule, prediction, score or count changes. The note of 25 November below is otherwise current.
\end{quote}

\begin{quote}
\textbf{Status note, 25 November 2026 (v1.4.0).} This version simplifies how the framework is stated. It is now five rules and one bet (Section 2): (1) everything is one matrix of relationships, read one perspective at a time, in which held states (matter) sit on the diagonal, the links between them are the pulls (the twelve carriers), the link across one state between its two sides is the Higgs (the resistance that gives weight), and the geometry of the relationships is space and time; (2) whatever holds together is a pull balanced by a resistance, and mass is that balance; (3) maximum uncertainty between self and relationship (the 45\ensuremath{^\circ} rule, leptons only); (4) nothing extra; (5) the universe started smooth. The bet is unchanged: the far side of the relationship matrix is real. How the start ran is left to measurement. A short stretching driven by the stiffness of space (Starobinsky, 1980) would explain the measured flatness and the uniform temperature of the sky and predicts a tilt of the first ripples n\textsubscript{s} = 0.963 with primordial gravitational waves at r = 0.0040; it agrees with the CMB alone but sits 3.2--4.2\ensuremath{\sigma} below the tilt found when the CMB is combined with DESI's baryon acoustic oscillations, past the 3\ensuremath{\sigma} line used throughout this paper. The main line therefore follows the combined data, as it does for dark energy, whose reading here allows the change DESI hints at: it does not adopt that stretching phase, leaves the form of the start open as in v1.3.3, keeps flatness and the tilt as measured inputs, and keeps the start phase on record as the alternative, which returns if the tilt settles near 0.963 or r \ensuremath{\approx} 0.004 is seen (Sections 5 and 9). No particle-physics number changes. On the main line the count stays 32 measured numbers today and 28 if every pending measurement agrees, with six choices made with today's data in view; with the alternative start it is 33 and 28, with a seventh choice, the Higgs field's coupling to curvature (Appendix A.2). Also new: the exact right-angled quark triangle as a dated lead and a tested charge rule for the far-side weight, which does not hold (Section 10); Appendix A.7 records the other tests run since 18 November, all of which failed, and Appendix A.8 gives where each open problem stands. What changes is how much is assumed: the rules of quantum theory and the stage of space and time are no longer borrowed whole, and about 25 separate starting assumptions reduce to six. The note of 18 November below is kept as it was; where it differs, this note is current.
\end{quote}

\begin{quote}
\textbf{Status note, 18 November 2026.} This draft is current to 18 November 2026. The framework is stated as three ideas, one rule and one bet (Section 2), and takes about 28 measured numbers, against about 34 for the Standard Model with gravity and cosmology counted the same way (the 28 assumes every pending measurement agrees; counting only savings already confirmed it is 32, plus six stated choices; Appendix A); the Higgs mass follows from a rule at the grain of space that awaits measurement (Section 9). Seven of its nine criteria (listed in Appendix A.0) are met and two are partly met: the quark link strengths, which await a rule, and the force groups, where only the force strengths remain and await a measurement of the strong coupling. JUNO's first precision results (2026) favour normal ordering at 2.1\ensuremath{\sigma} and match the predicted small neutrino splitting, but put the TM1 mixing value 2.3\ensuremath{\sigma} above the measured sin\textsuperscript{2}\ensuremath{\theta}\textsubscript{12} (Section 9). Results that follow within the framework but rest on a measurement not yet returned are marked \emph{awaiting measurement} and are not counted as established until that measurement returns. Dated predictions: the strong coupling 0.1178, with the grain of space fixed in advance as the Planck energy divided by \ensuremath{\surd}2; the quark triangle angle \ensuremath{\gamma} = 64.1--67.6\ensuremath{^\circ}; the first-row quark-mixing shortfall shrinking to nothing (Section 9); the neutrino CP-violating phase \ensuremath{\delta} = 253--292\ensuremath{^\circ}; a neutrinoless double-beta mass of 1.3--4.1 meV; the tau mass 1776.969 MeV; the neutrino mass sum 0.0591 to 0.0594 eV with normal ordering, given today's limit from cosmology; and a quark pairing gap of about 82--97 MeV (for a strange-quark mass of 100 MeV; 88--109 MeV for a running mass of 120--140 MeV) in the cores of compact stars, the band that keeps the heaviest stars inside the observed window of 2.0--2.15 solar masses (the observations set this band, so the stars are not evidence for it; Appendix A.6). If dark matter is the cold majoron wave there is no X-ray line and no measurable primordial gravitational waves; both belong to that lead and not to the core. The heavy quark link strengths and mixing, the matter excess and dark matter are leads, each with one named missing rule; the twist is a measured input. The majoron's mass is not fixed; it follows m\_J = 0.027 eV \ensuremath{\times} (1.5 \ensuremath{\times} 10\textsuperscript{12} GeV \ensuremath{\div} f)\textsuperscript{4} \ensuremath{\div} angle\textsuperscript{4}. Where older passages below differ, this note and the Results page at trpgphysics.org are current. Appendix A gathers the bookkeeping in one place: every quantity labelled as input, fitted, derived or predicted; the input count as a range (32 counting only savings already confirmed, 28 if every pending measurement agrees, plus six choices made with today's data in view); and the measurements that would show the framework false.
\end{quote}

\hypertarget{abstract}{%
\subsection{Abstract}\label{abstract}}

We present Dynamic Relational Transitioning Geometry (TRPG), a framework in which the universe is described by the relationships between its parts rather than by the parts alone. It rests on five rules and one bet. (1) Everything is one matrix of relationships, seen one perspective at a time, with entanglement first: matter is held states on its diagonal, forces are the links between them, the Higgs is the link across one state, and space and time are the geometry of the relationships. (2) Every bound structure is a pull balanced by a resistance (containment), and mass is that balance. (3) Within the lepton families a particle's tie to itself balances its ties to the others. (4) Nothing extra. (5) The universe started smooth; whether a short stretching followed is left to measurement (Section 5). The bet is that the far side of the relationship matrix is real. Section 2 unpacks these into the working principles; earlier versions stated rules 1--4 as three ideas and one rule.

Combined with established physics, these principles reproduce the following.

\begin{itemize}
\tightlist
\item
  Einstein's gravity, through the entanglement area law and Jacobson's thermodynamic argument. This includes the full deflection of starlight (1.75\ensuremath{^{\prime\prime}}) and the Shapiro delay.
\item
  The expansion of a flat universe, with an age of 13.79 billion years; with the alternative start phase, also the flatness itself and the uniform temperature of the sky.
\item
  A containment pressure B = 60 (57 to 64) MeV/fm\textsuperscript{3}, obtained from the proton and \ensuremath{\Delta} masses alone.
\item
  Carried to compact stars (observed as neutron stars, modelled here as strange quark stars), that same pressure gives maximum masses of 2.0 to 2.15 solar masses and a tidal deformability \ensuremath{\Lambda}\textsubscript{1.4} = 440 to 580, consistent with observation, provided quarks in the core pair with a gap of about 82 to 97 MeV. Without pairing the stars fail; the gap band is the one the observed stars require, so the star masses are not evidence, and the gap is the prediction for an independent determination.
\item
  The three charged-lepton masses as the eigenvalues of a single relational matrix with two numbers, a scale and a twist, under the 45\ensuremath{^\circ} balance.
\item
  The strong force SU(3) with three colours, once the medium's inner space is required to have no gaps and to conserve the size of relationships; and, reading the forces from the relationship matrix seen from one perspective, with stated readings, which other forces exist and, reading the matrix's states three at a time, every particle of a generation with its measured charge and one-handed matter. The force strengths are taken from measurement.
\item
  Fractional quark charges, which anomaly cancellation fixes once the number of colours is given.
\end{itemize}

The framework makes predictions that can fail and that experiments can test. The first is normal neutrino mass ordering (given today's limit from cosmology), with a lightest mass of 0.36 meV and a sum of 0.0591 to 0.0594 eV. The second is a quark pairing gap of about 82 to 97 MeV in the cores of compact stars (for a strange-quark mass of 100 MeV; what the stars fix is 4\ensuremath{\Delta}\textsuperscript{2} \ensuremath{-} m\textsubscript{s}\textsuperscript{2}). A third is a strong coupling of 0.1178, on which the meeting of the force strengths rests and which awaits measurement. A fourth, beyond the reach of present and planned experiments, is that neutrinos are their own antiparticles, with neutrinoless double beta decay at an effective mass of 1.3 to 4.1 meV. Tested strictly, two quark-mass relations in the particle matrices pass: the down quark through the Cabibbo angle and the up quark through the decay \ensuremath{\eta} \ensuremath{\to} 3\ensuremath{\pi}, without lattice input. The force strengths together with gravity's strength, the top, bottom, charm and strange masses with quark mixing, the lepton twist (0.222222, close to 2/9 but measured, and the same open number as the electron's mass) remain open, as do the mass and amount of dark matter (the matrix allows one candidate, a chargeless wave, the majoron, as a lead) and the size and tilt of the seeds of galaxies, which are measured; the alternative start phase would predict the tilt, n\textsubscript{s} = 0.963, and primordial gravitational waves at r = 0.0040 \ensuremath{\pm} 0.0003, and is favoured only if the tilt settles near the CMB-only value. The matter--antimatter excess is not predicted; it is allowed. Where a result was already established by others (the W and Z mass relation, atomic binding, lepton lifetimes, the Cabibbo angle, hadron and nuclear masses), the framework reaches it by a second route, and each such check is reported with its status: confirmed, fitted, or failed.

\hypertarget{introduction}{%
\subsection{1. Introduction}\label{introduction}}

Physics has several very successful descriptions of nature. General relativity describes gravity as geometry. Quantum mechanics describes matter as waves of probability. Quantum field theory describes the forces. Cosmology describes the history of the whole. Each is confirmed to remarkable precision, and yet they are usually presented as separate subjects.

This work starts from a different question. What if the great physicists were each correct, and each was looking at the same universe from a different side? Then the task is not to replace any of them. It is to find the single picture in which each of their results is one view.

The picture proposed here is relational. A particle is real on its own, but what it does, and what the universe does, comes from how particles relate to one another. Write down how every particle relates to every other and you obtain a matrix. The structure of physics is the structure of such matrices.

\textbf{Purpose and method.} This work starts from existing conclusions. Any idea that disagreed with determined measurements was discarded, and the record of each one is kept. Its purpose is to show that the Standard Model, gravity and quantum mechanics already fit together, building on the work of those who found them. Where a choice was made with today's data in view it is labelled as fitted, not predicted. Its predictions are only for measurements that have not yet been made.

Section 2 states the principles. Section 3 writes them as two actions, a fundamental one and an effective one for hadrons. Sections 4 to 8 show what follows, each compared with measurement. Section 9 gives the predictions, Section 10 the open questions, and Section 11 concludes.

\hypertarget{principles}{%
\subsection{2. Principles}\label{principles}}

\textbf{The statement of v1.4.0: five rules and one bet.} The framework is one matrix of relationships read in three roles. \emph{Held states} sit on the diagonal: the 16 states of a generation, placed three times. \emph{Links between different states} are the pulls: 8 colour links, 3 weak links and 1 hypercharge phase, the twelve carriers. \emph{The link across one state}, from its left-handed side to its right-handed side, is the Higgs: the resistance that gives weight, which is also the perspective rule of handedness (Section 8). The strengths of these links (the masses and mixings) are where every open number of the framework lives. This layout is the one used to compute the proton on a lattice (Wilson, 1974): matter on the points, forces on the links; relabelling each point's colours independently leaves every term unchanged (\texttt{analysis/forces\_resistances\_matrix\_2026-11-19.py}; \texttt{figures/relationship\_matrix\_diagonal\_links.png}). It is a reading of the same equations and changes no number.

\begin{enumerate}
\def\labelenumi{\arabic{enumi}.}
\tightlist
\item
  \textbf{One matrix of relationships, seen one perspective at a time, with entanglement first.} The geometry of the relationships is the stage: four numbers in a 2 \ensuremath{\times} 2 grid of complex relationships have the determinant t\textsuperscript{2} \ensuremath{-} x\textsuperscript{2} \ensuremath{-} y\textsuperscript{2} \ensuremath{-} z\textsuperscript{2}, unchanged by any change of perspective, which is 3 + 1 space and time with light speed built in (Penrose and Rindler, 1984); quaternions would give 5 + 1 and octonions 9 + 1. Gravity follows from entanglement and heat flow (Jacobson, 1995), and space can be built from entanglement in model universes (Van Raamsdonk, 2010), though this is not yet shown for ours. The rules of quantum theory reduce to two: complex waves, now required by experiment (Renou et al., 2021), and entanglement first, from which quantum theory has been derived (Hardy, 2001; Chiribella, D'Ariano and Perinotti, 2011). Uncertainty is what any wave does; the probability rule follows from the space of states (Gleason, 1957); spin--statistics and CPT are theorems. Time can be read from relationships, as change seen by any part entangled with a clock (Page and Wootters, 1983). Charges are which links touch a state (Noether, 1918).
\item
  \textbf{Balance.} Whatever holds together is a pull balanced by a resistance, and mass is the balance point. The weak force is the electric family of pull cut short by the Higgs resistance (Anderson, 1963); the strong force is a pull whose carriers resist each other, squeezing the field into a tube, and its resistance scale follows from its strength. The same balance sets the proton's containment pressure and, on the largest scales, the radius at which a galaxy group's pull equals the tension of dark energy (about 1.6 Mpc for the Local Group, just outside its observed edge of about 1 Mpc).
\item
  \textbf{Maximum uncertainty between self and relationship} (the 45\ensuremath{^\circ} rule, leptons only; Section 10).
\item
  \textbf{Nothing extra.}
\item
  \textbf{Smooth start.} The universe began smooth (Penrose, 1979), which sets the direction of time. Whether the stiffness of space then drove a short stretching (Starobinsky, 1980) is decided by measurement; the main line follows today's combined data and does not take it (Section 5).
\end{enumerate}

\textbf{The bet} is unchanged: the far side of the relationship matrix is real.

Of these, rules 1 and 2 restate earlier ideas 1 and 2, rules 3 and 4 are unchanged, and rule 5 is new; its stretching phase is the alternative, not the main line. What is assumed without explanation falls from the earlier rules plus the borrowed rules of quantum theory, the borrowed stage of general relativity and about nine start facts, to five rules and one bet. Measured numbers are not reduced by this restatement (Appendix A.2). Not testable as such: the matrix layout and the pull-and-resistance reading (they restate the same equations), space and time from relationships in our universe (not yet), and why the start was smooth. The earlier statement follows, unchanged in content.

\textbf{The earlier statement (v1.3.3).} Stated most simply, the framework is three ideas, one rule and one bet. \textbf{The ideas.} (1) Space is one full medium of relationships, seen one perspective at a time, with entanglement first. (2) Every structure is a balance of outward push and inward pull (containment). (3) Within the lepton families, a particle's tie to itself and its ties to the others are balanced (45\ensuremath{^\circ}); for leptons this follows from the family symmetry together with maximum uncertainty between self and relationship (Section 10), and it does not hold for quarks (Section 7). \textbf{The rule.} Nothing extra: nothing beyond the 27 states of each generation's matrix (the 16 known, the 10 of the far side and one neutral state), the twelve carriers and one Higgs pair. So no proton decay, no fourth generation, no new forces, no mirror families and no second Higgs pair. \textbf{The bet.} The far side of the relationship matrix is real (with two stated assumptions, Section 10); it gives the strong coupling and the Higgs-mass rule of Section 9, and if it fails the ideas and the rule still stand. Which forces exist follows from the perspective principle by a theorem (Section 8). A whole generation's charges and one-handed matter are reproduced with one further reading of the matrix, which is stated and not derived. The full medium is consistent with the measured flat universe but does not by itself require it. The principles below unpack the three ideas.

\begin{enumerate}
\def\labelenumi{\arabic{enumi}.}
\tightlist
\item
  \textbf{One medium.} Space is like liquid filling a sealed bag. There is no gap between the liquid and anything in it, because the liquid is space itself.
\item
  \textbf{Always full.} The proportions of matter, radiation and dark energy change, but together they always fill the whole. Energy is conserved, though conservation can be checked only locally.
\item
  \textbf{Resistance with direction.} Space has a baseline resistance. Mass and energy raise it, equally for time and for space. Waves are energy spreading through regions of different resistance.
\item
  \textbf{Containment.} Every bound structure is a balance of outward push against inward pull. If the push wins, the structure leaks, as a star does. If the pull wins, it holds everything in, as a black hole does. Inside a proton, quarks push out and gluons pull in. The rule is the same at every size, while the forces differ.
\item
  \textbf{Relational scale.} Each particle sets its own scale of view, and no scale is privileged by the laws themselves. The proton sits at the scale of the spectrum we physically observe, so it sets our view. The only scale-fair distance between two masses is S = \ensuremath{\frac{1}{2}} ln(m\textsubscript{1}/m\textsubscript{2}).
\item
  \textbf{Entanglement first.} Entanglement is the defining feature of quantum physics, as Schrödinger said in 1935. Entangled partners are joined through the geometry, so nothing needs to cross the distance between them. Gravity follows from entanglement, and uncertainty means that we see one perspective of a single geometry at a time.
\item
  \textbf{Transitioning geometry.} Space behaves like crossing waves. Locally it curves as peaks, saddles and troughs, and the pattern keeps changing, but its total curvature averages to zero.
\item
  \textbf{Smooth start (v1.4.0).} The universe began smooth. A short stretching driven by the stiffness of space against bending is kept as the alternative start (Section 5).
\end{enumerate}

\hypertarget{the-action}{%
\subsection{3. The action}\label{the-action}}

The framework is written as two actions: a fundamental one, valid from the grain of space down, and an effective one for the inside of hadrons and compact stars. Units are natural (\ensuremath{\hbar} = c = 1) and the metric signature is (+, \ensuremath{-}, \ensuremath{-}, \ensuremath{-}). M\textsubscript{Pl} = G\textsuperscript{\ensuremath{-}1/2} = 1.22 \ensuremath{\times} 10\textsuperscript{19} GeV is the Planck energy, \ensuremath{\bar{M}}\textsubscript{P} = (8\ensuremath{\pi}G)\textsuperscript{\ensuremath{-}1/2} = 2.44 \ensuremath{\times} 10\textsuperscript{18} GeV the reduced Planck mass, and the grain of space is E\textsubscript{grain} = M\textsubscript{Pl}/\ensuremath{\surd}2 = 8.63 \ensuremath{\times} 10\textsuperscript{18} GeV (Section 9). Every term has mass dimension four.

S = \ensuremath{\int} d\textsuperscript{4}x \ensuremath{\surd}(\ensuremath{-}g) {[} \ensuremath{\frac{1}{2}} \ensuremath{\bar{M}}\textsubscript{P}\textsuperscript{2} R \ensuremath{-} \ensuremath{\frac{1}{4}} \ensuremath{\Sigma} F\textsuperscript{2} + \ensuremath{\Sigma} \ensuremath{\bar{\ensuremath{\psi}}} i\ensuremath{\gamma}\ensuremath{^{\mu}}D\ensuremath{_{\mu}} \ensuremath{\psi} + \textbar DH\textbar{}\textsuperscript{2} \ensuremath{-} V(H) \ensuremath{-} ( (\ensuremath{\surd}2/v\textsubscript{H}) \ensuremath{\bar{L}}\textsubscript{L} S\textsubscript{\ensuremath{\ell}}\textsuperscript{2} e\textsubscript{R} H + quark terms + far-side terms + h.c. ) \ensuremath{-} \ensuremath{\frac{1}{2}} (M\textsubscript{N} \ensuremath{\bar{N}}\textsuperscript{c}N + h.c.) {]}

S\textsubscript{eff} = \ensuremath{\int} d\textsuperscript{4}x \ensuremath{\surd}(\ensuremath{-}g) {[} \ensuremath{\bar{q}} i\ensuremath{\gamma}\ensuremath{^{\mu}}\ensuremath{\partial}\ensuremath{_{\mu}} q + \ensuremath{\frac{1}{2}}(\ensuremath{\partial}P)\textsuperscript{2} \ensuremath{-} (\ensuremath{\lambda}/4)(P\textsuperscript{2} \ensuremath{-} v\textsuperscript{2})\textsuperscript{2} \ensuremath{-} g\textsubscript{P} P \ensuremath{\bar{q}}q {]}

\textbf{Gravity (R).} The curvature term is not put in by hand. It emerges from entanglement, as shown in Section 4. The entanglement route (Jacobson, 1995) yields only the trace-free part of Einstein's equations, with the cosmological constant appearing as a constant of integration; the metric is therefore varied at fixed volume element (unimodular gravity: Einstein, 1919; van der Bij, van Dam and Ng, 1982; Henneaux and Teitelboim, 1989; Weinberg, 1989). In this form the constant part of V(H) does not curve space (Section 5). Changes of V(H) with time, such as its drop at the electroweak crossover near 160 GeV, still gravitate in the ordinary way; at that time they are a small part of the energy, which radiation dominates.

\textbf{Containment (P).} P is an effective field of the strong force, in the manner of Friedberg and Lee (1977): it describes the gluon vacuum inside hadrons and compact stars and is not added to the fundamental action, so the strong force is not counted twice. It has mass dimension one, so g\textsubscript{P} is a pure number. The depth of its potential, \ensuremath{\lambda}v\textsuperscript{4}/4, is the containment pressure B. Its stiffness, \ensuremath{\surd}(2\ensuremath{\lambda}) v, is the lightest glueball, whose mass from lattice calculations is 1.71 to 1.73 GeV. With B from the proton (Section 6), \ensuremath{\lambda} \ensuremath{\approx} 1160 to 1210 and v \ensuremath{\approx} 35 MeV. No parameter is tuned. A quartic coupling this large is far beyond the range where loop corrections are small (about 8\ensuremath{\pi}/3 for unitarity, 16\ensuremath{\pi}\textsuperscript{2} at most), so S\textsubscript{eff} is used only as a classical, mean-field description of the bag, as Friedberg and Lee used it, and not as a weakly coupled quantum field theory.

\textbf{Leptons (S\textsubscript{\ensuremath{\ell}}).} S\textsubscript{\ensuremath{\ell}} = M\textsubscript{0} \ensuremath{\cdot} circulant(1, e\textsuperscript{i\ensuremath{\delta}}/\ensuremath{\surd}2, e\textsuperscript{\ensuremath{-}i\ensuremath{\delta}}/\ensuremath{\surd}2), whose eigenvalues are the square roots of the charged-lepton masses (so M\textsubscript{0} carries units of \ensuremath{\surd}mass). The term pairs the left-handed doublet L\textsubscript{L} with the right-handed singlet e\textsubscript{R}, as gauge symmetry requires, and S\textsubscript{\ensuremath{\ell}}\textsuperscript{2}/v\textsubscript{H} is a pure number. The matrix has two numbers, the scale M\textsubscript{0} and the twist \ensuremath{\delta}, with links of 1/\ensuremath{\surd}2 (Section 7). It is a fixed relational matrix, not a field: the action records the pattern and does not supply a mechanism that produces it.

\textbf{Far side and neutrinos.} Each generation's 27 holds, besides the 16 known states, the 10 of the far side and one neutral state. The far-side 10 are spin-\ensuremath{\frac{1}{2}} matter: a colourless weak pair and a coloured pair of charge \ensuremath{\mp}1/3 (the D states), each with both handednesses, so they take ordinary masses M\textsubscript{L} and M\textsubscript{D} (Section 10). The Higgs doublet H is the separate spin-0 copy of the colourless part of the same numbers, as in the standard E\textsubscript{6} construction; its coloured spin-0 parts are absent (Section 8). M\textsubscript{N} is the Majorana mass of the hidden neutrino partners, written with N as a right-handed Weyl field and N\textsuperscript{c} its conjugate.

\textbf{What earlier drafts had.} A scale field \ensuremath{\sigma}, which coupled to nothing, is dropped: the relational scale S = \ensuremath{\frac{1}{2}} ln(m\textsubscript{1}/m\textsubscript{2}) is a coordinate on mass ratios, not a field.

\textbf{Inputs.} Put together as one calculation (Section 10), the framework takes about 28 measured numbers, against about 34 for the Standard Model with gravity and cosmology counted the same way (the 28 assumes every pending measurement agrees; counting only savings already confirmed it is 32, plus six stated choices; Appendix A). They are: the three force strengths; the Higgs value; the charged-lepton scale and twist; five quark masses (the down quark follows from the Cabibbo angle); four quark-mixing numbers and the strong-interaction angle; two neutrino splittings; the four neutrino mixing numbers that the mixing pattern does not fix; gravity's strength; the amount of dark energy; the amount of dark matter; the matter excess; and the two numbers for the seeds of galaxies (their size and tilt). Saved against the standard list and confirmed by measurement are the tau mass and the down quark mass. The lightest neutrino mass, the Higgs mass (Section 9) and two neutrino mixing numbers are saved only if pending measurements agree, so the figure of about 28 holds only if they all do, and counting only confirmed savings it is 32 (Appendix A); and one force strength is saved as well if the strong coupling lands at 0.1178. The proton's two containment numbers are effective numbers of the strong force, like the constants of the nuclear-mass formula, and are counted on neither side. The up quark's mass is counted as measured: the \ensuremath{\eta} \ensuremath{\to} 3\ensuremath{\pi} route trades it for the measured ratio Q, so it is a check and not a saving. The numbers attached to the leads (the majoron's mass and the start conditions of the early universe) belong to those leads and are not part of the core.

\hypertarget{gravity}{%
\subsection{4. Gravity}\label{gravity}}

\textbf{Resistance and the bending of light.} Give space a resistance n = 1 + 2GM/(rc\textsuperscript{2}) around a mass M. Light passing the Sun is then deflected by 1.75\ensuremath{^{\prime\prime}}, as measured. The radar echo from planets is delayed by exactly the Shapiro amount. The factor of two matters. A resistance acting on time alone would give half the deflection (0.876\ensuremath{^{\prime\prime}}), and one with no direction would give none. The resistance must act on time and space equally. This is the optical form of weak-field general relativity (Eddington, Gordon). Feynman, Weinberg and Deser showed that a long-range field which couples to everything, and which has this tensor character, must be Einstein's gravity.

\textbf{Gravity from entanglement.} In a filled medium, the entanglement between a region and its surroundings grows with the area of the boundary rather than with the volume inside. We confirmed this on a lattice of coupled oscillators: a boundary law fits the computed entanglement far better than a volume law. Jacobson (1995) showed that if this area law holds and the relation \ensuremath{\delta}Q = T dS holds at every local horizon, Einstein's equations follow. Gravity is then a consequence of entanglement, not an extra ingredient. Related work by Ryu and Takayanagi, and by Van Raamsdonk, connects entanglement directly to the shape of space.

\textbf{Gravity's strength.} Entanglement fixes the form of gravity (Bekenstein, 1973; Hawking, 1975; Jacobson, 1995) but not its strength, which is the fineness of the links of space: ln(E\textsubscript{grain}/m\textsubscript{p}) = 43.7, with E\textsubscript{grain} = M\textsubscript{Pl}/\ensuremath{\surd}2 (Section 3). It is tied to the force strengths, since the proton is light compared with the grain because the strong force is weak there and strengthens only slowly (Gross and Wilczek, 1973; Politzer, 1973; Wilczek, 2001); the simple count 2\ensuremath{\pi}/(7\ensuremath{\alpha}\textsubscript{s}) at the grain gives 47.5 with the Standard Model's states alone and 43.6 with the far side included. The second is close, but the count uses the six-quark rate of running all the way down while the actual rate changes where the far-side states enter, so the closeness is not counted as evidence. The tie trades one measured number for another. Setting gravity equal to the other forces where they meet predicts a grain energy 62 times too low (Kaplunovsky, 1988), and counting the framework's 191 states (124 for the Standard Model with three hidden partners, 66 for the far side and one majoron; Sakharov, 1967) does not give the size. The Higgs field cannot set gravity's strength (Zee, 1979; the required curvature coupling exceeds the LHC limit of Atkins and Calmet, 2013, by 4 \ensuremath{\times} 10\textsuperscript{16}); conversely, a Higgs self-coupling vanishing at the grain (Shaposhnikov and Wetterich, 2010) needs m\textsubscript{H} = 128.8 \ensuremath{\pm} 1.5 GeV at two loops with the Standard Model's states alone against 125.13 (Buttazzo et al., 2013; Particle Data Group, 2026), a strict fail at 2.6\ensuremath{\sigma} (Section 9).

\textbf{Black holes.} For the remnant of the first observed merger, GW150914, the ringdown is 250 Hz with a damping time of 4.0 ms. This agrees with the Kerr solution and with the measurement. A black hole is fully contained, so no echoes are expected, and none have been seen.

\hypertarget{cosmology}{%
\subsection{5. Cosmology}\label{cosmology}}

\textbf{The full bag.} The density parameters are measured to sum to one, \ensuremath{\Omega}\textsubscript{m} + \ensuremath{\Omega}\textsubscript{r} + \ensuremath{\Omega}\textsubscript{\ensuremath{\Lambda}} = 1 to within 0.2\% (\ensuremath{\Omega}\textsubscript{k} = 0.0007 \ensuremath{\pm} 0.0019; Planck 2018 with BAO): on the largest scales the universe is spatially flat. A full medium is consistent with this but does not require it, since closed and open universes are also full everywhere; flatness is taken from measurement. With the alternative start phase (below) it would also be explained: the stretching leaves \textbar \ensuremath{\Omega}\textsubscript{k}\textbar{} near 10\textsuperscript{\ensuremath{-}5} or less. Solving the Friedmann equation with Planck 2018 parameters gives an age of 13.79 billion years (measured 13.787 \ensuremath{\pm} 0.020).

\textbf{No centre.} The Big Bang is the release of an over-contained singularity. It is not an explosion into empty space: the same amount is everywhere, spreading as waves. Because there is no empty space outside the bag, the ``centre'' is a moment in time, not a place. Every observer sees the same background glow.

\textbf{The start (v1.4.0): main line and alternative.} The main line follows today's combined data. With the CMB alone the tilt of the first ripples is 0.9638 \ensuremath{\pm} 0.0040; combined with DESI's baryon acoustic oscillations it is 0.969--0.974 \ensuremath{\pm} 0.003, and the framework's dark-energy reading allows the change in dark energy that DESI hints at (below), so the main line takes that direction: it does not adopt this stretching phase and leaves the form of the start open, as in v1.3.3, and flatness, the size and the tilt of the first ripples stay measured inputs. The alternative, kept on record and decided by the tilt, is a stretching phase. Space resists being curved; the simplest form of that stiffness is a term in R\textsuperscript{2} in the gravitational action (Starobinsky, 1980), which adds no particle. In this alternative, in the first instant it stretches space for about 54--55 e-folds, then hands its energy to the Higgs field as heat. The measured size of the first ripples, A\textsubscript{s} = 2.1 \ensuremath{\times} 10\textsuperscript{\ensuremath{-}9}, fixes its one number, the scalaron mass of 3.1--3.2 \ensuremath{\times} 10\textsuperscript{13} GeV; then the tilt of the ripples, n\textsubscript{s} = 1 \ensuremath{-} 2/N = 0.9630--0.9636, and the ratio of gravitational waves to ripples, r = 12/N\textsuperscript{2} = 0.0040 \ensuremath{\pm} 0.0003, follow (Mukhanov and Chibisov, 1981; Bezrukov and Gorbunov, 2012), together with flatness, a uniform temperature, no separate dark-matter ripples and nearly Gaussian seeds. The expansion rate during the start is 1.6 \ensuremath{\times} 10\textsuperscript{13} GeV. The heat afterwards is about 5 \ensuremath{\times} 10\textsuperscript{9} GeV, or 0.5--2 \ensuremath{\times} 10\textsuperscript{9} GeV if the Higgs couples to curvature with \ensuremath{\xi} \ensuremath{\approx} 0.10--0.15 (Vilenkin, 1985; Bezrukov and Gorbunov, 2012). With the measured top mass the Higgs self-coupling turns negative far below the start's expansion rate, so the Higgs is safe during the start only if \ensuremath{\xi} \ensuremath{\gtrsim} 0.06 (Herranen et al., 2014; Espinosa et al., 2015); \ensuremath{\xi} is the seventh choice made with today's data in view (Appendix A.2). Consequences for the framework's own pieces: the coloured far-side states (4.6 \ensuremath{\times} 10\textsuperscript{13} GeV) are heavier than the scalaron and are never made; the far-side weak pairs (2.4 \ensuremath{\times} 10\textsuperscript{8} GeV) are made and must decay (Section 10); the majoron's starting angle is averaged, because the expansion rate exceeds its settling value (Section 8). Why the start was smooth is not testable. The tilt and r are tested in Section 9. This alternative makes two dated predictions (Section 9); it returns if a CMB-only tilt or a reconciled CMB--DESI tilt settles near 0.963, or if primordial gravitational waves near r = 0.004 are seen.

\textbf{Expansion as tension.} Dark energy has a pressure equal to minus its energy density, which is tension: a pull spread everywhere, like a stretched rubber band, not a pull from an edge (there is no edge or centre). In gravity tension drives faster expansion. With \ensuremath{\Omega}\textsubscript{\ensuremath{\Lambda}} = 0.685 against \ensuremath{\Omega}\textsubscript{m} = 0.315 the two are equal at redshift 0.30 and the speed-up begins at 0.63, as measured. Bound systems hold where pull exceeds tension: about 110 pc for the Sun and 1.6 Mpc for the Local Group, just outside the observed radius of about 1 Mpc inside which the group has stopped expanding (Karachentsev et al., 2009).

\textbf{Flat on average, curved locally.} The simplest crossing-wave surface, z = cos x + cos y, has one peak, two saddles and one trough per cell. By the Gauss--Bonnet theorem its total curvature over a cell is exactly zero. This shows how local curvature can keep changing while averaging to zero; it is a picture on a two-dimensional surface and does not by itself fix the curvature of three-dimensional space.

\textbf{Dark energy and dark matter.} With gravity from entanglement, Einstein's equations come out in their trace-free form (Section 3), so the constant part of stored field energy does not gravitate as vacuum energy, and the cosmological constant is a constant of integration fixed once (Einstein, 1919; Jacobson, 1995). Whether this protection survives quantum corrections is debated, and it is not claimed here. Otherwise the Higgs balance energy, 5 \ensuremath{\times} 10\textsuperscript{54} times the observed value, would dominate (Weinberg, 1989). On the main line the size of dark energy is set by the two ends of space together, the grain and the horizon: \ensuremath{\surd}(grain \ensuremath{\times} horizon) = 3.5 meV against 2.24 meV measured, the relation of Cohen, Kaplan and Nelson (1999). This agreement in size follows from the Friedmann equation and is not counted as a derivation; the amount stays one measured number. On this reading dark energy need not be constant, and no law for its change is derived: the holographic formula of Li (2004) gives today's value but changes in the direction opposite to the hint reported by DESI (2025), and is not used. A fixed amount that never changes is kept on record as the alternative. The closeness of the scale to the neutrino masses is noted and not used. No state of the matrix itself is dark matter: its only chargeless state, the hidden neutrino partner, decays in 10\textsuperscript{\ensuremath{-}15} s or less, and dark matter as an entanglement response (Verlinde, 2017) matches the galactic acceleration scale (McGaugh, Lelli and Schombert, 2016) but fails published tests of galaxies, the Solar System and clusters. Reading dark matter as the medium's baseline resistance holds for galaxies: cH\textsubscript{0}/2\ensuremath{\pi} = 1.04 \ensuremath{\times} 10\textsuperscript{\ensuremath{-}10} m/s\textsuperscript{2} against Milgrom's scale of about 1.2 \ensuremath{\times} 10\textsuperscript{\ensuremath{-}10} m/s\textsuperscript{2} (Milgrom, 1983), and voids push (Hoffman, Pomarède, Tully and Courtois, 2017); it fails for the Bullet Cluster (Clowe et al., 2006) and the CMB (Planck Collaboration, 2020). The real part is the majoron of Section 8, a cold, chargeless wave.

\hypertarget{containment-the-proton-and-neutron-stars}{%
\subsection{6. Containment: the proton and neutron stars}\label{containment-the-proton-and-neutron-stars}}

\textbf{The proton.} Model the proton and the \ensuremath{\Delta} as bags in containment balance. The pressure B pulls in, and the motion of the quarks pushes out. The pion cloud around the bag is its leakage (the cloudy bag model). We fitted only the measured masses of the proton (938.92 MeV, isospin average) and the \ensuremath{\Delta} (1232 MeV), using the centre-of-mass correction of Donoghue and Johnson and static pion couplings. This gives B = 60.1 MeV/fm\textsuperscript{3}, or 57.1 to 63.7 for a \ensuremath{\pm}5\% pion--nucleon coupling. The proton's quark core then has an rms radius of 0.725 fm. As a check, the same calculation with the 1975 settings returns B = 58.4 MeV/fm\textsuperscript{3}, against the published value of about 57.5.

\textbf{Neutron stars.} In a neutron-star core, the same pressure B holds the quarks in, and the measured strong interaction pushes out, with a\textsubscript{4} = 1 \ensuremath{-} 2\ensuremath{\alpha}\textsubscript{s}/\ensuremath{\pi} and \ensuremath{\alpha}\textsubscript{s} run to the relevant scale. We solved the Tolman--Oppenheimer--Volkoff equations with Hinderer's tidal equation, with the surface correction for a self-bound star. The star is modelled as quark matter all the way to its surface (a strange quark star, Alcock, Farhi and Olinto, 1986), with a strange-quark mass of 100 MeV, \ensuremath{\alpha}\textsubscript{s} run at two loops from 0.1180 to twice to four times the quark chemical potential, and no hadronic layer. This model is consistent only if three-flavour quark matter is the ground state of matter while two-flavour matter is not (Bodmer, 1971; Witten, 1984; Farhi and Jaffe, 1984). With the pairing gap below, it is: at zero pressure strange matter has 820 to 880 MeV per baryon (below iron's 930) and up--down matter 941 to 1021 MeV (above 934, so ordinary nuclei are safe), the narrowest margin being 7 MeV at the lowest B and largest gap. Without pairing strange matter is not the ground state (940 to 962 MeV), which is a second reason the gap is needed. The proton fit uses no free zero-point term (Z\textsubscript{0} = 0), the centre-of-mass correction, a fitted bag colour coupling and the measured pion--nucleon coupling f\textsuperscript{2}/4\ensuremath{\pi} = 0.0778 with the bag form factor 3j\textsubscript{1}(kR)/(kR). A hybrid star with a hadronic outer core would change the numbers and is not computed here. The stars actually fix the combination 4\ensuremath{\Delta}\textsuperscript{2} \ensuremath{-} m\textsubscript{s}\textsuperscript{2}, which enters the pressure: 1.7 to 2.8 \ensuremath{\times} 10\textsuperscript{4} MeV\textsuperscript{2} at the central B. The strange mass runs from 93.5 MeV at 2 GeV to about 120 to 140 MeV at the densities of a star's core, which moves the gap window to about 88 to 109 MeV; at m\textsubscript{s} = 130 MeV and the shifted gaps, strange matter stays the ground state (819 to 877 MeV per baryon) and up--down matter stays above 934 MeV within each B's own window, the narrowest margin shrinking to about 5 MeV at the lowest B. Two caveats go with a strange star. Every compact star would then be a strange star, and large pulsar glitches and the vibrations seen in magnetar flares are hard to explain without a thick nuclear crust, unless the core holds a crystalline paired phase (Mannarelli, Rajagopal and Sharma, 2007). Observed neutron stars are called compact stars below where the strange-star model is meant.

\begin{itemize}
\tightlist
\item
  With the strong interaction alone, the heaviest star is 1.54 to 1.80 solar masses. This is below the 2.08 \ensuremath{\pm} 0.07 solar masses measured for PSR J0740+6620, so it fails.
\item
  Quarks in such cores are expected to pair (colour superconductivity). At the central B = 60.1 MeV/fm\textsuperscript{3}, a pairing gap \ensuremath{\Delta} of about 82 to 97 MeV gives a heaviest star of 2.0 to 2.13 solar masses, a radius of 10.8 to 11.3 km for a 1.4-solar-mass star, and a tidal deformability \ensuremath{\Lambda}\textsubscript{1.4} = 440 to 580.
\item
  These values meet both the pulsar mass and the GW170817 limit (\ensuremath{\Lambda}\textsubscript{1.4} below about 580). Across the full range of B allowed by the proton (57.1 to 63.7 MeV/fm\textsuperscript{3}), the window shifts to \ensuremath{\Delta} \ensuremath{\approx} 75 to 108 MeV, higher B requiring a larger gap.
\end{itemize}

The containment pressure comes from the proton, and the push comes from the measured strong force. The one remaining quantity, the pairing gap, becomes a prediction (Section 9).

\textbf{Stars and black holes.} For stars the measure of containment is the compactness C = GM/(Rc\textsuperscript{2}). The Sun has C \ensuremath{\approx} 0.000002, a neutron star about 0.18, the last stable configuration 4/9 (Buchdahl's limit), and a black hole 1/2. A star is under-contained and leaks energy. A black hole is over-contained. The same equations describe every step in between.

\hypertarget{relational-masses}{%
\subsection{7. Relational masses}\label{relational-masses}}

\textbf{The relational distance.} For any two masses, the quantity S = \ensuremath{\frac{1}{2}} ln(m\textsubscript{1}/m\textsubscript{2}) is the only measure that is unchanged when all masses are multiplied by the same factor. Their combined ratio is (m\textsubscript{1} + m\textsubscript{2})/\ensuremath{\surd}(m\textsubscript{1}m\textsubscript{2}) = 2 cosh S. Applying this to every pair of the known particles gives the particle-to-particle matrix. The antiparticle and decay matrices are built in the same way.

\textbf{The charged leptons.} Write the square roots of the three charged-lepton masses as the eigenvalues of one circulant matrix, M\textsubscript{0} \ensuremath{\cdot} circulant(1, re\textsuperscript{i\ensuremath{\delta}}, re\textsuperscript{\ensuremath{-}i\ensuremath{\delta}}). With r = 1/\ensuremath{\surd}2 and \ensuremath{\delta} = 2/9, the matrix gives the electron, muon and tau masses from a scale and the twist, to 0.003\%. Fitted with r left free, the twist is 0.222239 \ensuremath{\pm} 0.000040, against 2/9 = 0.222222 (0.4\ensuremath{\sigma}); with r fixed at 1/\ensuremath{\surd}2, the electron-to-muon ratio pins it at 0.22222205, which is 2/9 to about one part in a million (Section 10). The relation r = 1/\ensuremath{\surd}2 is equivalent to Koide's formula (Koide, 1983), in the geometric form of Foot and the parametrisation of Brannen. In the relational reading, the shared part and the relational part of the mass vector are equal, which is exactly r = 1/\ensuremath{\surd}2. Each part is 941.6 MeV, within 0.35\% of the proton mass. We note this coincidence without claiming an explanation.

\textbf{Quarks.} Running all six quark masses to a common scale gives balance angles of 52.15\ensuremath{^\circ} for the up-type quarks and 48.06\ensuremath{^\circ} for the down-type, against 45\ensuremath{^\circ} for the leptons. These angles do not depend on the scale chosen, so they are the targets any theory of quark masses must meet. A simple pattern of link strengths (3/6, 4/6, 5/6 for leptons, down-type and up-type) misses the data by about seven standard deviations and is rejected. The relation \ensuremath{\surd}(m\textsubscript{d}/m\textsubscript{s}) \ensuremath{\approx} sin \ensuremath{\theta}\textsubscript{C} (Gatto, Sartori and Tonin, 1968) holds to 0.04\%.

\textbf{Quark-mass relations, tested strictly.} Each relation contained in the particle matrices was tested with measured masses on both sides and judged against the measurement's own error bar. A full Higgs coupling of the top quark (y\textsubscript{t} = 1; Hill, 1981) gives m\textsubscript{t} = 174.1 GeV against 172.60 \ensuremath{\pm} 0.27 GeV (+0.9\%, but 5.6\ensuremath{\sigma}). The 45\ensuremath{^\circ} balance applied to the triples (c, b, t) and (s, c, b) (Rivero, 2011) gives m\textsubscript{c} = 1353 MeV against 1272.9 \ensuremath{\pm} 4.5 MeV (+14.8\ensuremath{\sigma}) and m\textsubscript{s} = 84.5 MeV against 92.9 \ensuremath{\pm} 0.7 MeV (\ensuremath{-}10\ensuremath{\sigma}); no standard definition of the charm mass rescues it. An earlier chain that started from the bottom mass alone fed its own high charm value into the strange step and so appeared to work to about 1\%. The Cabibbo relation, in its sine form sin \ensuremath{\theta}\textsubscript{C} = \ensuremath{\surd}(m\textsubscript{d}/m\textsubscript{s}), gives m\textsubscript{s}/m\textsubscript{d} = 19.9 against 19.8 \ensuremath{\pm} 0.3 from the measured masses (m\textsubscript{s} = 92.9 \ensuremath{\pm} 0.7 MeV and m\textsubscript{d} = 4.70 \ensuremath{\pm} 0.07 MeV; Particle Data Group, 2026), +0.3\ensuremath{\sigma}, which passes; earlier versions compared it with a value of 20.4 \ensuremath{\pm} 0.3 that has no published source. The tangent form gives 18.9 (\ensuremath{-}2.7\ensuremath{\sigma}); the sine form was chosen with both in view, so the form is a fitted choice (Appendix A.6). The up quark is reached through the decay \ensuremath{\eta} \ensuremath{\to} 3\ensuremath{\pi}, in which electromagnetic effects nearly cancel (Sutherland, 1966; Bell and Sutherland, 1968). Leutwyler's relation (m\textsubscript{u}/m\textsubscript{d})\textsuperscript{2} + (m\textsubscript{s}/m\textsubscript{d})\textsuperscript{2}/Q\textsuperscript{2} = 1 (Leutwyler, 1996), with the Cabibbo value of m\textsubscript{s}/m\textsubscript{d} and Q = 22.1 \ensuremath{\pm} 0.7 from the measured decays (Colangelo, Lanz, Leutwyler and Passemar, 2018), gives m\textsubscript{u}/m\textsubscript{d} = 0.437 \ensuremath{\pm} 0.060 (measured 0.462 \ensuremath{\pm} 0.020) and m\textsubscript{u} = 2.04 \ensuremath{\pm} 0.28 MeV (measured 2.16 \ensuremath{\pm} 0.07 MeV), with no lattice input. With the uncertainty of Q included, both lie 0.4\ensuremath{\sigma} from the measurement. Because Q is known to only 3\%, the prediction is wide, so this is a loose test. The result also depends on the choice of Q: the lattice value Q = 23.4 \ensuremath{\pm} 0.5 (Particle Data Group, 2026) gives m\textsubscript{u}/m\textsubscript{d} = 0.53 \ensuremath{\pm} 0.03 and m\textsubscript{u} = 2.47 \ensuremath{\pm} 0.14 MeV, 2.0\ensuremath{\sigma} above the measurement, on the line between a pass and a partial (Appendix A.6). Each relation is published separately; linking them through the matrices is the contribution here. The 45\ensuremath{^\circ} relations among the heavier quarks are numerical patterns good to 1--10\%, not laws, so the bottom, charm and strange masses remain measured inputs.

\textbf{Cone or spiral.} Written as exact statements and tested at \ensuremath{\mu} = 162.5 GeV, the lepton ladder is a cone, \ensuremath{\surd}m\textsubscript{k} \ensuremath{\propto} 1 + \ensuremath{\surd}2 cos(\ensuremath{\delta} + 2\ensuremath{\pi}k/3) with opening 45.000\ensuremath{^\circ} (Koide, 1983; Foot, 1994; Brannen, 2006). The quark openings are 48.150\ensuremath{^\circ} \ensuremath{\pm} 0.015\ensuremath{^\circ} (down-type) and 52.147\ensuremath{^\circ} \ensuremath{\pm} 0.018\ensuremath{^\circ} (up-type); they differ from 45\ensuremath{^\circ} by 213\ensuremath{\sigma} and 405\ensuremath{\sigma} and from each other by 174\ensuremath{\sigma}, and a single 45\ensuremath{^\circ} cone for all six quarks (Kartavtsev, 2011) fails by 13\ensuremath{\sigma}. Two cones about a common axis would give zero mixing, so the CKM tilt measures the misalignment of the two cone axes. The spiral (geometric steps) fails in all five forms tested, including Froggatt--Nielsen powers of the Cabibbo parameter (Froggatt and Nielsen, 1979), whose fitted exponents (2.0, 2.7, 4.2, 3.7) are near but not at integers; the quarks are nevertheless much closer to equal logarithmic steps than the leptons. The down-type opening is scale-independent under QCD and remains a fixed value to explain. No quark mass moves.

\textbf{Mutual pushes of quarks and leptons.} Renormalization-group running (Gell-Mann and Low, 1954; Machacek and Vaughn, 1984) couples the two sectors in three ways. Through the Higgs, the top shifts every lepton Yukawa coupling by 1.67\% per e-fold, the leptons shift the quarks by 6 \ensuremath{\times} 10\textsuperscript{\ensuremath{-}7}; this is flavour-universal and changes a ladder's scale, not its shape. Within a weak doublet the top drives the bottom by \ensuremath{-}0.83\% per e-fold with opposite sign, the route by which the two quark ladders bend each other. QED self-running separates the lightest and heaviest lepton by 2.84\%, and direct photon exchange between sectors is 10\textsuperscript{\ensuremath{-}5} to 10\textsuperscript{\ensuremath{-}4}. None closes a near miss: matching the lepton scale to m\textsubscript{p}/3 requires a common scale of about 1.3 GeV, where the cone opening is 45.05\ensuremath{^\circ} rather than 45\ensuremath{^\circ}; quark balance ratios move away from 5/4 and 5/3; and the 6.8\% charm and 8.9\% strength misses are far larger. The Koide relation holds for pole masses to 10\textsuperscript{\ensuremath{-}5} although QED running should spoil it by 0.35\% (Sumino, 2009), so any quark rule must specify its scale before it is tested.

\textbf{The three layers of mass.} Every composite mass is a Higgs layer (the elementary masses), a containment layer (constituent masses and hadrons; Sakharov and Zeldovich, 1966; De Rújula, Georgi and Glashow, 1975; 8 light baryons from 3 fitted numbers to 0.56\% rms) and a nuclear layer (the semi-empirical binding formula; Weizsäcker, 1935; Bethe and Bacher, 1936; 5 fitted numbers). Together they give the masses of 2,899 nuclides to 0.008\% rms (AME2020; Wang et al., 2021). For \textsuperscript{56}Fe the layers are 595 MeV (1.1\%), 52,001 MeV (99\%) and \ensuremath{-}494 MeV (\ensuremath{-}492 MeV measured). The proton is about 9\% quark mass, 32\% quark kinetic energy, 37\% gluon energy and 23\% trace anomaly (Ji, 1995; Yang et al., 2018), so about 91\% is containment energy. The structure holds; its numbers are fitted, and the bottom layer, the Higgs links, is open. Balance points of particle pairs fall on other particles' masses at chance level.

\textbf{Sizes and the proton's home.} The matrix of sizes S\textsubscript{ij} holds one number per particle. Five patterns in it fail chance tests, including whole powers of one parameter (Froggatt and Nielsen, 1979). Seen from the proton, each particle's balance mass is m\textsubscript{p}/(1 + e\textsuperscript{\ensuremath{-}2S}), a logistic curve (Verhulst, 1838): lighter particles at their own mass, heavier ones capped at the proton, the proton with itself at one half. This is the scale-view principle in exact form and moves no number.

\textbf{Scaling backwards.} Run at two loops from the top mass to the Planck energy (Machacek and Vaughn, 1984; Arason et al., 1992), the down-type opening moves from 48.15\ensuremath{^\circ} to 47.72\ensuremath{^\circ} and the up-type from 52.15\ensuremath{^\circ} to 52.31\ensuremath{^\circ}, while the leptons stay at 45.05\ensuremath{^\circ}; neither quark ladder is a 45\ensuremath{^\circ} cone at any energy, and m\textsubscript{b}/m\textsubscript{\ensuremath{\tau}} reaches only 0.62. The top Yukawa coupling, run down from large values (Pendleton and Ross, 1981; Hill, 1981), lands at 1.28--1.32 against 0.937: right in kind, 36--41\% high. A rule for the elementary masses must act on ratios, which do not run.

\hypertarget{quantum-structure}{%
\subsection{8. Quantum structure}\label{quantum-structure}}

\textbf{Entanglement.} Measured entangled pairs show the correlation E(a, b) = \ensuremath{-}cos(a \ensuremath{-} b), stronger than any local hidden-variable model allows (Bell). In the relational picture the two particles are one structure seen from two places. This is the content of the ER = EPR proposal of Maldacena and Susskind, and no signal travels faster than light.

\textbf{Uncertainty.} Any wave localised to a width \ensuremath{\Delta}x has a spread in wavenumber \ensuremath{\Delta}k with \ensuremath{\Delta}x\ensuremath{\cdot}\ensuremath{\Delta}k \ensuremath{\ge} \ensuremath{\frac{1}{2}}. Heisenberg's relation is this property of waves, applied to matter.

\textbf{Charges.} For one generation with N colours, the four anomaly-cancellation conditions fix the hypercharges uniquely. The up-quark charge is \ensuremath{\frac{1}{2}} + 1/(2N) and the down-quark charge is \ensuremath{-}\ensuremath{\frac{1}{2}} + 1/(2N). For N = 3 these are +2/3 and \ensuremath{-}1/3. Anomalies do not fix N; it is measured through the neutral-pion decay rate and the ratio of hadron to muon production in electron--positron collisions, and it follows from the inner algebra below.

\textbf{Forces, colours and generations.} The measured carrier counts (eight carriers with two neutral, and four with two neutral) allow only SU(3) and SU(2) \ensuremath{\times} U(1) in the classification of Lie algebras (Killing; Cartan), and the global form S(U(3) \ensuremath{\times} U(2)) passes the \ensuremath{\mathbb{Z}}\textsubscript{6} check for every field. A medium with no gaps has an inner algebra without zero divisors, and relationships that keep their size when combined require a conserved norm; by Hurwitz's theorem (1898) only \ensuremath{\mathbb{R}}, \ensuremath{\mathbb{C}}, \ensuremath{\mathbb{H}} and \ensuremath{\mathbb{O}} qualify. A force richer than the weak force requires \ensuremath{\mathbb{O}}, whose symmetry group G\textsubscript{2} leaves SU(3) when any one direction is fixed, as every observation fixes one; the remaining six directions form three complex ones, which are the three colours. \ensuremath{\mathbb{H}} acting from the left and from the right carries SU(2)\textsubscript{L} \ensuremath{\times} SU(2)\textsubscript{R}, but only under a different meaning of symmetry from the one that gives colour: on the number systems alone no single meaning yields the strong, weak and hypercharge carriers together, and constructions that join \ensuremath{\mathbb{C}}, \ensuremath{\mathbb{H}} and \ensuremath{\mathbb{O}} into one algebra (Dixon, 1994; Furey, 2016) reintroduce zero divisors, which the no-gaps rule forbids. The strong force is therefore derived. Read two at a time, the weak force, hypercharge and left-handedness would be taken from measurement; on the main line, with the relationship matrix read three at a time from one perspective (Section 8), they follow with one stated reading, and only the strengths are measured. The minimal particle content also needs at least one further Higgs field to match the measured masses (Georgi and Jarlskog, 1979). Hermitian octonionic matrices satisfy the Jordan identity only up to 3 \ensuremath{\times} 3 (Jordan, von Neumann and Wigner, 1934), while an irremovable CP-violating phase needs at least three generations (Kobayashi and Maskawa, 1973), so on this picture the number of generations is exactly three. These steps rest on an interpretive identification of the physical picture with algebraic rules. That forces are local symmetries (Weyl; Yang and Mills) follows from relativity and quantum mechanics, which the framework inherits. Completed with the 4 \ensuremath{\times} 4 group of Pati and Salam, this route is kept on record as the alternative picture; the main line is the relationship-matrix picture below.

\textbf{Mass as a balance point.} Mass is the energy at which containment and leakage balance. For the proton this gives a mass of four times the containment energy, about 99\% of the proton's mass when only the up and down quark masses are counted as Higgs mass; counting the strange quark the Higgs share is about 9\%, and with the heavy quarks' indirect effect about 29\% (Shifman, Vainshtein and Zakharov, 1978); the rest is set by one Higgs balance point (246.22 GeV), and one containment pressure (B = 60.1 MeV/fm\textsuperscript{3}) serves both the proton and the compact star. Fission is a shift of the balance (166 MeV against 173.3 MeV measured). The shape of each force follows from what space does to its field: spreading (1/r\textsuperscript{2}), confinement in a tube (constant pull), or screening (short reach); the nuclear force is shared leakage carried by pions. As strict numbers the balance falls short: the pull between quarks is 0.155 against 0.18--0.19 GeV\textsuperscript{2}, and the deconfinement temperature 105 against 156.5 \ensuremath{\pm} 1.5 MeV. It does not give elementary-particle masses, the strengths of the forces, or which forces exist.

\textbf{The main picture of the forces: the relationship matrix.} Where gravity emerges from entanglement, an exact symmetry must come with a force (Harlow and Ooguri, 2019; proven for spaces with a boundary), so forces are expected rather than added. Treating matter itself as motion in an inner space fails: a charged particle would weigh about 5 \ensuremath{\times} 10\textsuperscript{17} GeV, and a smooth inner space gives no left-handed matter (Witten, 1983). The 3 \ensuremath{\times} 3 Hermitian octonionic matrix, read as the relationship matrix seen by an observer who stands at one point and faces one way, has exactly 12 carriers, 4 neutral, in blocks 8 \textbar{} 3 \textbar{} 1, with one meaning of symmetry and no zero divisors (Todorov and Dubois-Violette, 2018; Dubois-Violette and Todorov, 2019; Krasnov, 2019; Baez, 2020--21; reproduced here). The weak pairs carry hypercharges in the ratio 3 : 1, and anomaly cancellation then fixes their partners at +2/3, \ensuremath{-}1/3 and \ensuremath{-}1. The matrix lacks exactly two of the partners a whole generation needs, those of the up quark and the electron (four states, one octonion seen from one perspective), and these partners and the strengths are one problem. Counting only what the matrix holds, one common strength predicts sin\textsuperscript{2}\ensuremath{\theta}\textsubscript{W} = 0.75 where the forces meet, against about 0.43 from the measured values; counting the whole generation (Chamseddine and Connes, 1996) gives 0.375, and, run to the Z mass at two loops, 0.2107 against 0.2312, 8.9\% off, the near miss of Georgi, Quinn and Weinberg (1974) and still a fail (Georgi and Glashow, 1974; Baez and Huerta, 2010). The strengths therefore remain inputs. Four homes for the missing partners built from two-way combinations (relationships between two links, the matrix's other turns, the doubled link space and a 4 \ensuremath{\times} 4 matrix, which is not a consistent algebra) all fail, and the three-way reading below supplies them; the generations cannot be slots inside the same matrix, so they are taken as three placements of the one shape, leaving ``at least three'' as a theorem and ``at most three'' resting on the nothing-extra rule and measurement; and a right-handed force cannot coexist with it, which removes the left--right scale and leaves about 0.4 of the observed matter excess. On this picture handedness is the side the observer faces: with the no-gaps rule applied to spin (a second added reading), mirror-image partners would be a gap, so matter must be one-handed, and anomaly consistency leaves the measured arrangement (Lee and Yang, 1956; Wu, 1957). The lepton 45\ensuremath{^\circ} balance holds to 15 parts in a million and follows, for leptons only, from the family symmetry and maximum uncertainty (Section 10), and the twist and the electron's mass are one open number. Any new force carrier would exclude it. Which forces exist is not a choice within this picture. Take the perspective principle as one place of the matrix kept and one imaginary direction kept. The symmetry group of the place is Spin(9), acting on the 16 real numbers that link it to the other two, and the part of Spin(9) that respects the direction is exactly 12-dimensional, 11 + 1, where 11 can only be 8 + 3; the result is the same for every direction (computed here for all seven basis directions and a random one). The eight complex link states carry the remaining charge in the ratio 3 : 1 with multiplicities 6 and 2, a quark pair in three colours and a lepton pair (Günaydin and Gürsey, 1973; Todorov and Dubois-Violette, 2018; Krasnov, 2019). The three forces therefore follow from the perspective principle by a theorem. What the theorem does not give are the strengths and the missing partners. This picture is carried as the main line, with the 4 \ensuremath{\times} 4 picture kept on record as the alternative.

\textbf{A relationship among three.} The first principle is read at two levels: two at a time for the numbers in the matrix, which forces the octonions, and three at a time for the states built from them, so that nothing combined with its own mirror image and itself gives nothing. The three-way product \{x, y, z\}, with the middle state entering as its mirror image (Jordan, 1933; Loos, 1977), has no zero divisors of this kind on the matrix's link states; the remaining two-way ones, ``the same state twice gives nothing'', are for matter the exclusion principle (Pauli, 1925). With the perspective read as the spot kept exactly and the facing kept as a direction (overall phase left out), the turns of the three-way product that respect it number exactly 12, four neutral, in blocks 8 \textbar{} 3 \textbar{} 1, and the 16 link states are a whole generation with the measured charges: quark pair +1/6, lepton pair \ensuremath{-}1/2, down partner +1/3, up partner \ensuremath{-}2/3, electron partner +1 and neutral partner 0. All anomaly sums vanish, no mirror states exist, and the two states picked by the perspective are the neutral and electron partners. This is the whole-generation pattern of Georgi (1975) and Fritzsch and Minkowski (1975), without their extra carriers. Other readings give 11, 13 or 14 turns or inconsistent charges, so the reading is an added step; applied to \ensuremath{\mathbb{C}} and \ensuremath{\mathbb{H}} it gives 3 and 5 kept turns, and only \ensuremath{\mathbb{O}} gives 8 \textbar{} 3 \textbar{} 1. Only 9 of the 12 turns coincide with the two-way count, so the weak force is re-identified: half the states are weak pairs and half singles, as measured, replacing the two-way reading above, under which every state was a pair. The strengths are unchanged: the geometry gives the whole-generation count, 0.375 with strong equal to weak, 8.9\% off at the Z mass at two loops. Charges as literal bound sets of three (Harari, 1979; Shupe, 1979) fail; the count of colour switches holds as a pattern (Georgi and Glashow, 1974; Furey, 2016). The framework adopts this reading. With this content the whole particle side, the Higgs, the mass links, the Majorana mass and the majoron included, can be defined exactly on a lattice of points if the construction of Craig (2026) holds (Appendix B).

\textbf{Perspective in place of handedness.} The three-way structure has exactly one i: only 1 and the facing direction commute with all 46 of its turns, so every state takes its screw sense from it and the 16 states cannot be wound independently. The perspective rule reads: one spot, one facing, one i throughout. Reversing the facing (i \ensuremath{\to} \ensuremath{-}i) exchanges matter with antimatter and left with right together, and the charge list of the 16 states differs from its mirror image, so the world is not mirror-symmetric (Weyl, 1929; Lee and Yang, 1956; Wu, 1957). A separate i for spin would double every state into mirror pairs, excluded by the nothing-extra rule; the earlier reading that spin uses the same i becomes part of the perspective rule. Mass is the link between a paired and a single state through the far side of the matrix, whose 10 parts carry the numbers of a colourless weak pair of charge \ensuremath{\pm}1/2 (the Higgs field's numbers) and six coloured parts (Freudenthal, 1954). These numbers appear twice, as in the standard E\textsubscript{6} construction (Gürsey, Ramond and Sikivie, 1976): as spin-\ensuremath{\frac{1}{2}} matter in each generation's 27 (the far-side weak pairs and D states of Sections 9 and 10, each with both handednesses, so not mirror families), and once as the spin-0 Higgs doublet that links paired to single states. Only the colourless part of the spin-0 copy is a particle; its coloured spin-0 parts must be read as matrix entries rather than particles (Dimopoulos and Georgi, 1981), and why is not derived (the doublet--triplet problem). Each state has exactly one mass partner, but the four links come out equal, so their sizes are not given. The same link gives no quark mixing, the right starting point, but not the small measured tilt.

\textbf{The forces that build the atom.} The largest set of the 12 turns under which the 16 states are their own mirror image is unique: the 8 strong turns plus one mixture of the single turn with a weak turn, which is electric charge, the two mirror-symmetric forces that bind (Vafa and Witten, 1984). Settling the colourless far-side doublet keeps exactly these 9 turns, the other three becoming W\textsuperscript{\ensuremath{\pm}} and Z; this is one statement rather than two, since a mass link maps the world onto its mirror image, and that the far side settles, in the direction keeping most turns (Michel and Radicati, 1971), is not derived. The kept charges give proton +1, neutron 0 and hydrogen 0 exactly (measured to 10\textsuperscript{\ensuremath{-}21}; Bressi et al., 2011). The state count gives the direction of running: colour confined near the proton scale, the weak force hidden 26 orders of magnitude before it would confine, electric charge free. The Higgs link is a Yukawa force whose charge is mass (Yukawa, 1935), the link that gives mass rather than the third binding force, which is B \ensuremath{-} L (below), with range 1/534 of the proton radius; it binds nothing and is not gravity, since universality of free fall holds to 1.5 \ensuremath{\times} 10\textsuperscript{\ensuremath{-}15} (MICROSCOPE, 2022) while the Higgs charge per unit mass is 1 for the electron and about 0.29 for the proton. A Higgs coupling to all energy would raise the LHC Higgs rate about 130-fold (ATLAS and CMS, 2022). With two settled directions, a complex far side allows a free top-to-bottom ratio (Lee, 1973), but bottom = \ensuremath{\tau} still fails. The whole matrix read three-way has exactly three independent sizes and no larger such matrix exists (Jordan, von Neumann and Wigner, 1934; Albert, 1934), a hint for three generations but not a derivation, since one matrix cannot hold both forces and generations.

\textbf{The energy the Higgs does not give, and matter number.} Under the 12 turns exactly one Higgs-free mass term is allowed: the hidden neutrino partner with itself, the state that the perspective rule keeps exactly and that therefore carries no charge. With no perspective, or with the spot held only as a direction, there is none. This is the Majorana mass of the seesaw (Majorana, 1937; Minkowski, 1977; Yanagida, 1979; Gell-Mann, Ramond and Slansky, 1979); its size is an input. Keeping the spot exactly breaks one further generator, a 13th commuting with the 12, whose charges are +1/3 on quarks and \ensuremath{-}1 on leptons: B \ensuremath{-} L (Davidson, 1979; Marshak and Mohapatra, 1980). Strong, electric and B \ensuremath{-} L form the largest facing-blind subset of the 13. The far side settling breaks the weak generators and gives the Dirac masses; the spot breaks B \ensuremath{-} L and gives the Majorana mass; together they leave the 9 of SU(3) \ensuremath{\times} U(1)\textsubscript{em}. Two readings are possible: (ii) B \ensuremath{-} L is gauged above about 10\textsuperscript{11} GeV, with fifth-force limits (Wagner et al., 2012) requiring a heavy carrier; or (i) B \ensuremath{-} L is never a force. The main line adopts (i); (ii) is recorded as the alternative. Because gravity allows no exact global symmetry (Banks and Seiberg, 2011; Harlow and Ooguri, 2019), reading (i) keeps the spot at every scale except the grain of space: B \ensuremath{-} L is broken there, very slightly, and that small breaking is what gives the majoron its mass (Akhmedov, Berezhiani, Mohapatra and Senjanović, 1993; Rothstein, Babu and Seckel, 1993). In reading (i) the phase of the field that settles the Majorana mass survives as a massless or nearly massless Goldstone boson, the majoron (Chikashige, Mohapatra and Peccei, 1981), which is part of a piece the framework already has rather than an extra state. The Higgs-free mass does not separate the generations (no such term touches a charged state; a common value gives m\textsubscript{2}/m\textsubscript{3} = 0.0035 against 0.172), and a generation-blind Higgs fails the LHC muon-to-tau comparison by about 68\ensuremath{\sigma} (ATLAS and CMS, 2022). The Higgs value, G and the cosmological constant are all properties of space rather than turns, but the Higgs value melts above 159.5 GeV (D'Onofrio and Rummukainen, 2016) and trial relations among them miss by factors of 90 to 450. The matrix supplies a \ensuremath{\mathbb{Z}}\textsubscript{2}, (\ensuremath{-}1)\textsuperscript{3(B\ensuremath{-}L)}, under which all 16 states are odd and all 10 far-side parts even (Krauss and Wilczek, 1989). Among the 16 it protects only the lightest neutrino, which gives 0.5\% of the dark matter (Kadastik, Kannike and Raidal, 2010). For the far-side parts the sign cannot be exact: if it were, the lightest neutral far-side weak state would be stable, and a stable state of that kind is excluded by direct-detection searches by about two orders of magnitude; it would also forbid the mixing with ordinary matter through which the far-side parts decay. The sign is therefore approximate there, and no prediction rests on it. The one dark matter candidate the matrix allows is the majoron, as a lead. With the settling scale near the heavy partners (about 10\textsuperscript{11} to 10\textsuperscript{12} GeV), its lifetime is at least 10\textsuperscript{25} s; below about 37 eV it behaves as wave dark matter (Hu, Barkana and Gruzinov, 2000; Hui et al., 2017), allowed above about 2 \ensuremath{\times} 10\textsuperscript{\ensuremath{-}20} eV (Rogers and Peiris, 2021). Its mass and abundance are not derived (Berezinsky and Valle, 1993; Garcia-Cely and Heeck, 2017). Produced by misalignment (Preskill, Wise and Wilczek, 1983; Abbott and Sikivie, 1983; Dine and Fischler, 1983) after a start colder than the settling value f (the lightest partners may share the heat, but majorons made by their pairs are suppressed by about (M\textsubscript{N}/f)\textsuperscript{2}), it is all of the dark matter when m\_J = 0.027 eV \ensuremath{\times} (1.5 \ensuremath{\times} 10\textsuperscript{12} GeV \ensuremath{\div} f)\textsuperscript{4} \ensuremath{\div} angle\textsuperscript{4} (equivalently 8.8 keV \ensuremath{\times} (6.3 \ensuremath{\times} 10\textsuperscript{10} GeV \ensuremath{\div} f)\textsuperscript{4} \ensuremath{\div} angle\textsuperscript{4}), where f is the settling value and the angle is where the field started, so its mass is not fixed, and it is cold. Because the field settled before the start, the start's own fluctuations are imprinted on it as fluctuations in dark matter uncorrelated with the rest; the CMB limit on these (Planck Collaboration, 2020; fraction below 0.038) caps the start's expansion rate, which leaves primordial gravitational waves at r \ensuremath{\lesssim} 0.08 (f \ensuremath{\times} angle \ensuremath{\div} \ensuremath{\bar{M}}\textsubscript{P})\textsuperscript{2}, with \ensuremath{\bar{M}}\textsubscript{P} the reduced Planck mass, below about 10\textsuperscript{\ensuremath{-}11} for any settling value used here; this needs B \ensuremath{-} L to stay broken through the start and afterwards (f above both the start's expansion rate \ensuremath{\div} 2\ensuremath{\pi} and the reheating temperature) (Fox, Pierce and Thomas, 2004; Hertzberg, Tegmark and Wilczek, 2008). This is the origin of the gravitational-wave test in Appendix A.4, and it belongs to the majoron lead; a hot start leaves about 73\% of it hot, which fails. A non-thermal start that hands its energy directly to the heavy partners (Lazarides and Shafi, 1991) avoids this and keeps leptogenesis, at the cost of one number. \emph{With the alternative start phase} (Section 5), the expansion rate during the start, 1.6 \ensuremath{\times} 10\textsuperscript{13} GeV, divided by 2\ensuremath{\pi}, exceeds every settling value used here, so the majoron's starting angle is averaged rather than chosen: there are no separate dark-matter ripples and the angle is no longer a start condition. The bound r \ensuremath{\lesssim} 10\textsuperscript{\ensuremath{-}11} above then applies only to the cold start, which is the main line; the alternative predicts r = 0.0040. Keeping the majoron cold needs the heat after the start below about 1.6 \ensuremath{\times} 10\textsuperscript{9} GeV, which the Higgs--curvature coupling \ensuremath{\xi} \ensuremath{\approx} 0.12--0.15 gives; whether hot majorons from the partners stay under a few per cent at that heat has not been recalculated.

\textbf{Neutrinos and the matter--antimatter excess.} If the neutrino shares the relation matrix of the charged leptons, its Dirac link is the charged-lepton mass matrix, and the seesaw mechanism (Minkowski; Yanagida; Gell-Mann, Ramond and Slansky; Glashow; Mohapatra and Senjanović) fixes the heavy partners' mass matrix as M\textsubscript{N} = \ensuremath{-}D m\textsubscript{\ensuremath{\nu}}\textsuperscript{\ensuremath{-}1} D, with D = diag(m\textsubscript{e}, m\textsubscript{\ensuremath{\mu}}, m\textsubscript{\ensuremath{\tau}}) and m\textsubscript{\ensuremath{\nu}} built from the light masses and the measured mixing. Computed as the full 3 \ensuremath{\times} 3 matrix with TM1 mixing, the partners weigh 0.6 to 1.9 \ensuremath{\times} 10\textsuperscript{5}, 2.2 to 6.2 \ensuremath{\times} 10\textsuperscript{8} and 1.3 to 1.65 \ensuremath{\times} 10\textsuperscript{12} GeV across \ensuremath{\theta}\textsubscript{23} and the unmeasured phases; earlier drafts quoted 7.3 \ensuremath{\times} 10\textsuperscript{5}, 1.3 \ensuremath{\times} 10\textsuperscript{9} and 6.3 \ensuremath{\times} 10\textsuperscript{10} GeV from the ratio m\textsubscript{\ensuremath{\ell}}\textsuperscript{2}/m\textsubscript{\ensuremath{\nu}}, which holds only without mixing. The heaviest then needs a right-handed coupling of about 3 for left--right breaking at 10\textsuperscript{11.7} GeV, the scale at which the three measured force strengths unify in the minimal Pati--Salam model (proton lifetime about 10\textsuperscript{36} years); that is near the limit where the calculation stays reliable (about \ensuremath{\surd}(4\ensuremath{\pi}) \ensuremath{\approx} 3.5), so the tie to Pati--Salam is strained rather than comfortable. Flavoured leptogenesis from these partners (Fukugita and Yanagida, 1986), solved with the full rate equations (Buchmüller, Di Bari and Plümacher; Blanchet and Di Bari; Nardi, Nir, Roulet and Racker), does not predict the observed baryon asymmetry of 6.1 \ensuremath{\times} 10\textsuperscript{\ensuremath{-}10}: with the shared-matrix rule read literally (this calculation always used the full matrix, so it is unchanged) and only the unmeasured phases free, the largest value found is about 1.5 \ensuremath{\times} 10\textsuperscript{\ensuremath{-}13}, some 4,000 times short. One narrow region survives every test tried and allows the observed value: \ensuremath{\delta} \ensuremath{\approx} 274\ensuremath{^\circ}, consistent with current neutrino data, a hot start within a window of about five in temperature (capped by magnetic monopoles; Zeldovich and Khlopov, 1978; Preskill, 1979), and a lightest partner almost unlinked to the electron. The same shared-matrix rule fails for quark pairs (m\textsubscript{t}/m\textsubscript{b} = 41), so it is stated for leptons only.

\textbf{Lifetimes.} Lepton universality, written relationally, gives the tau lifetime from the muon's as \ensuremath{\tau}\textsubscript{\ensuremath{\tau}} = \ensuremath{\tau}\textsubscript{\ensuremath{\mu}} \ensuremath{\cdot} BR \ensuremath{\cdot} e\textsuperscript{\ensuremath{-}10S}. This equals 289.7 fs, against 290.3 fs measured.

\textbf{Carrier and energy matrices.} Every interaction joins two matter particles through one carrier. Written as matrices, the photon, Z, W, gluon and Higgs couplings reproduce the measured charges, and the main line predicts \ensuremath{\rho} = m\textsubscript{W}\textsuperscript{2}/(m\textsubscript{Z}\textsuperscript{2} cos\textsuperscript{2}\ensuremath{\theta}\textsubscript{W}) = 1 exactly, because the far side holds exactly one weak doublet of hypercharge \ensuremath{\frac{1}{2}} (Veltman, 1977); measured, 1.0003 \ensuremath{\pm} 0.0005 after known corrections. Veltman's condition for cancelling quadratic divergences (Veltman, 1981) needs a 313 GeV Higgs, texture zeros (Fritzsch, 1977) miss the mixing, and equal entries miss: the open masses are the Higgs-matrix entries and the tilt is the W matrix. The pair link value (m\textsubscript{1} + m\textsubscript{2})/\ensuremath{\surd}(m\textsubscript{1}m\textsubscript{2}) is exactly the inverse reduced-mass factor, so the binding energies of hydrogen (13.598 eV), positronium (6.803 eV), muonium (13.540 eV) and muonic hydrogen (2.528 keV) are reproduced, as in Bohr's treatment (Bohr, 1913) with the reduced mass; this is a reformulation, not an independent test. The same matrices give the top width (1.35 GeV against 1.42 (+0.19 \ensuremath{-}0.15) GeV) and light hadron masses to a few per cent with five fitted numbers (De Rújula, Georgi and Glashow, 1975). Gravity written as a perspective fails in matrix form; for every elementary pair the Higgs-to-gravity force ratio is 1.96 \ensuremath{\times} 10\textsuperscript{32}. The Higgs cannot be a top-quark condensate (Bardeen, Hill and Lindner, 1990).

\hypertarget{predictions}{%
\subsection{9. Predictions}\label{predictions}}

\textbf{Awaiting measurement.} Each result in this section follows within the framework but rests on a measurement that has not yet returned. None is counted as established until it does; the deciding measurement is named with each.

\textbf{The strong coupling (awaiting measurement).} The force strengths meet at 4.6 \ensuremath{\times} 10\textsuperscript{13} GeV if the three colourless far-side weak pairs of the relationship matrix sit at 2.4 \ensuremath{\times} 10\textsuperscript{8} GeV (Section 10). At three loops this requires \ensuremath{\alpha}\textsubscript{s}(M\textsubscript{Z}) = 0.1178, with the grain of space fixed in advance as the Planck energy divided by \ensuremath{\surd}2. The grain was fixed on 17 November 2026, before the deciding measurement. It, and the form of the rule for the pairs' weight, were chosen with today's data in view (fitted, not predicted), so only later measurements test the prediction; the unreduced Planck energy would have given 0.1196, which is not the framework's value. The world average is 0.1180 \ensuremath{\pm} 0.0009; the most precise single determination is 0.11873 \ensuremath{\pm} 0.00056 (Dalla Brida and others, 2025), which lies 1.7 standard deviations above it and does not yet decide. Lattice determinations at about 0.3\% will. An observed proton decay would also falsify this part.

\textbf{The Higgs mass (awaiting measurement).} The framework's balance of containment, applied at the grain of space, says that the Higgs holds itself together down to the grain and no further, so its self-coupling vanishes there. This is the rule of Shaposhnikov and Wetterich (2010), which assumes that gravity stays well behaved at the grain (Weinberg, 1979; Reuter, 1998). Run at two loops (Machacek and Vaughn, 1984; Luo, Wang and Xiao, 2003) from the starting values of Buttazzo et al.~(2013), the Standard Model's states alone give m\textsubscript{H} = 128.8 \ensuremath{\pm} 1.5 GeV, which is 2.6 standard deviations above the measured 125.13 \ensuremath{\pm} 0.11 GeV (Particle Data Group, 2026). With the far-side states of the relationship matrix included, three colourless weak pairs at 2.4 \ensuremath{\times} 10\textsuperscript{8} GeV and their three coloured partners at 4.6 \ensuremath{\times} 10\textsuperscript{13} GeV, both fixed beforehand by the force strengths and not adjusted here, the result is 127.3 \ensuremath{\pm} 1.5 GeV, within 1.4 standard deviations. Estimating the three-loop step from the Standard Model's own two-to-three-loop shift (Buttazzo et al., 2013), a step decided before it was run and accepted whichever way it went, moves this by \ensuremath{-}0.1 GeV, to 127.2 \ensuremath{\pm} 1.5 GeV, and changes no verdict (\texttt{analysis/higgs\_three\_loop\_check\_2026-11-18.py}). This is an estimate: the far-side states enter only above 2.4 \ensuremath{\times} 10\textsuperscript{8} GeV, while the shift comes mostly from the top-quark and strong-force matching near the weak scale, but a full three-loop run that includes the far-side states has not been done. The test and its pass line were fixed before it was run. The Higgs mass therefore leaves the list of measured inputs, while the Higgs value stays measured. This is not counted as established until measurement returns. The rule needs a top pole mass of 171.6 GeV against 172.60 \ensuremath{\pm} 0.27 GeV measured (Particle Data Group, 2026); the quoted error allows 0.5 GeV for the ambiguity in what the measured top mass means, and without that allowance the gap is 2.1 standard deviations, just outside the pass line of 2 set before the test, so the pass rests on that allowance; on the top mass alone (171.6 needed against 172.60 \ensuremath{\pm} 0.27 measured) it is 3.7. A sharper top mass or strong coupling that pushes the gap beyond 3 standard deviations would falsify it. Any extra weak matter at high energy would help in the same way, so the agreement does not single out this framework.

\textbf{No proton decay.} In the 27 states of each generation (Gürsey, Ramond and Sikivie, 1976; Slansky, 1981) every charge is exact and all five anomaly sums vanish. With the framework's content, the twelve known carriers, one Higgs pair, the majoron and no coloured spin-0 states, no renormalisable interaction changes baryon number, since that takes three quarks at once (Weinberg, 1979b; Wilczek and Zee, 1979). The only remaining source is the grain of space itself, which gives a lifetime near 10\textsuperscript{44} years against the measured floor of 2.4 \ensuremath{\times} 10\textsuperscript{34} years (Takenaka et al., 2020). The far side's spin-\ensuremath{\frac{1}{2}} D states do not change this: with only a colourless spin-0 field, every renormalisable term they enter (their mass, their mass link to the ordinary down quarks and their Higgs link) keeps baryon number when D carries 1/3, as d does. The statement rests on the nothing-extra rule: coloured spin-0 partners of the Higgs would bring the decay back.

\textbf{Quark mixing stays unitary.} The only states that mix with known quarks are the D states, at 4.6 \ensuremath{\times} 10\textsuperscript{13} GeV. Their mixing with right-handed down quarks causes no flavour-changing Z couplings at tree level, since both carry the same charges, and their mixing with left-handed quarks is suppressed by (v\textsubscript{H}/M\textsubscript{D})\textsuperscript{2} \ensuremath{\approx} 10\textsuperscript{\ensuremath{-}23}. The first row of the quark-mixing matrix therefore sums to one: \textbar V\textsubscript{ud}\textbar{}\textsuperscript{2} + \textbar V\textsubscript{us}\textbar{}\textsuperscript{2} + \textbar V\textsubscript{ub}\textbar{}\textsuperscript{2} = 1 to far better than any measurement. The present shortfall, 0.0016 \ensuremath{\pm} 0.0007 (2.3 standard deviations; Particle Data Group, 2026; the Cabibbo-angle anomaly), must shrink to nothing as nuclear and radiative corrections improve; a confirmed shortfall would falsify the far-side reading.

\textbf{Neutrino masses.} The same relational rule, applied to neutrinos with the two measured mass splittings, gives, with today's limit from cosmology, normal ordering. With JUNO's 2026 splittings (\ensuremath{\Delta}m\textsuperscript{2}\textsubscript{21} = 7.388 \ensuremath{\times} 10\textsuperscript{\ensuremath{-}5}, \ensuremath{\Delta}m\textsuperscript{2}\textsubscript{31} = 2.509 \ensuremath{\times} 10\textsuperscript{\ensuremath{-}3} eV\textsuperscript{2}) the masses are m\textsubscript{1} = 0.36 meV, m\textsubscript{2} = 8.6 meV and m\textsubscript{3} = 50.1 meV, summing to 0.0591 eV; with the global fit NuFIT 6.0 (including atmospheric data) m\textsubscript{2} = 8.7 meV, m\textsubscript{3} = 50.3 meV and the sum is 0.0594 eV. Earlier versions quoted 8.7 meV and 0.0592 eV, from the global fit without atmospheric data; the change in v1.3.3 is only the input. The rule also has an inverted solution, with a sum of 0.103 eV; it is excluded by the current cosmological bound, not by the rule itself. The sum lies just above the minimum normal ordering allows (0.0588 eV), so the sharpest number is the lightest mass. JUNO's first ordering analysis favours normal ordering at 2.1\ensuremath{\sigma} (JUNO Collaboration, 2026), not yet decisive; it will reach 3\ensuremath{\sigma} within a few years, and cosmological surveys are approaching the required sensitivity to the sum. Inverted ordering, or a sum clearly different from 0.059 eV, would falsify this part of the framework.

\textbf{Neutrino mixing (conditional).} Exact tribimaximal (Harrison, Perkins and Scott, 2002) and trimaximal mixing are excluded by the data, while the TM1 pattern (Xing and Zhou, 2007; Albright and Rodejohann, 2009) gives sin\textsuperscript{2}\ensuremath{\theta}\textsubscript{12} = 0.318. That agreed with the global fits within 1\ensuremath{\sigma}, but JUNO's first precision measurement gives 0.3092 \ensuremath{\pm} 0.0087 from its first 59 days (JUNO Collaboration, 2025) and 0.3036 \ensuremath{\pm} 0.0064 with more data (JUNO Collaboration, 2026), which puts TM1 2.3 standard deviations high. This is not yet a failure, which needs 3\ensuremath{\sigma}, but it is the sharpest pressure on the framework today: if JUNO's next data keep this central value, TM1 fails, and with it the savings of \ensuremath{\theta}\textsubscript{12} and the neutrino CP phase (Appendix A.2) and the phase prediction below. Awaiting measurement: not 100\% until JUNO's next release. If the neutrino flavours are perspectives on one relational structure of this kind, the CP-violating phase lies near 270\ensuremath{^\circ}, close to maximal CP violation; current global fits weakly prefer 180\ensuremath{^\circ} to 215\ensuremath{^\circ}, and DUNE and Hyper-Kamiokande will decide.

\textbf{Quark pairing in compact stars.} With the containment pressure set by the proton, compact stars require a quark pairing gap of about 82 to 97 MeV for a strange-quark mass of 100 MeV (75 to 108 MeV across the proton's range of B); what they fix is 4\ensuremath{\Delta}\textsuperscript{2} \ensuremath{-} m\textsubscript{s}\textsuperscript{2} \ensuremath{\approx} 1.7 to 2.8 \ensuremath{\times} 10\textsuperscript{4} MeV\textsuperscript{2}, which gives 88 to 109 MeV for a running strange mass of 120 to 140 MeV. Theoretical estimates span 50 to 150 MeV, and the gap has not been measured. A firm determination outside this band, from lattice or effective-theory calculations, or from compact-star cooling and merger signals, would falsify this part. The range of 2.0 to 2.15 solar masses is the observed window that the gap must respect, not a separate prediction. A star clearly above about 2.15 solar masses would break this part; PSR J0952\ensuremath{-}0607, at 2.35 \ensuremath{\pm} 0.17 solar masses (Romani and others, 2022), and the 2.6 solar-mass object in GW190814 (Abbott and others, 2020) are the cases to watch.

\textbf{Neutrinoless double beta decay.} The Higgs-free mass of Section 8 exists only if neutrinos are Majorana particles (Majorana, 1937). Neutrinoless double beta decay (Furry, 1939) then occurs with an effective mass of 1.3 to 4.1 meV with TM1 mixing and both Majorana phases free (1.3 to 4.0 with JUNO's splittings, 1.4 to 4.1 with the global fit; 1.8 to 2.3 meV if the relational matrix fixes them at zero, a reading not yet established), below the present limit of 28 to 122 meV (KamLAND-Zen, 2024) and below the 10 to 20 meV reach of planned experiments. Proof that neutrinos are Dirac particles, or an effective mass clearly outside this range, would falsify it.

\textbf{Primordial gravitational waves and the tilt of the first ripples (alternative start; awaiting measurement).} On the main line no primordial gravitational waves are expected at a measurable level (r \ensuremath{\lesssim} 10\textsuperscript{\ensuremath{-}11} within the majoron lead), and the tilt is a measured input. The alternative start phase (Section 5) predicts r = 0.0040 \ensuremath{\pm} 0.0003 for N = 54--55, below today's limit r \textless{} 0.036 (BICEP/Keck, 2021). Simons Observatory and LiteBIRD reach about 0.003 and 0.001 in the early 2030s; not seeing primordial gravitational waves at that level would falsify the start phase. The tilt is n\textsubscript{s} = 0.9630--0.9636. Against published constraints (McDonough and Ferreira, 2025): Planck 0.9638 \ensuremath{\pm} 0.0040 (\ensuremath{-}0.1\ensuremath{\sigma}); Planck with DESI baryon acoustic oscillations 0.9690 \ensuremath{\pm} 0.0035 (\ensuremath{-}1.6\ensuremath{\sigma}); Planck, ACT and SPT combined 0.9693 \ensuremath{\pm} 0.0029 (\ensuremath{-}2.1\ensuremath{\sigma}); ACT DR6 with Planck and DESI 0.9743 \ensuremath{\pm} 0.0034 (\ensuremath{-}3.2\ensuremath{\sigma}; ACT Collaboration, 2025); Planck, ACT and SPT with DESI 0.9737 \ensuremath{\pm} 0.0025 (\ensuremath{-}4.2\ensuremath{\sigma}). It therefore agrees with the CMB alone and fails once DESI is combined; that shift is driven by a disagreement of about 2\ensuremath{\sigma} between DESI and the CMB (McDonough and Ferreira, 2025), which these combinations assume away. This is why the stretching phase is the alternative and not the main line: a CMB-only tilt above about 0.972 at 3\ensuremath{\sigma}, or DESI and the CMB brought into agreement at the higher value, would falsify the start phase in this form (Bezerra-Sobrinho and Medeiros, 2025, discuss the same tension). Not 100\% until those measurements return.

\textbf{Unmeasured values.} The carrier and energy matrices give values not yet measured: the Higgs self-couplings (three-point 191 GeV, four-point 0.776), the HHWW coupling (0.213), Higgs couplings to the charm quark, the electron and the light quarks equal to m/246 GeV, and the binding energies of true muonium (1.407 keV), tauonium (23.66 keV) and the proton--tau atom (16.35 keV). Most are standard physics in the framework's form, so they test its bookkeeping rather than new physics. The dark matter candidate is a cold, chargeless spin-0 wave decaying only to neutrinos (its mass is not fixed); dark matter that feels the weak force would falsify it.

\hypertarget{open-questions}{%
\subsection{10. Open questions}\label{open-questions}}

\textbf{Quark masses.} Of the relations in Section 7, only the down and up links pass strict tests; the bottom, charm and strange masses are measured inputs. The up-quark step uses no lattice QCD, but its value of Q comes from a dispersion analysis of measured decays, whose \ensuremath{\pm}3\% uncertainty dominates. Every relation tried for the bottom mass fails, including b = \ensuremath{\tau} at the left--right scale (26\% off) and the 45\ensuremath{^\circ} colour rule, which misses charm by 15--19\ensuremath{\sigma}. The scale-free balance angles of Section 7 remain the targets (the down-type cone opening is 48.15\ensuremath{^\circ} \ensuremath{\pm} 0.015\ensuremath{^\circ} with current lattice masses); the cone fits the leptons only, and the spiral fits no ladder exactly. Quark mixing belongs with these masses: the simplest mass link gives none, and the small measured tilt is not produced. The strongest lead is a Schur (Feshbach) fold of one heavy far-side colour-triplet state of charge \ensuremath{-}1/3 (the D quark of the 27 of E\textsubscript{6}; Gürsey, Ramond and Sikivie, 1976) along the TM1 direction (\ensuremath{-}1, \ensuremath{-}1, 2) (Harrison, Perkins and Scott, 2002; Schur, 1917; Feshbach, 1958): with one fitted strength all three mixing angles come within 19\%, and 2\% of random directions do as well. The direction deepens the ladder's existing break; the 19\% excess is far outside the errors, so it is a lead, not a result.

\textbf{Force groups (awaiting measurement).} Section 8 obtains the strong force and three colours from the full-medium rules (building on Hurwitz; Günaydin and Gürsey; Dixon; Furey; Connes and Chamseddine). In the main picture, the relationship matrix read three at a time from one perspective holds every particle of a generation with its measured charge, the measured carrier counts, the weak hypercharges and one-handed matter, at the cost of stated readings; the 4 \ensuremath{\times} 4 picture, in which the weak force, hypercharge and handedness had to be taken from measurement, is kept on record as the alternative. Only the force strengths remain under this heading. With the standard two-loop method the Standard Model alone misses the measured weak mixing by \ensuremath{-}8.9\% \ensuremath{\pm} 0.1\%. The gap closes if the three colourless far-side weak pairs of the matrix (spin \ensuremath{\frac{1}{2}}, no colour, one per generation) sit at 2.2 \ensuremath{\times} 10\textsuperscript{8} GeV, with the coloured far-side states at the meeting, 4.6 \ensuremath{\times} 10\textsuperscript{13} GeV; the doublet--triplet split is the known one (Dimopoulos and Wilczek, 1981). These states add nothing to the matrix. Their weight is the meeting energy squared divided by the grain of space, 2.4 \ensuremath{\times} 10\textsuperscript{8} GeV, close to the 2.2 \ensuremath{\times} 10\textsuperscript{8} GeV the strengths require. This rule was written down after the required weight was known, so the agreement is not evidence; it is one of the fitted choices (Appendix A.2, A.6). Stated exactly, the bet is one assumption: the weak pairs have no direct link to the background value that weighs their coloured partners, only the gravitational one. A single shared E\textsubscript{6} coupling would weigh both alike, put the pairs at the meeting energy and fail to close the strengths; the two couplings are not tied together here because the forces never merge into one group (Masiero, Nanopoulos, Tamvakis and Yanagida, 1982). This is assumed, not derived. They cannot be stable, since any produced during reheating would exceed the dark matter many times over (Chung, Kolb and Riotto, 1999; Giudice, Kolb and Riotto, 2001); they decay through mixing with the ordinary lepton doublets, the colourless counterpart of the mixing that tilts the down quarks, which adds one small coupling (about 4 \ensuremath{\times} 10\textsuperscript{\ensuremath{-}16}) to the count. If the universe never reheated above 2.4 \ensuremath{\times} 10\textsuperscript{8} GeV, as the majoron lead requires, none were made and that coupling is not needed. With the alternative start phase the heat reaches about 0.5--5 \ensuremath{\times} 10\textsuperscript{9} GeV, so they would be made, and the coupling would be required: it must exceed about 9 \ensuremath{\times} 10\textsuperscript{\ensuremath{-}12} so that they decay (within about 10\textsuperscript{\ensuremath{-}9} s) before they would dominate near 23 GeV and thin out the matter excess (\texttt{analysis/start\_phase\_starobinsky\_2026-11-19.py}). Not derived: the meeting energy itself, the grain factor, and gravity's strength. The test is the strong coupling (Section 9). The meeting is a crossing, not a merging. Computed at two loops, the three strengths agree at 4.6 \ensuremath{\times} 10\textsuperscript{13} GeV to 0.02 in 1/\ensuremath{\alpha} and then separate again; at the grain of space they are 27, 41 and 49. No placement of the far-side states makes them equal at the grain (the smallest spread found is 17 in 1/\ensuremath{\alpha}). In a unified group the equality would follow from the merging (Georgi and Glashow, 1974). The framework has no merged group, so it needs a boundary condition instead: the three strengths are equal where each generation's 27 becomes complete, that is, where the coloured far-side states enter. This is the second assumption inside the bet, alongside the weight rule, and it is not derived; the prediction \ensuremath{\alpha}\textsubscript{s} = 0.1178 rests on both. This criterion stays partly met until that measurement returns.

\textbf{The twist.} The measured twist, 0.222222 to about one part in a million, matches 2/9 but is not derived. It cannot be exactly 2/9 if the 45\ensuremath{^\circ} balance is also exact, since the two together miss the electron-to-muon mass ratio by 10 parts per million; reading 2/9 as a fraction of a ring is not an explanation. Rules that reproduce it for the leptons (\ensuremath{\delta} = Q/3 and related forms) fail for both quark families.

\textbf{The 45\ensuremath{^\circ} balance, derived for leptons.} Two of the framework's principles give it. Three placements of one shape (Section 8) mean the family matrix of amplitudes, the square roots of the masses, is unchanged by a turn of 120\ensuremath{^\circ}, which makes it circulant, with a the tie to itself and b the ties to the others (Brannen, 2006). Maximum uncertainty between self and relationship means the self share of the matrix's weight, p = 3a\textsuperscript{2}/(3a\textsuperscript{2} + 6\textbar b\textbar{}\textsuperscript{2}), is \ensuremath{\frac{1}{2}}; every symmetric measure of uncertainty peaks there (Shannon, 1948). And p = \ensuremath{\frac{1}{2}} is exactly \textbar b\textbar{}\textsuperscript{2} = a\textsuperscript{2}/2, which is Koide's Q = 2/3 and the 45\ensuremath{^\circ} angle (Koide, 1983; Foot, 1994). The charged leptons measure p = 0.50000, and the tau mass follows at 1776.969 MeV, +0.4\ensuremath{\sigma} against 1776.93 \ensuremath{\pm} 0.09 MeV (PDG 2024). The rule holds for the pole masses to about 10\textsuperscript{\ensuremath{-}5}, although running from a high scale should spoil it by 0.35\% (Sumino, 2009); the framework does not explain this. For the neutrinos p = 0.5004 with the lightest amplitude negative, which is a consistency check only, because their masses were built from the rule. For the quarks p = 0.456 (down-type) and 0.393 (up-type), and no choice of signs reaches \ensuremath{\frac{1}{2}}, so the principle as stated does not hold for them; colour is the named difference and why it changes the balance is open. Maximum uncertainty is itself a starting principle of the framework and is not derived from anything deeper. Scripts: \texttt{analysis/derive\_45\_and\_sqrt2\_2026-11-18.py}, which also shows that the grain factor \ensuremath{\surd}2 is not derived (five equally natural readings give five factors), and \texttt{analysis/matrix\_reading\_from\_rule4\_2026-11-18.py}.

\textbf{The matter--antimatter excess.} Not predicted; allowed in one narrow region (Section 8). With majoron dark matter the hot start that region needs fails, so the excess is carried by non-thermal leptogenesis (Lazarides and Shafi, 1991), with one number; the field that settles the Majorana mass cannot itself be that start. The cold start squeezes the numbers. The literal partners, the lightest near 10\textsuperscript{5} GeV, cannot make the excess from the lightest one's decays, whose asymmetry is capped near 10\textsuperscript{\ensuremath{-}11} (Davidson and Ibarra, 2002). In the region that works the perspective is tilted, the partners are near 3.6 \ensuremath{\times} 10\textsuperscript{4}, 1.7 \ensuremath{\times} 10\textsuperscript{10} and 10\textsuperscript{11} GeV, the source is the middle one, and the lightest one's later erasing is included. The start must hand about 0.2\% of its energy to that partner and reheat the universe to between 10\textsuperscript{6} and 1.6 \ensuremath{\times} 10\textsuperscript{9} GeV, so the lightest partner does meet the heat; the majoron stays cold because the settling field never re-enters it (hot majorons made by partner pairs stay under a few per cent below about 10\textsuperscript{9} GeV). Above 2.4 \ensuremath{\times} 10\textsuperscript{8} GeV the far-side weak pairs are also made and their decay coupling is needed. The same CMB limit that keeps the majoron cold caps the start's expansion rate at about 3 \ensuremath{\times} 10\textsuperscript{\ensuremath{-}5} f \ensuremath{\times} angle, that is 3 \ensuremath{\times} 10\textsuperscript{6} to 5 \ensuremath{\times} 10\textsuperscript{7} GeV for f between 10\textsuperscript{11} and 1.5 \ensuremath{\times} 10\textsuperscript{12} GeV, while handing energy to the middle partner needs a start field heavier than twice its mass, 3.4 \ensuremath{\times} 10\textsuperscript{10} GeV; the start field must therefore be 700 to 10\textsuperscript{4} times heavier than the expansion rate, which simple single-field starts do not give and a steep or two-field start does. The window is narrow; closing it would remove the lead, not the core. \emph{With the alternative start phase} the start's own field (the scalaron, 3.2 \ensuremath{\times} 10\textsuperscript{13} GeV) can hand energy to the heaviest partner (about 3\% of its decays at \ensuremath{\xi} = 0.12); at the largest asymmetry allowed (Davidson and Ibarra, 2002) this gives 1.0 \ensuremath{\times} 10\textsuperscript{\ensuremath{-}10} against 0.87 \ensuremath{\times} 10\textsuperscript{\ensuremath{-}10} needed, before erasure by the middle partners, which are in the heat. This is borderline, not a result; the excess stays a measured number.

\textbf{Two leads with numbers.} The quark triangle angle: if the quark link carries a quarter turn, the unitarity triangle is right-angled at \ensuremath{\alpha} (Antusch, King, Malinsky and Spinrath, 2010), and with the measured mixing the best fits give \ensuremath{\alpha} = 89 to 93\ensuremath{^\circ} and \ensuremath{\gamma} = 64.1 to 67.6\ensuremath{^\circ} (66.1 best; 67.4\ensuremath{^\circ} exactly right-angled), against \ensuremath{\gamma} = 66.4 (+2.7 \ensuremath{-}2.8)\ensuremath{^\circ} measured; the chance of such a match by luck is 24 to 80\%, so it is a lead. The direct measurement of \ensuremath{\alpha}, 84.1 (+3.7 \ensuremath{-}3.0)\ensuremath{^\circ} (Particle Data Group, 2026), sits 1.6\ensuremath{\sigma} below 90\ensuremath{^\circ}. The small neutrino splitting: the hint m\textsubscript{2}/m\textsubscript{3} = 3 \ensuremath{-} 2\ensuremath{\surd}2 = tan\textsuperscript{2} 22.5\ensuremath{^\circ}, with m\textsubscript{1} = 0.36 meV and the measured large splitting, gives \ensuremath{\Delta}m\textsuperscript{2}\textsubscript{21} = 7.37 \ensuremath{\pm} 0.08 \ensuremath{\times} 10\textsuperscript{\ensuremath{-}5} eV\textsuperscript{2} with JUNO's large splitting (2.509 \ensuremath{\times} 10\textsuperscript{\ensuremath{-}3} eV\textsuperscript{2}), against JUNO's measured 7.388 \ensuremath{\pm} 0.078 (JUNO Collaboration, 2026), within 0.2\ensuremath{\sigma}. Versions up to v1.2.2 gave 7.40 \ensuremath{\pm} 0.06: the relation is the same and has no free number, and the value moved in v1.3.2 (18 November 2026) only because the large-splitting input changed from the global fit's 2.519 \ensuremath{\times} 10\textsuperscript{\ensuremath{-}3} to JUNO's 2.509 \ensuremath{\times} 10\textsuperscript{\ensuremath{-}3} eV\textsuperscript{2} (29 September 2026). Both values lie within 0.2\ensuremath{\sigma} of JUNO's, so the change gained nothing. The lightest mass used here was itself obtained with the measured small splitting, but setting it to zero moves the result by only 0.01, so the reuse does not decide the agreement. The ratio 3 \ensuremath{-} 2\ensuremath{\surd}2 comes from a geometric picture, the neutrinos at the midpoint of one side of the mass triangle, which was proposed when the measured m\textsubscript{2}/m\textsubscript{3} (0.172--0.173) was already known; the chance of so close a match by luck was put at about 10\%. It is therefore a hint, not a rule, and JUNO's full precision will decide it.

\textbf{The exact right-angled triangle (lead, 25 November 2026).} Taking the right angle at \ensuremath{\alpha} as exact and choosing simple values for the other numbers, \ensuremath{\lambda} = 1/\ensuremath{\surd}20 = 0.22361, A = \ensuremath{\surd}(2/3), \ensuremath{\beta} = 22.5\ensuremath{^\circ} and \ensuremath{\gamma} = 67.5\ensuremath{^\circ}, gives an exactly unitary mixing matrix (Jarlskog invariant 2.94 \ensuremath{\times} 10\textsuperscript{\ensuremath{-}5}) whose seven directly measured quantities all lie within 1.6\ensuremath{\sigma} of measurement: \textbar V\textsubscript{us}\textbar{} \ensuremath{-}0.8\ensuremath{\sigma}, \textbar V\textsubscript{cb}\textbar{} \ensuremath{-}0.2\ensuremath{\sigma}, \textbar V\textsubscript{ub}\textbar{} \ensuremath{-}1.2\ensuremath{\sigma}, sin 2\ensuremath{\beta} = 1/\ensuremath{\surd}2 \ensuremath{-}0.2\ensuremath{\sigma}, \ensuremath{\alpha} +1.6\ensuremath{\sigma}, \ensuremath{\gamma} +0.4\ensuremath{\sigma} and \textbar V\textsubscript{td}/V\textsubscript{ts}\textbar{} +1.0\ensuremath{\sigma}, with \ensuremath{\chi}\textsuperscript{2} = 6.0 for the seven (Particle Data Group, 2026, and 2024 for the quantities not recalibrated here). The four numbers were chosen, not derived, so the \ensuremath{\chi}\textsuperscript{2} is not evidence and the form is a lead. Because it is exactly unitary it shares the first-row prediction of Section 9: \textbar V\textsubscript{ud}\textbar{} = 0.97467, 3.1\ensuremath{\sigma} above today's 0.97367 \ensuremath{\pm} 0.00032. A firm sin 2\ensuremath{\beta} away from 0.707, \ensuremath{\alpha} firmly away from 90\ensuremath{^\circ}, or a confirmed first-row shortfall would end it (Belle II, LHCb). A 3 \ensuremath{\times} 3 down-quark mass matrix with set weights built to reproduce this triangle does not (Appendix A.7). Script: \texttt{analysis/quark\_triangle\_and\_press\_2026-11-25.py}.

\textbf{A charge rule for the far-side weight (tested, does not hold).} The weight rule of the far-side weak pairs, the meeting energy squared divided by the grain, would follow from a discrete six-step charge (\ensuremath{\mathbb{Z}}\textsubscript{6}) carried by the spin-0 field \ensuremath{\Phi} that gives the coloured far-side states their mass. If the two halves of each coloured state differ by one step and those of each weak pair by two, the coloured states take their mass from \ensuremath{\Phi} directly and the weak pairs only from \ensuremath{\Phi}\textsuperscript{2} divided by the grain, which gives 2.45 \ensuremath{\times} 10\textsuperscript{8} GeV; the leptons then carry three steps, which forbids every operator that would turn a proton into a lepton. But a discrete charge that gravity respects must be free of anomalies (Krauss and Wilczek, 1989; Ibáñez and Ross, 1991), and this one is not. Counted field by field over three generations, the colour and weak sums are 3 and 15, not multiples of 6; with the half-strength form of the condition sometimes used, those two pass but the gravitational sum, 21, fails, and the Standard Model's own fields pass without the far side, so the far side is not required either way. No discrete charge of any order works (\ensuremath{\mathbb{Z}}\textsubscript{2} to \ensuremath{\mathbb{Z}}\textsubscript{18} were searched with every flavour-blind assignment): the quarks' own Higgs links make their colour sum vanish, so the colour sum is three times the coloured states' step, \ensuremath{\pm}3, which vanishes only for \ensuremath{\mathbb{Z}}\textsubscript{3}, and there \ensuremath{\Phi}\textsuperscript{2} carries the charge of \ensuremath{\Phi}*, so the weak pairs would take their mass at the same step as the coloured ones. The weight rule therefore stays a fitted choice (Appendix A.2, A.6). Script: \texttt{analysis/farside\_charge\_rule\_2026-11-25.py}.

\textbf{Gravity's strength and dark matter.} Gravity's strength and the three force strengths are one open problem, the strengths at the grain of space (Section 4). Dark matter is read as partly perspective (Section 5) and partly the majoron (Section 8); its mass and amount are open. A gravitational breaking of the kept spot does not fix the mass: counted correctly it gives about 1\% of the dark matter at the settling value used here, and some power of the breaking lands for almost any settling value, so it is not evidence. The majoron as dark matter makes no directly observable number; it is a candidate that fits every limit, not a tested result.

\textbf{The Higgs side, the generations and the proton's mass.} Why the far side settles, and at 246 GeV, is not derived; nor is how the Higgs link distinguishes the three placements, which is the single question behind the heavy quark masses, the mixing tilt and the twist. Whether the entanglement that gives gravity and the colour balance that gives the proton's mass are one relationship is open. The generations are unequal in a simple way: the heaviest member carries 99.62\% of the up-type masses, 97.94\% of the down-type and 94.36\% of the charged-lepton masses (quark masses run to 162.5 GeV); this describes the measured masses and is not derived.

\textbf{The seeds of galaxies and quantum probabilities.} On the main line the size and tilt of the first ripples are measured and not derived, and what laid them down is open. The alternative start phase lays them down: the size stays a measured number and the tilt is predicted (Section 9), at 0.963, which today's combined data disfavour. The probabilities of quantum mechanics (Born, 1926) are inherited; v1.4.0 adopts the published results that derive them from the space of states (Gleason, 1957) and from information rules with entanglement first (Chiribella, D'Ariano and Perinotti, 2011), which the framework does not derive by itself.

\textbf{Summary of status.} Derived from the rules: the octonions and the strong force with three colours; which other forces exist (by theorem from the perspective principle); for leptons, the 45\ensuremath{^\circ} balance, given the starting principle of maximum uncertainty between self and relationship, which was chosen knowing Koide's relation (Appendix A.2); with stated readings, every particle of a generation (16 states, the up and electron partners included) with its measured charge, one-handed matter with no mirror states and one i throughout, which mass links exist, and the same forces for all generations; at least three generations; the form of gravity; mass as balance energy. Measured rules, not derived: the twist, and the quark balance angles. Passing strict tests: the down quark, the up quark (with the decay value of Q) and the anomaly sums. Reproduced with fitted numbers: fission (166 against 173.3 MeV). Consistent with today's bound but not yet tested: the neutrino mass total. Met by construction, not tests: the compact-star masses, which set the pairing gap band (Appendix A.6). Taken from measurement: the three force strengths, the Higgs balance point, the lepton scale and twist, the top, bottom, charm and strange masses, quark mixing, the proton's two containment numbers, the strength of gravity and the amount of dark energy; strictly counted, one lepton number and one quark number fewer than the Standard Model, with the neutrino masses tied to the leptons. Open, each needing a new idea: the strengths at the grain of space (the force strengths, 8.9\% off at the Z mass for the Standard Model alone and closed by the far-side weak pairs if the strong coupling is 0.1178, which awaits measurement, together with gravity's strength); why the Higgs side settles, and at 246 GeV; how the Higgs link distinguishes the three placements, which holds the heavy quark masses, the mixing tilt and the twist (2/9, the same number as the electron's mass); why three generations and no more (the matrix's three sizes are a hint only); the matter excess, allowed with one number; dark matter's mass and amount, with the majoron as the candidate; and whether gravity's entanglement and the proton's balance are one relationship. Joining parts the theory already has brings the strengths, the heavy quarks with the tilt, the matter excess and dark matter each down to one missing rule for one number; a free number set to make a match is not a derivation, so none of them passes.

\textbf{Established results reached by a second route.} Confirmed exactly, with no measured number: which forces exist, the charges of a generation, neutral atoms and a colour-neutral proton, and \ensuremath{\rho} = 1 (Glashow, 1961; Weinberg, 1967; Salam, 1968; Veltman, 1977). Confirmed within measurement: atomic bindings (Bohr, 1913, in relational form), the tau lifetime (Fermi, 1934), the top width, the tau mass (Koide, 1983), the pion decay ratio (Marciano and Sirlin, 1993), the Cabibbo angle (Gatto, Sartori and Tonin, 1968), the down and up quarks, and the proton's mass as mostly binding energy (Ji, 1995; Yang et al., 2018). Reproduced with fitted numbers, as in the standard treatments: light baryons and nuclear binding; and the heaviest compact stars, whose masses fix the pairing gap. Holding in part: the galactic acceleration scale (Milgrom, 1983) and the top fixed point (Pendleton and Ross, 1981; Hill, 1981). Failing: the common strength of the forces (Georgi, Quinn and Weinberg, 1974) and the Higgs mass from asymptotic safety (Shaposhnikov and Wetterich, 2010).

\hypertarget{conclusion}{%
\subsection{11. Conclusion}\label{conclusion}}

A single relational picture, in which space is a full medium, structures are set by containment, and entanglement comes first, reproduces gravity, cosmic expansion, the proton's containment pressure, the masses of the heaviest compact stars and the charged-lepton masses, and, through two quark-mass relations, the down and up quarks. Required to have no gaps, its inner space yields the strong force with three colours, and its relationship matrix, seen from one perspective, gives which other forces exist, their charges and one-handed matter, each with a stated reading. It does this with about 28 measured numbers, against about 34 for the Standard Model with gravity and cosmology counted the same way (the 28 assumes every pending measurement agrees; counting only savings already confirmed it is 32, plus six stated choices; Appendix A). It makes predictions that experiments will test, among them the neutrino mass ordering and scale, the quark pairing gap and the strong coupling, and one, neutrinoless double beta decay at 1.3 to 4.1 meV, that lies beyond present reach. Read three at a time, its states are every particle of a generation with the measured charges. Its binding forces are exactly those blind to the observer's facing, giving a colour-neutral proton and a neutral atom. Where others established a result first, from the W and Z mass relation to atomic binding and nuclear masses, it reaches the same numbers by a second route, and says plainly where that route is fitted or fails. Dark matter's amount is measured; it is read as partly perspective, and the matrix allows one candidate for the rest, a cold, chargeless wave, the majoron, as a lead. The force strengths with gravity's, the settling of the Higgs side, the heavy quark masses and quark mixing, the lepton twist, the upper limit of three generations, the seeds of galaxies and what dark matter is remain open, as they are for all of physics, and the matter excess is allowed but not predicted. The framework does not ask any established result to be given up. It asks that each be seen as one view of the same universe.

In v1.4.0 the same framework is restated as five rules and one bet: one matrix of relationships read in three roles (held states, links between them, the link across each), a single balance of pull and resistance from the proton to the Local Group, maximum uncertainty for the leptons, nothing extra, and a smooth start. The stage of space and time and the rules of quantum theory are no longer borrowed whole. A stretching start driven by the stiffness of space is kept as the alternative, with two predictions, r = 0.0040 and n\textsubscript{s} = 0.963; the main line follows today's combined data, which disfavour that tilt. The measured numbers that remain are almost all link strengths, and they are where the work goes next.

\hypertarget{appendix-a.-ledger-of-inputs-and-predictions}{%
\subsection{Appendix A. Ledger of inputs and predictions}\label{appendix-a.-ledger-of-inputs-and-predictions}}

This appendix adds no new result. It gathers in one place what the sections above take from measurement, what they choose with today's data in view, and what they predict, so that a reader can see at once which agreements were possible only after the fact. The count, and the break tests for the results that move it, are rerun by \texttt{analysis/input\_ledger\_2026-11-18.py} (main line) and \texttt{analysis/input\_ledger\_2026-11-19.py} (with the alternative start), and the full table, with a column saying whether another published idea gives the same number, is \texttt{docs/ledger/TRPG\_Ledger.md} in the code archive.

\textbf{Labels.} \emph{Input}: taken from measurement. \emph{Fitted}: a choice made with today's data in view. \emph{Derived}: follows from the rules and inputs and agrees with a measurement already made. \emph{Predicted}: follows from the rules and inputs, but the deciding measurement has not returned; it is not 100\% until it does. \emph{Lead}: a relation with one named missing rule, outside the core. \emph{Effective}: an effective number of the strong force, counted on neither side.

\textbf{A.0 The nine criteria.} (1) Reduces to Einstein's gravity at large scales: met, from entanglement through Jacobson's argument. (2) Reduces to quantum theory and the Standard Model, with every charge, force group and generation: partly met, only the force strengths remain. (3) Conservation laws follow from symmetries: met. (4) Matches the cosmos: met with measured inputs. (5) Fewer free numbers than results: partly met, the heavy quark links await a rule. (6) Every equation derivable, units consistent and code that reruns: met; an outside review found unit errors in the earlier written action, corrected in Section 3 of this version without changing any number. (7) At least one prediction that could fail: met. (8) Internally consistent: met; the same review found the far side used both as a spin-0 and a spin-\ensuremath{\frac{1}{2}} object, resolved in Sections 3 and 8 by the E\textsubscript{6} reading. (9) One written action: met, as a fundamental action with an effective one for hadrons (Section 3).

\textbf{A.1 The headline prediction.} Normal neutrino ordering, with m\textsubscript{1} = 0.36 meV, m\textsubscript{2} = 8.6 to 8.7 meV and m\textsubscript{3} = 50.1 to 50.3 meV, summing to 0.0591 to 0.0594 eV (Section 9; the range is JUNO 2026 against the NuFIT 6.0 global fit). It uses only the 45\ensuremath{^\circ} rule that gives the tau mass, applied with one negative root, and the two measured splittings. The same relation for neutrinos was proposed by Brannen (2006); the framework supplies a reason for it but not a different number, so this test can falsify both and cannot separate them. JUNO is expected to decide the ordering at 3\ensuremath{\sigma} after about six and a half years of running (JUNO Collaboration, 2024), and DESI, Euclid and the next CMB experiments approach the needed sensitivity to the sum. Inverted ordering, or a sum measured clearly different from 0.059 eV, would falsify it.

\textbf{A.2 The count.} Counted like for like, the Standard Model with gravity and cosmology takes 34 free numbers. The framework removes six: the tau mass (Koide, 1983; 1776.969 MeV against 1776.93 \ensuremath{\pm} 0.09, +0.4\ensuremath{\sigma}) and the down quark (Gatto, Sartori and Tonin, 1968; m\textsubscript{s}/m\textsubscript{d} = 19.9 against 19.8 \ensuremath{\pm} 0.3, +0.3\ensuremath{\sigma}), which measurement has confirmed; and the Higgs mass, the lightest neutrino mass, \ensuremath{\theta}\textsubscript{12} given \ensuremath{\theta}\textsubscript{13} and the neutrino CP phase, which wait on measurements not yet made. The count is therefore 32 today and 28 if every pending measurement agrees; the figure of about 28 used above is that ceiling. If the strong coupling also lands at 0.1178, one force strength is saved as well, but only once the Higgs mass lands too, since both rest on the same fitted grain. Six choices were made with today's data in view: the grain as the Planck energy divided by \ensuremath{\surd}2 (one of five equally natural readings); the weight rule of the far-side weak pairs; the coloured far-side states placed where the strengths cross, with the strengths taken as equal there; TM1 rather than another mixing pattern; the sign pattern (\ensuremath{-}, +, +) of the neutrino roots; and maximum uncertainty between self and relationship as the reason for the 45\ensuremath{^\circ} rule, a starting principle chosen knowing that Koide's relation holds for the charged leptons. The containment pressure B and the stiffness of the containment field are effective numbers of the strong force and are counted on neither side. The far side's costs: the weight of the weak pairs is set by the weight rule (one of the six choices above), the weight of the D states by the meeting of the strengths, and their decay coupling is one more number only if they were ever produced (Section 10). The numbers of the leads (the majoron's breaking value and starting angle, the non-thermal start, the strength of the D fold) belong to the leads and are not in the count; a reader who charges them to the framework should add at least four. \emph{v1.4.0, alternative start only:} the start phase would turn the tilt of the first ripples into a prediction (saved only if it lands, so the ceiling stays 28), makes the far-side weak pairs' decay coupling a required number (one more today; once the tilt lands the two cancel and the ceiling stays 28), and adds a seventh choice made with today's data in view, the Higgs field's coupling to curvature \ensuremath{\xi} \ensuremath{\approx} 0.06--0.15 (the lower end for the Higgs's safety during the start, the upper band to keep the majoron cold). The majoron's starting angle would stop being a start condition. With the alternative the count is 33 today and 28 if every pending measurement agrees, plus seven choices (\texttt{analysis/input\_ledger\_2026-11-19.py}); on the main line it stays 32 and 28, plus six choices (\texttt{analysis/input\_ledger\_2026-11-18.py}).

\textbf{A.3 Labels tightened in gathering the ledger.} The lepton masses come from two numbers, the scale and the twist, with the 45\ensuremath{^\circ} rule, not from one. The Cabibbo angle from the masses and the down quark from the Cabibbo angle are one relation and are counted once. Light bending, the black-hole ringdown and the age of the universe are standard results reproduced from measured inputs and are not counted as predictions. The compact-star mass window of 2.0 to 2.15 solar masses is the observation the pairing gap must respect; the gap is the prediction. The Higgs-mass rule of Shaposhnikov and Wetterich fails with the Standard Model's states alone (2.6\ensuremath{\sigma}, Section 9) and passes at 1.4\ensuremath{\sigma} with the far-side states (Section 9), while needing a top mass 3.7\ensuremath{\sigma} below today's on the top mass alone.

\textbf{A.4 The framework is false if} any of the following is measured. Inverted neutrino ordering (JUNO, early 2030s; DUNE and Hyper-Kamiokande). A neutrino mass sum that excludes 0.0591--0.0594 eV at 3\ensuremath{\sigma} (DESI, Euclid). The strong coupling more than 3\ensuremath{\sigma} from 0.1178 in a lattice determination at about 0.3\%. A top-quark mass that puts the Higgs rule more than 3\ensuremath{\sigma} off. The neutrino CP phase firmly outside 253--292\ensuremath{^\circ} (DUNE, Hyper-Kamiokande), or sin\textsuperscript{2}\ensuremath{\theta}\textsubscript{12} more than 3\ensuremath{\sigma} from the TM1 value 0.318 (JUNO; today 2.3\ensuremath{\sigma}). A quark pairing gap firmly outside the band 4\ensuremath{\Delta}\textsuperscript{2} \ensuremath{-} m\textsubscript{s}\textsuperscript{2} = 1.7--2.8 \ensuremath{\times} 10\textsuperscript{4} MeV\textsuperscript{2} (82--97 MeV at m\textsubscript{s} = 100 MeV), or a compact star clearly above about 2.15 solar masses. Any proton decay (Hyper-Kamiokande). A fourth generation, a new force carrier, mirror families, a second Higgs pair or a coloured spin-0 particle. A confirmed shortfall in the first row of the quark-mixing matrix. Neutrinos proven to be Dirac particles. The tau mass firmly away from 1776.969 MeV (Belle II, BESIII). \emph{Added in v1.4.0, for the alternative start only:} no primordial gravitational waves at r \ensuremath{\approx} 0.004 once Simons Observatory and LiteBIRD reach about 0.001, or a CMB-only tilt above about 0.972 at 3\ensuremath{\sigma} (today it is 2.1\ensuremath{\sigma} low against the CMB with ACT and SPT and 3.2--4.2\ensuremath{\sigma} low once DESI is combined). Four leads would fail on their own without touching the core: primordial gravitational waves at r far above 10\textsuperscript{\ensuremath{-}11} under the cold start of the main line (the alternative start predicts them instead), the small neutrino splitting more than 3\ensuremath{\sigma} from 7.37 \ensuremath{\pm} 0.08 \ensuremath{\times} 10\textsuperscript{\ensuremath{-}5} eV\textsuperscript{2} once JUNO's own error is combined with the prediction's (the \ensuremath{\pm} comes from the large splitting), which at JUNO's full precision means outside about 7.13 to 7.62 (versions up to v1.3.2 printed 7.32 to 7.48, and a draft of v1.3.3 printed 7.29 to 7.45; both were one-standard-deviation bands, too narrow to count as a failure) (JUNO's 2026 value, 7.388 \ensuremath{\pm} 0.078, sits within 0.2\ensuremath{\sigma} of the predicted 7.37), and the quark triangle angle \ensuremath{\gamma} firmly outside 64.1--67.6\ensuremath{^\circ} (LHCb, Belle II), which follows from the lead \ensuremath{\alpha} = 90\ensuremath{^\circ} (compare Antusch, King, Malinsky and Spinrath, 2010), and the exact right-angled form of Section 10 if sin 2\ensuremath{\beta} is measured firmly away from 1/\ensuremath{\surd}2 = 0.707 (Belle II, LHCb).

\textbf{A.6 Tests set after an earlier attempt failed (audit of 18 November 2026).} Every test in this paper was checked for one question: was a number, a form or a model chosen after an earlier attempt missed, so that the result could not have failed? Those found are listed here with what they mean. None of them is counted as evidence; each is either a fitted choice already in A.2 or is now labelled so in the text.

\begin{enumerate}
\def\labelenumi{\arabic{enumi}.}
\tightlist
\item
  \emph{The grain of space, M\textsubscript{Pl}/\ensuremath{\surd}2.} Five readings of the grain were written down before running (\texttt{analysis/leak\_establish\_three\_loop\_2026-11-16.py}). The plain Planck energy also passed (+1.6\ensuremath{\sigma} in the weak mixing); \ensuremath{\surd}2 fitted best (\ensuremath{-}0.5\ensuremath{\sigma}) and was kept. Choosing the best of five is a fit, so \ensuremath{\alpha}\textsubscript{s} = 0.1178 is a prediction only against later measurements; with the plain Planck energy the framework would predict 0.1196. Three-loop running was used for the final number after the two-loop result was known; the two differ by less than 0.0002.
\item
  \emph{The weight of the far-side weak pairs.} The strengths require 2.2 \ensuremath{\times} 10\textsuperscript{8} GeV; the rule ``meeting energy squared \ensuremath{\div} grain'' (2.4 \ensuremath{\times} 10\textsuperscript{8} GeV) was found afterwards. Fitted choice 2.
\item
  \emph{Equal strengths at the crossing.} After the strengths were shown not to meet at the grain for any placement (v1.3.1), the boundary condition ``equal where the 27 is complete'' was added. Fitted choice 3.
\item
  \emph{The compact stars.} With no pairing the stars fail (1.54--1.80 solar masses). The pairing gap was then introduced, and its band has been re-set each time the model changed: 70--90 MeV (14 November), 80--100 and 82--97 MeV (finer scans), and 88--109 MeV for a running strange mass in the self-bound strange-star model adopted after outside review (v1.3.0--v1.3.1). The band is by construction the one that fits the observed stars, so the star masses are not evidence; only an independent determination of the gap tests it, and the Bodmer--Witten margin left by the band is thin (5--7 MeV).
\item
  \emph{The Higgs-mass rule.} It fails with the Standard Model's states alone and passes with the far side, whose placement was already fixed by the strengths and the fitted grain. Of the rules borrowed from other theories and tested on 17 November, it is the only one that set a number and passed; the others that set a number failed (the loop-quantum-gravity grain, the top mass from the bottom, Nelson--Barr), so choosing it is a selection among several. Its pass now also rests on the 0.5 GeV allowance for the meaning of the top mass: without it the gap is 2.1\ensuremath{\sigma} with the 2026 values, outside the pass line of 2\ensuremath{\sigma} fixed before the test.
\item
  \emph{TM1 mixing} was chosen after tribimaximal and TM2 were excluded, and \emph{the sign pattern (\ensuremath{-}, +, +)} because the all-positive one has no solution. Fitted choices 4 and 5.
\item
  \emph{Maximum uncertainty between self and relationship} was adopted knowing Koide's relation holds. Fitted choice 6; the tau mass is derived from a principle chosen with data in view.
\item
  \emph{The small neutrino splitting.} The ratio 3 \ensuremath{-} 2\ensuremath{\surd}2 (the midpoint of a side of the mass triangle) was proposed when the measured m\textsubscript{2}/m\textsubscript{3} was already known. A hint, outside the core.
\item
  \emph{The Cabibbo relation's form.} The sine form passes (+0.3\ensuremath{\sigma}) and the tangent form does not (\ensuremath{-}2.7\ensuremath{\sigma}); the sine form was chosen with both in view. This is a seventh choice made with data in view; it is not added to the count of six because the down-quark saving would be lost with the other form rather than a number added.
\item
  \emph{The up quark's Q.} The decay value Q = 22.1 \ensuremath{\pm} 0.7 gives a pass (0.4\ensuremath{\sigma}, with the uncertainty of Q included); the lattice value Q = 23.4 \ensuremath{\pm} 0.5 gives 2.0\ensuremath{\sigma}, on the line between a pass and a partial. The decay value was preferred to avoid lattice input. Because Q is known to only 3\%, the test is loose either way.
\item
  \emph{Earlier fits since withdrawn,} kept on record: the golden-ratio mass ladder (it fits random masses as well); the bottom-mass quark chain (it fed its own charm into the strange step); the 4.4 \ensuremath{\times} 10\textsuperscript{\ensuremath{-}10} matter excess (solver-dependent); the link pattern 3/6, 4/6, 5/6 (7\ensuremath{\sigma}).
\end{enumerate}

Not affected: the strong force and three colours, which forces exist, the charges of a generation, the anomaly sums, \ensuremath{\rho} = 1, gravity's form, the tau lifetime, the age of the universe and the light bending. The near agreements already marked in the text as not evidence (gravity's count 43.6, the dark-energy size, the galactic acceleration scale cH\textsubscript{0}/2\ensuremath{\pi}, whose factor 2\ensuremath{\pi} was taken from Milgrom's known coincidence) stay so.

\textbf{A.7 Tests run after 18 November 2026 that failed.} These were tried between 19 and 25 November and rechecked with independent code before this version. None is part of the framework. Several of them had been checked more than once, in different ways, before an independent recheck found the mistake (a miscounted anomaly sum, measured inputs that had drifted from the cited values, and a target moved after a failure). They are kept here so that they are not repeated, and to show how every result in this paper is tested. Scripts: \texttt{analysis/farside\_charge\_rule\_2026-11-25.py}, \texttt{analysis/quark\_triangle\_and\_press\_2026-11-25.py}, \texttt{analysis/cfl\_gap\_and\_tides\_2026-11-25.py} and \texttt{analysis/refinement\_checks\_2026-11-25.py}.

\begin{enumerate}
\def\labelenumi{\arabic{enumi}.}
\tightlist
\item
  \emph{A six-step charge for the far side} (Section 10). It is anomalous, and no discrete charge of any order gives the weight rule without an anomaly (\texttt{analysis/refinement\_checks\_2026-11-25.py}). It therefore does not select three generations either: with the sums counted fully, no number of generations from one to five passes, and six, which does, ends the asymptotic freedom of the strong force.
\item
  \emph{Removing that charge's walls by gravity.} If the charge were broken by gravity at its first allowed power (\ensuremath{\Phi}\textsuperscript{5} divided by the grain), the walls it leaves would vanish near 1.3 \ensuremath{\times} 10\textsuperscript{13} GeV, before they could dominate at 1.1 \ensuremath{\times} 10\textsuperscript{11} GeV. But a charge that gravity breaks gives no protection, so this cannot be combined with item 1. (\ensuremath{\Phi}\textsuperscript{6} carries no charge and breaks nothing.)
\item
  \emph{A very light majoron as all of the dark matter, without inflation} (settling value 2.1 \ensuremath{\times} 10\textsuperscript{16} GeV, mass 3 \ensuremath{\times} 10\textsuperscript{\ensuremath{-}19} to 10\textsuperscript{\ensuremath{-}16} eV). The amount comes out right for starting angles of 0.1 to 0.6 (0.2 to 1.0 with another standard estimate; Hui, Ostriker, Tremaine and Witten, 2017). But without inflation the starting angle differs between regions that were never in contact, so one angle cannot be chosen; a mass this small from gravity needs the breaking to begin near the 36th to 38th power, not the 5th of item 2; and black holes would lose their spin to such a field, which spin measurements of supermassive black holes already largely exclude from 10\textsuperscript{\ensuremath{-}22} to 10\textsuperscript{\ensuremath{-}17} eV (Zu and others, 2020; with the modelling caveats of Hoof and Marsh, 2024). The majoron lead of Section 8, with settling values of 10\textsuperscript{11} to 1.5 \ensuremath{\times} 10\textsuperscript{12} GeV, is not affected.
\item
  \emph{A larger pairing gap, 90 to 125 MeV,} chosen to reach the 2.35 solar masses of PSR J0952\ensuremath{-}0607. The stars reach 2.20 to 2.61 solar masses, but their tidal deformability at 1.4 solar masses is 652 to 1124, above the GW170817 limit of about 580 (Abbott and others, 2018). Placing the 1.4-solar-mass stars on a separate, ordinary branch is the published two-family picture (Drago, Lavagno and Pagliara, 2014), not a result of this framework. If up and down quarks pair with the same gap, ordinary nuclei would also be unstable over much of that band. The band of Section 9 is unchanged.
\item
  \emph{The matter excess with the Dirac link aligned to the neutrino masses.} The alignment gives weak erasure (K\textsubscript{1} = m\textsubscript{1}/m\textasciitilde*\textasciitilde{} = 0.33) but makes the heavy partners' decay asymmetry vanish, so the whole excess rests on a starting rotation of the majoron field set by hand. This is not a derivation; Section 10 stands.
\item
  \emph{Dark energy from the containment pressure.} \ensuremath{\rho}\textsubscript{\ensuremath{\Lambda}} = B\textsuperscript{3/2} / {[}3 (16\ensuremath{\pi}\textsuperscript{2})\textsuperscript{2} M\textsubscript{Pl}\textsuperscript{2}{]} gives 2.13 to 2.23 meV against 2.24 meV measured, but it was the closest of ten forms tried, with the factor 16\ensuremath{\pi}\textsuperscript{2} put in by hand. It is a variant of Zel'dovich's relation (1967) and is not evidence.
\item
  \emph{A 3 \ensuremath{\times} 3 down-quark mass matrix with set weights} (2r\textsubscript{d}, r\textsubscript{s} + 4r\textsubscript{s}\textsuperscript{2}, \ensuremath{\surd}(r\textsubscript{d}r\textsubscript{s}), 2r\textsubscript{s} and \ensuremath{\surd}(2/3) \ensuremath{\surd}(r\textsubscript{d}r\textsubscript{s}) at a phase of \ensuremath{-}67.5\ensuremath{^\circ}). It gives \ensuremath{\chi}\textsuperscript{2} = 11.5 for the seven mixing quantities, with \textbar V\textsubscript{us}\textbar{} 2.2\ensuremath{\sigma} high, and 26 when the masses are taken at the Z mass instead of 162.5 GeV. The weights were chosen by trial.
\item
  \emph{Simple fractions for the tilt of the seeds of galaxies} (1 \ensuremath{-} 1/29 and similar). No equation produces them.
\item
  \emph{The Higgs mass as half of 246 GeV} (123.1 GeV, 1.7\% below the measurement), that is a self-coupling of exactly 1/8. Run at two loops with the 2026 values, the self-coupling is 0.1260 at the top mass, 0.1259 at v/\ensuremath{\surd}2 and 0.1192 at v; it passes 1/8 near 182 GeV, which is no scale of the framework, and 1/8 at v would need m\textsubscript{H} = 127.9 GeV against 125.13 \ensuremath{\pm} 0.11.
\item
  \emph{Three generations in one octonionic matrix.} Three 16-dimensional subspaces of its 24-dimensional off-diagonal part overlap pairwise in 8 dimensions and together span only 24, so one matrix cannot hold three independent generations. This rules out that embedding; it does not show that three copies are required.
\item
  \emph{Loop rungs for the masses.} Measured in steps of ln(16\ensuremath{\pi}\textsuperscript{2}) below the wall v/\ensuremath{\surd}2, the bottom, tau, charm, muon, strange, down, up and electron sit at 0.81, 0.91, 1.11, 1.47, 1.60, 2.19, 2.34 and 2.53; their mean distance from a whole step, 0.28, is matched by random masses over the same span 81\% of the time. The ordering of the generations is real; the steps are not.
\end{enumerate}

\textbf{A.8 Where each open problem stands (25 November 2026).} Every problem worked on in November is listed once, with what holds after the independent rechecks of 25 November and what is still missing. ``Holds'' means reproduced with independent code against today's measurements; it does not mean derived.

\begin{longtable}[]{@{}
  >{\raggedright\arraybackslash}p{(\columnwidth - 6\tabcolsep) * \real{0.2500}}
  >{\raggedright\arraybackslash}p{(\columnwidth - 6\tabcolsep) * \real{0.2500}}
  >{\raggedright\arraybackslash}p{(\columnwidth - 6\tabcolsep) * \real{0.2500}}
  >{\raggedright\arraybackslash}p{(\columnwidth - 6\tabcolsep) * \real{0.2500}}@{}}
\toprule\noalign{}
\begin{minipage}[b]{\linewidth}\raggedright
Problem
\end{minipage} & \begin{minipage}[b]{\linewidth}\raggedright
What holds
\end{minipage} & \begin{minipage}[b]{\linewidth}\raggedright
What failed (A.6, A.7)
\end{minipage} & \begin{minipage}[b]{\linewidth}\raggedright
Still missing
\end{minipage} \\
\midrule\noalign{}
\endhead
\bottomrule\noalign{}
\endlastfoot
Heavy quark masses and mixing & Down and up quarks pass strict tests (Section 7); the exact right-angled triangle fits all seven mixing quantities within 1.6\ensuremath{\sigma} (Section 10, lead); one far-side D fold gives all three angles within 19\% (lead) & 45\ensuremath{^\circ} cones for quarks; fixed ladders; a 3 \ensuremath{\times} 3 down-quark matrix with set weights & a rule for the heavy masses and the fold's strength \\
The lepton twist and the 45\ensuremath{^\circ} balance & 45\ensuremath{^\circ} balance derived for leptons from maximum uncertainty (fitted choice 6); tau mass +0.4\ensuremath{\sigma}; \ensuremath{\delta} = 2/9 gives m\textsubscript{e} 9.8 ppm low & rules that also fit the quarks (\ensuremath{\delta} = Q/3 and related forms) & the twist itself and its 10-ppm offset \\
Why three generations, and their mass ladder & Three placements of one shape; the heaviest member carries 99.6\%, 97.9\% and 94.4\% of each family's mass & a discrete charge that fixes the number; integer loop rungs; three generations in one octonionic matrix & why three, and why the steps differ \\
Absolute scales (246 GeV, the far-side weights, the force strengths) & M\textsubscript{L} = M\textsubscript{D}\textsuperscript{2}/E\textsubscript{grain} = 2.45 \ensuremath{\times} 10\textsuperscript{8} GeV matches what the strengths need (fitted choice 2); strong coupling 0.1178 and Higgs 127.2 \ensuremath{\pm} 1.5 GeV await measurement & the weight rule as a discrete symmetry (any \ensuremath{\mathbb{Z}}\textsubscript{N}); m\textsubscript{H} = v/2 & a derivation of v and of the meeting scale \\
Matter--antimatter excess & allowed through non-thermal leptogenesis with one start number (Section 10) & literal thermal leptogenesis; the Dirac link aligned to the neutrinos & a rule for the start number \\
Dark matter & the majoron as a cold wave with settling value 10\textsuperscript{11} to 1.5 \ensuremath{\times} 10\textsuperscript{12} GeV passes every limit (lead) & a very light majoron without inflation; a hot start & its mass and amount \\
Dark energy & unimodular form keeps the vacuum energy from weighing; \ensuremath{\surd}(grain \ensuremath{\times} horizon) gives 3.5 meV against 2.24 (not evidence) & a formula from the containment pressure, the closest of ten tried & its size \\
Seeds of galaxies & size and tilt measured, used as inputs; the alternative start phase predicts n\textsubscript{s} = 0.963 and r = 0.0040, fine against the CMB alone, 3.2--4.2\ensuremath{\sigma} low with DESI & simple fractions for the tilt & what laid them down on the main line \\
Compact stars & the band 82--97 MeV (m\textsubscript{s} = 100 MeV) reproduced by an independent solver: at B = 60.1 MeV/fm\textsuperscript{3} the heaviest star is 2.00--2.13 solar masses and \ensuremath{\Lambda}\textsubscript{1.4} = 459--575 & unpaired quark stars; a 90--125 MeV gap & an independent measurement of the gap \\
\end{longtable}

\textbf{A.5 The derivation tree.} \texttt{figures/derivation\_tree.png} in the code archive shows the chain from the rules (one full medium seen one perspective at a time; containment; maximum uncertainty between self and relationship; nothing extra; and, in v1.4.0, the smooth start, with the stretching phase as the alternative) and the bet (the far side is real) to each output, with the points where measured inputs and fitted choices enter.

\hypertarget{appendix-b.-the-particle-side-on-a-lattice-v1.4.2}{%
\subsection{Appendix B. The particle side on a lattice (v1.4.2)}\label{appendix-b.-the-particle-side-on-a-lattice-v1.4.2}}

\textbf{Question.} Can the framework's particles and forces be defined exactly on a grid of points, the way Craig (2026) has proposed for the Standard Model? A grid version makes a theory mathematically well defined at every finite spacing, without approximations. It derives no new numbers.

\textbf{Result in one line.} If Craig's theorem holds, the framework's whole particle side has an exact, gauge-invariant definition on a finite grid. That covers all 27 states per generation, the Higgs, the mass links, the Majorana mass and the majoron. Craig's theorem already covers all 27 states; two lemmas, proved below and checked numerically, give the far-side pair an explicit phase-free measure and add the global symmetry, the scalars and the mass terms. The computable steps of Craig's own proof were rechecked independently (B.7). Reflection positivity and the zero-spacing limit stay open, exactly as for the Standard Model.

\textbf{B.1 Craig's construction.} Craig (2026) puts one-handed fermions on a lattice using overlap fermions (Narayanan and Neuberger, 1994; Neuberger, 1998). These keep an exact remnant of chiral symmetry (Ginsparg and Wilson, 1982) and avoid the mirror-image doublers that a naive grid produces (Nielsen and Ninomiya, 1981). With H = \ensuremath{\gamma}\ensuremath{_{5}}(D\_W \ensuremath{-} 1), a left-handed field obeys \ensuremath{\hat{P}}\ensuremath{_{-}}\ensuremath{\psi} = \ensuremath{\psi}, where \ensuremath{\hat{P}}\ensuremath{_{-}} = (1 + sign H)/2 depends on the gauge field. The antifield obeys \ensuremath{\bar{\ensuremath{\psi}}}P\ensuremath{_{+}} = \ensuremath{\bar{\ensuremath{\psi}}} with the fixed P\ensuremath{_{+}} = (1 + \ensuremath{\gamma}\ensuremath{_{5}})/2 (Lüscher, 1999).

The open problem since Lüscher (1999), and for the electroweak sector Kadoh and Kikukawa (2008), was to fix the phase of the fermion measure gauge-invariantly for non-abelian groups. Craig fixes it by comparing each gauge field with smooth copies of itself on grids of spacing 1/3, 1/9, and so on, and summing the phase changes. The Standard Model is embedded in SU(5), with left-handed fermions in 10 \ensuremath{\oplus} \ensuremath{\bar{5}} per generation. The construction needs three things:

\begin{itemize}
\tightlist
\item
  the gauge anomalies cancel;
\item
  under the SU(2) acting on two of the five colours, half-integer isospins appear an even number of times;
\item
  the plaquettes obey a bound (admissibility).
\end{itemize}

The result is a smooth, gauge-invariant measure, fixed up to a constant phase on each topological sector. The paper was prepared with the aid of OpenAI's GPT-6 Astra, which developed the detailed proof. Reflection positivity at finite spacing is not yet established. M. Golterman, A. Kronfeld and Y. Shamir examined the proof before release, which is not journal review.

\textbf{B.2 The framework's content.} On the main line each generation is the 27 of Section 3 (Gürsey, Ramond and Sikivie, 1976). In SU(5) language it reads

10 \ensuremath{\oplus} \ensuremath{\bar{5}} \ensuremath{\oplus} 1 (the known states and the hidden partner N) and 5 \ensuremath{\oplus} \ensuremath{\bar{5}} \ensuremath{\oplus} 1 (the far side and the further neutral state).

The forces are SU(3) \ensuremath{\times} SU(2) \ensuremath{\times} U(1) only. The scalars are the Higgs doublet and the complex field whose phase is the majoron. All anomaly sums vanish for the 27, the SU(5) cubic sum is 1 \ensuremath{-} 1 + 1 \ensuremath{-} 1 = 0, and Craig's even-count condition passes: 4 doublets in the known 16 states, 6 in the full 27 (\texttt{analysis/lattice\_fit\_check\_2026-11-27.py}). Craig's covered families (his Proposition 11.1 and Table 3) include any content with zero cubic anomaly and an even count of half-integer isospins, combined with conjugates and gauge singlets, and explicitly any pair Q \ensuremath{\oplus} Q*. His theorem therefore covers the full 27 directly. An earlier draft of this appendix said the far-side pair lay outside it; the recheck in B.7 found that this was wrong. What his theorem does not address is the scalar and mass terms. Lemma A below gives the far-side pair an explicit form.

\textbf{B.3 Lemma A: an explicit, phase-free measure for the far-side pair.} \textbf{Statement.} In four dimensions, for any gauge field U and any representation r, the left-handed projector in the conjugate representation is the image of the one in r under a fixed antiunitary map:

\ensuremath{\hat{P}}\ensuremath{_{-}}{[}\ensuremath{\bar{r}}{]}(U) = A \ensuremath{\cdot} conj(\ensuremath{\hat{P}}\ensuremath{_{-}}{[}r{]}(U)) \ensuremath{\cdot} A\ensuremath{^{-}}\ensuremath{^{1}}, with A = \ensuremath{\gamma}\ensuremath{_{5}}C\ensuremath{^{-}}\ensuremath{^{1}},

where C is charge conjugation (C\ensuremath{\gamma}\_\ensuremath{\mu}C\ensuremath{^{-}}\ensuremath{^{1}} = \ensuremath{-}\ensuremath{\gamma}\_\ensuremath{\mu}\ensuremath{^{T}}).

\textbf{Proof.} A\ensuremath{\cdot}conj(\ensuremath{\gamma}\_\ensuremath{\mu})\ensuremath{\cdot}A\ensuremath{^{-}}\ensuremath{^{1}} = \ensuremath{\gamma}\_\ensuremath{\mu} and A\ensuremath{\cdot}conj(\ensuremath{\gamma}\ensuremath{_{5}})\ensuremath{\cdot}A\ensuremath{^{-}}\ensuremath{^{1}} = +\ensuremath{\gamma}\ensuremath{_{5}}, so A keeps handedness. Complex conjugation turns the links of r into those of \ensuremath{\bar{r}}. Hence A\ensuremath{\cdot}conj(D\_W{[}r{]})\ensuremath{\cdot}A\ensuremath{^{-}}\ensuremath{^{1}} = D\_W{[}\ensuremath{\bar{r}}{]}, and the same holds for H, for sign H and for \ensuremath{\hat{P}}\ensuremath{_{-}}. \ensuremath{\blacksquare}

\textbf{Consequence.} Take any orthonormal basis v\_j of the left-handed r fields and use v\ensuremath{^{\prime}}\_j = A\ensuremath{\cdot}conj(v\_j) for the \ensuremath{\bar{r}} fields. Changing the basis by a unitary Q multiplies the pair's measure by det Q\ensuremath{^{-}}\ensuremath{^{1}} \ensuremath{\cdot} conj(det Q)\ensuremath{^{-}}\ensuremath{^{1}} = \textbar det Q\textbar\ensuremath{^{-}}\ensuremath{^{2}} = 1. The pair's measure therefore needs no phase choice at all: it is canonical, gauge invariant (the map commutes with gauge transformations) and smooth. Its return phases cancel exactly, \ensuremath{\mathcal{T}}{[}\ensuremath{\bar{r}}{]} = conj(\ensuremath{\mathcal{T}}{[}r{]}). Its current is zero, so Craig's integrability condition W = \ensuremath{\mathcal{T}} holds for the sum because both sides multiply.

This gives a second route, which uses Craig's construction only for the Standard Model part: the full measure is Craig's measure for 3 \ensuremath{\times} (10 \ensuremath{\oplus} \ensuremath{\bar{5}}), times the canonical measure for 3 \ensuremath{\times} (5 \ensuremath{\oplus} \ensuremath{\bar{5}}), times constant measures for the six neutral states. A neutral state's projector does not depend on U, so its measure is constant. Which of the two identical \ensuremath{\bar{5}} copies is assigned to which factor changes the measure only by a constant phase per topological sector. That is the same freedom Craig's measure already has. Real representations on the lattice were treated earlier by Suzuki (2000), whom Craig cites.

\textbf{Checks} (\texttt{analysis/vectorlike\_pair\_measure\_check\_2026-11-27.py}):

\begin{itemize}
\tightlist
\item
  \textbf{The identity itself.} On a 3\ensuremath{^{4}} lattice with random SU(3) links it holds to 10\ensuremath{^{-}}\ensuremath{^{1}}\ensuremath{^{3}}.
\item
  \textbf{A closed loop of gauge fields.} The single-representation phase is \ensuremath{-}0.0053 rad, and the pair's phase is 1 \ensuremath{\times} 10\ensuremath{^{-}}\ensuremath{^{1}}\ensuremath{^{4}}.
\item
  \textbf{Control.} In two dimensions no handedness-keeping map exists, because there the conjugate of a left-mover is a left-mover and L(q) + L(\ensuremath{-}q) is anomalous. The check gives a nonzero phase (0.26 rad), as it must. The lemma is a property of four dimensions, not of the method of checking.
\end{itemize}

\textbf{B.4 Lemma B: the majoron's symmetry is exact on the grid.} \textbf{Statement.} Rotating each left-handed field by exp(i x\_f \ensuremath{\alpha}) multiplies the lattice measure by exp(\ensuremath{-}i\ensuremath{\alpha} \ensuremath{\Sigma} x\_f n\_f). Here n\_f is the index of the overlap operator in representation f.~By the lattice index theorem (Hasenfratz, Laliena and Niedermayer, 1998; Lüscher, 1998), \ensuremath{\Sigma} x\_f n\_f = 2n\ensuremath{_{3}}\ensuremath{\cdot}A(SU(3)\ensuremath{^{2}}X) + 2n\ensuremath{_{2}}\ensuremath{\cdot}A(SU(2)\ensuremath{^{2}}X) + n\ensuremath{_{1}}\ensuremath{\cdot}A(Y\ensuremath{^{2}}X) in every topological sector, with T = \ensuremath{\frac{1}{2}} for a fundamental (one fundamental field has index 1 = 2T in one unit of topology, hence the factor 2). A global charge is therefore exact on the lattice if and only if those three mixed sums vanish.

\textbf{Result.} For B \ensuremath{-} L on the 27 all three vanish, so the majoron's symmetry is exact on the grid in every sector. The only breaking is the deliberate grain-of-space term of Section 8. As a control, lepton number alone has SU(2)\ensuremath{^{2}}L = \ensuremath{\frac{1}{2}}. It is broken by weak-force topology, the sphaleron effect ('t Hooft, 1976), which the lattice reproduces as it should.

\textbf{Check} (\texttt{analysis/bl\_exact\_on\_lattice\_2026-11-27.py}). On an 8 \ensuremath{\times} 8 two-dimensional U(1) lattice the index is exactly q\ensuremath{\cdot}m for charges 1 to 3 and fluxes 0 to 2. Global charges with zero mixed sum leave the measure unchanged in every flux sector, and a charge with nonzero sum (6) changes it by 6 and 12 units.

\textbf{B.5 Corollary: the whole particle side at finite spacing.} The remaining terms are added as follows:

\begin{itemize}
\tightlist
\item
  \textbf{Gauge fields:} Lüscher's admissibility-respecting plaquette action for SU(3) \ensuremath{\times} SU(2) \ensuremath{\times} U(1), with the bound imposed on the SU(5)-embedded plaquettes.
\item
  \textbf{Scalars:} the Higgs doublet and the majoron field, each with nearest-neighbour covariant hopping and a quartic potential bounded below. These are standard lattice scalars, and on a finite grid their integrals converge.
\item
  \textbf{Fermion bilinears:} the Yukawa links of the known states and the lepton circulant, the mass links between D and the down quarks, the far-side masses, the Majorana mass of N and its coupling to the majoron field. In left-handed language these are all field--field or antifield--antifield bilinears.
\end{itemize}

A Berezin integral of exp(\ensuremath{-}S) against a fixed, gauge-invariant top form is a well-defined number for any polynomial action. So once the measure is fixed (by Craig's theorem for the whole 27, or by his theorem for 3 \ensuremath{\times} (10 \ensuremath{\oplus} \ensuremath{\bar{5}}) with Lemma A), the fermion integral is the Pfaffian of the mass-and-Yukawa matrix written in the measure's basis. It is gauge invariant whenever every term in S is. The Pfaffian's sign is fixed by the measure, so there is no square-root ambiguity. It is not positive in general, which matters for simulation but not for definition.

The containment field P is a classical description of the strong force (Friedberg and Lee, 1977) and is not part of the fundamental action. \textbf{Conclusion:} if Craig's theorem holds, every correlation function of the framework's particle side is defined and gauge invariant on every finite admissible lattice.

\textbf{B.6 Three further consequences.} \textbf{The strong-interaction angle fits the freedom left over.} Craig's measure is fixed up to one constant phase on each connected component of the admissible gauge fields. A phase of the form exp(i\ensuremath{\theta}n), with n the sector's winding number, is exactly the \ensuremath{\theta}-term, so the construction leaves room for the strong-interaction angle and neither fixes nor forbids it. A general choice of component phases is wider than this; requiring the phases to combine as a local term would reduce them to the \ensuremath{\theta}-terms, but that step is not proved here. This is consistent with the framework's choice to take the angle from measurement as one of its inputs (Section 3). It also means the lattice cannot be used to derive it.

\textbf{Baryon number changes only in steps of three.} Lemma B applied to baryon number B and lepton number L gives SU(3)\ensuremath{^{2}}B = SU(3)\ensuremath{^{2}}L = 0 and SU(2)\ensuremath{^{2}}B = SU(2)\ensuremath{^{2}}L = 3/2 for three generations. One unit of weak-force topology therefore changes B and L together by 2 \ensuremath{\times} 3/2 = 3, one per generation, and keeps B \ensuremath{-} L. That is the process of 't Hooft (1976). It cannot make a single proton decay, which needs \ensuremath{\Delta}B = 1, so the lattice reproduces the framework's ``no proton decay'' statement (Section 9) at the renormalisable level. The far-side states do not change the count, because they come in both handednesses (\texttt{analysis/extras\_theta\_baryon\_alternatives\_2026-11-27.py}).

\textbf{Only the main line is covered.} Both recorded alternatives place each generation's 16 in Spin(10):

\begin{itemize}
\tightlist
\item
  the gauged B \ensuremath{-} L of Section 8, reading (ii);
\item
  the 4 \ensuremath{\times} 4 picture of Pati and Salam.
\end{itemize}

Neither fits Craig's covered families. The one proposed lattice measure for the Spin(10) 16 (Kikukawa, as cited by Craig) has open smoothness and locality. So today only the main line has a complete lattice definition. That says which picture is better founded mathematically, not which is true.

\textbf{Already on the grid.} The framework uses lattice QCD for its containment numbers: the lightest glueball, 1.71 to 1.73 GeV, and the deconfinement temperature, 156.5 \ensuremath{\pm} 1.5 MeV. Those numbers come from the strong-force part of the same lattice theory, so they are the first place where its predictions meet the framework's balance picture. There, the balance gives 105 MeV against 156.5 MeV, a stated fail (Section 8).

\textbf{B.7 An independent check of Craig's proof.} Because everything here rests on Craig's theorem, we rechecked the parts of his proof that can be tested by direct computation, starting from the definitions in his paper (\texttt{analysis/craig\_proof\_recheck\_2026-11-27.py}, with its saved output). All 18 checks pass:

\begin{itemize}
\tightlist
\item
  \textbf{The colour-trace identity the cancellation rests on.} Ordered traces of up to four generators agree in 10 \ensuremath{\oplus} \ensuremath{\bar{5}} and in its conjugate to 10\ensuremath{^{-}}\ensuremath{^{1}}\ensuremath{^{3}}, and so they do for the full 27. At five generators they differ, so the cancellation stops exactly where the proof needs it to (the trace reductions follow van Ritbergen, Schellekens and Vermaseren, 1999). For an anomalous content (the 10 alone) they differ already at three. Behind this is the character identity \ensuremath{\chi}(R) \ensuremath{-} \ensuremath{\chi}(R*) = det(1 \ensuremath{-} g), with \ensuremath{\chi} the character, which holds to 10\ensuremath{^{-}}\ensuremath{^{1}}\ensuremath{^{4}} on random SU(5) elements.
\item
  \textbf{Craig's two tests (his Section 11.3).} Zero cubic sum and an even number of half-integer isospins: 4 for 10 \ensuremath{\oplus} \ensuremath{\bar{5}}, 6 for the 27, and his own failing SU(3) example fails as he states.
\item
  \textbf{The free-field facts.} The Wilson gap H\ensuremath{^{2}} \ensuremath{\ge} 1, the weights of the smooth copy, and no zero mode in the comparison of coarse and fine operators (least singular value 0.46 over 848 momenta and 21 interpolation values).
\item
  \textbf{The finite comparison identity (his Lemma 5.1).} On a random gapped family of projectors, the transported phase equals the ratio of return phases to 10\ensuremath{^{-}}\ensuremath{^{9}}, and the curvature identity db = \ensuremath{\Omega}\ensuremath{_{1}} \ensuremath{-} \ensuremath{\Omega}\ensuremath{_{0}} holds (0.015521 against 0.015520 by finite differences).
\item
  \textbf{The global step.} The spin and colour matrices of his global argument, the Kramers pairing that removes the last sign (determinant 1 to 10\ensuremath{^{-}}\ensuremath{^{1}}\ensuremath{^{5}}), the Standard Model embedding with its Z\ensuremath{_{6}} kernel, and the localisation-window identities.
\end{itemize}

Read and followed, but not re-derived: the global integrability argument and the reconstruction (his Sections 8 and 9) and the restriction to the Standard Model (his Section 11). Not checked: the analytic estimates in his appendices (locality, uniform constants, the order-h\ensuremath{^{8}} bounds, independence of the volume), the gap bound he cites from Hernández, Jansen and Lüscher (1999), and his Lean formalization and two-dimensional supplement, which we could not access. No error was found. This is an independent, partial check, not a second proof: every finite step we could compute holds, and the analytic estimates are left to specialists.

The check also corrected this appendix: the full 27 is in Craig's covered families (B.2), so Lemma A is a second, explicit route for the far-side pair. Lemma B and the corollary remain additions to his work.

\textbf{B.8 What stays open.} - \textbf{Craig's theorem itself.} It is a preprint with a Lean formalization, not yet a refereed result. Its computable steps pass our recheck (B.7), but its analytic estimates have not been checked by us. If it fails, this extension fails with it.
- \textbf{Reflection positivity.} It is needed to read the lattice theory as an ordinary quantum theory. It is open for the overlap chiral sector, for the framework exactly as for the Standard Model. The far-side pair and the scalars can be discretized in reflection-positive ways (Osterwalder and Seiler, 1978), but their couplings to the chiral sector have not been shown to keep it.
- \textbf{The zero-spacing limit.} U(1) and the scalar self-couplings are expected to be ``trivial'' in four dimensions (Aizenman, 1981; Fröhlich, 1982; Aizenman and Duminil-Copin, 2021). This is the same issue as for the Standard Model. The strong and weak forces remain asymptotically free with the far side included: b\ensuremath{_{3}} = 5 and b\ensuremath{_{2}} = 7/6, against 7 and 19/6 without it.
- \textbf{The numbers.} Nothing here derives a mass, strength or angle. Whether the lattice theory reproduces the framework's numbers is a separate, untested question.
- \textbf{Out of scope.} Gravity and the full-medium picture of space are not on the grid. The grid is a scaffold whose spacing goes to zero, so it says nothing for or against the physical grain of space. The relationship matrix fixes which states and charges exist; the lattice takes them as given, so it neither tests nor objects to that derivation.

\textbf{Verdict.} The finite-spacing definition is complete for the framework's particle side, conditional on Craig's theorem, with two lemmas proved and checked and every computable step of Craig's proof rechecked independently. Reflection positivity and the zero-spacing limit stay open, with the same status as the Standard Model's. Two of the framework's own features are what make the completion work: its anomaly-free charges, and its far side having both handednesses. The lemmas are standard lattice reasoning applied to this content, and they should be checked by a lattice specialist before being relied on. A general statement, independent of this framework, is given in companion paper F.

\hypertarget{data-and-code-availability}{%
\subsection{Data and code availability}\label{data-and-code-availability}}

Every number in this paper is produced by a script published with it under the MIT License, together with the particle-to-particle data. The scripts, their saved outputs and the data are archived at \url{https://doi.org/10.5281/zenodo.24329239} (all versions) and kept at github.com/NubbyNubkins/TRPG-Unified-Framework and trpgphysics.org. This paper is archived at \url{https://doi.org/10.5281/zenodo.24330763}, a DOI that always resolves to the latest version. Particle properties are from the Particle Data Group (2026; 2024 in scripts dated before 18 November 2026, which change no headline beyond those listed in Appendix A.6), constants from CODATA 2018, cosmological parameters from Planck 2018, neutrino parameters from JUNO (2026) and NuFIT 6.0, and the cosmological neutrino bound from DESI DR2. \texttt{analysis/recalibrate\_to\_2026\_values\_2026-11-18.py} recomputes every headline number from the 2026 values. The tests added in v1.4.0 (25 November 2026) are listed in Appendix A.7, and the status of each open problem in Appendix A.8.

\hypertarget{acknowledgements}{%
\subsection{Acknowledgements}\label{acknowledgements}}

This framework is a perspective built on the work of the physicists cited below. The lattice construction extended in Appendix B is due to N. Craig; we also acknowledge M. Golterman, A. Kronfeld and Y. Shamir, whose examination of it is reported by OpenAI (2026). Calculations were carried out in Python with NumPy, SciPy and SymPy. Due to medical necessity, AI language-model assistants were used, under the authors' direction, to draft and edit text from the authors' ideas, to carry out and check calculations, and to search the literature; every result was checked by rerunning the published code. The authors reviewed all content and take full responsibility for it. Matthew Gelfman and Melissa Behrends-Gelfman contributed equally to this work.

\hypertarget{references}{%
\subsection{References}\label{references}}

\begin{itemize}
\tightlist
\item
  Abbott, B. P. et al.~(LIGO and Virgo) (2016). Observation of gravitational waves from a binary black hole merger. \emph{Phys. Rev.~Lett.} 116, 061102.
\item
  Abbott, B. P. et al.~(LIGO and Virgo) (2018). GW170817: measurements of neutron star radii and equation of state. \emph{Phys. Rev.~Lett.} 121, 161101.
\item
  Abbott, L. F. and Sikivie, P. (1983). A cosmological bound on the invisible axion. \emph{Phys. Lett. B} 120, 133.
\item
  Abbott, R. et al.~(LIGO and Virgo) (2020). GW190814: gravitational waves from the coalescence of a 23 solar mass black hole with a 2.6 solar mass compact object. \emph{Astrophys. J. Lett.} 896, L44.
\item
  ACT Collaboration (Louis, T. et al.) (2025). The Atacama Cosmology Telescope: DR6 power spectra, likelihoods and \ensuremath{\Lambda}CDM parameters. arXiv:2503.14452.
\item
  Aizenman, M. (1981). Proof of the triviality of \ensuremath{\phi}\ensuremath{^{4}}\_d field theory and some mean-field features of Ising models for d \textgreater{} 4. \emph{Phys. Rev.~Lett.} 47, 1.
\item
  Aizenman, M. and Duminil-Copin, H. (2021). Marginal triviality of the scaling limits of critical 4D Ising and \ensuremath{\phi}\ensuremath{^{4}}\ensuremath{_{4}} models. \emph{Ann. Math.} 194, 163.
\item
  Akhmedov, E. Kh., Berezhiani, Z. G., Mohapatra, R. N. and Senjanović, G. (1993). Planck scale effects on the majoron. \emph{Phys. Lett. B} 299, 90.
\item
  Albert, A. A. (1934). On a certain algebra of quantum mechanics. \emph{Ann. Math.} 35, 65.
\item
  Albright, C. H. and Rodejohann, W. (2009). Comparing trimaximal mixing and its variants with deviations from tri-bimaximal mixing. \emph{Eur. Phys. J. C} 62, 599.
\item
  Alcock, C., Farhi, E. and Olinto, A. (1986). Strange stars. \emph{Astrophys. J.} 310, 261.
\item
  Alford, M., Braby, M., Paris, M. and Reddy, S. (2005). Hybrid stars that masquerade as neutron stars. \emph{Astrophys. J.} 629, 969.
\item
  Alford, M., Rajagopal, K. and Wilczek, F. (1998). QCD at finite baryon density: nucleon droplets and color superconductivity. \emph{Phys. Lett. B} 422, 247.
\item
  An, F. et al.~(JUNO) (2016). Neutrino physics with JUNO. \emph{J. Phys. G} 43, 030401.
\item
  Anderson, P. W. (1963). Plasmons, gauge invariance, and mass. \emph{Phys. Rev.} 130, 439.
\item
  Antusch, S., King, S. F., Malinsky, M. and Spinrath, M. (2010). Quark mixing sum rules and the right unitarity triangle. \emph{Phys. Rev.~D} 81, 033008.
\item
  Arason, H., Castaño, D. J., Keszthelyi, B., Mikaelian, S., Piard, E. J., Ramond, P. and Wright, B. D. (1992). Renormalization-group study of the standard model and its extensions: the standard model. \emph{Phys. Rev.~D} 46, 3945.
\item
  Aspect, A., Grangier, P. and Roger, G. (1982). Experimental realization of Einstein--Podolsky--Rosen--Bohm Gedankenexperiment. \emph{Phys. Rev.~Lett.} 49, 91.
\item
  Atkins, M. and Calmet, X. (2013). Bounds on the nonminimal coupling of the Higgs boson to gravity. \emph{Phys. Rev.~Lett.} 110, 051301.
\item
  ATLAS and CMS Collaborations (2022). A detailed map of Higgs boson interactions by the ATLAS experiment ten years after the discovery; A portrait of the Higgs boson by the CMS experiment ten years after the discovery. \emph{Nature} 607, 52; 607, 60.
\item
  Baez, J. C. and Huerta, J. (2010). The algebra of grand unified theories. \emph{Bull. Amer. Math. Soc.} 47, 483.
\item
  Banks, T. and Seiberg, N. (2011). Symmetries and strings in field theory and gravity. \emph{Phys. Rev.~D} 83, 084019.
\item
  Bardeen, W. A., Hill, C. T. and Lindner, M. (1990). Minimal dynamical symmetry breaking of the standard model. \emph{Phys. Rev.~D} 41, 1647.
\item
  Bekenstein, J. D. (1973). Black holes and entropy. \emph{Phys. Rev.~D} 7, 2333.
\item
  Bell, J. S. (1964). On the Einstein Podolsky Rosen paradox. \emph{Physics} 1, 195.
\item
  Bell, J. S. and Sutherland, D. G. (1968). Current algebra and \ensuremath{\eta} \ensuremath{\to} 3\ensuremath{\pi}. \emph{Nucl. Phys. B} 4, 315.
\item
  Berezinsky, V. and Valle, J. W. F. (1993). The keV majoron as a dark matter particle. \emph{Phys. Lett. B} 318, 360.
\item
  Bethe, H. A. and Bacher, R. F. (1936). Nuclear physics A. Stationary states of nuclei. \emph{Rev.~Mod. Phys.} 8, 82.
\item
  Bezerra-Sobrinho, J. and Medeiros, L. G. (2025). Starobinsky inflation and the latest CMB data: a subtle tension? arXiv:2511.06640.
\item
  Bezrukov, F. and Gorbunov, D. (2012). Distinguishing between R\ensuremath{^{2}}-inflation and Higgs-inflation. \emph{Phys. Lett. B} 713, 365.
\item
  BICEP/Keck Collaboration (2021). Improved constraints on primordial gravitational waves using Planck, WMAP, and BICEP/Keck observations through the 2018 observing season. \emph{Phys. Rev.~Lett.} 127, 151301.
\item
  Blanchet, S. and Di Bari, P. (2007). Flavor effects on leptogenesis predictions. \emph{JCAP} 0703, 018.
\item
  Bodmer, A. R. (1971). Collapsed nuclei. \emph{Phys. Rev.~D} 4, 1601.
\item
  Bohr, N. (1913). On the constitution of atoms and molecules. \emph{Phil. Mag.} 26, 1.
\item
  Bombelli, L., Koul, R. K., Lee, J. and Sorkin, R. D. (1986). Quantum source of entropy for black holes. \emph{Phys. Rev.~D} 34, 373.
\item
  Born, M. (1926). Zur Quantenmechanik der Stoßvorgänge. \emph{Z. Phys.} 37, 863.
\item
  Bose, S. N. (1924). Plancks Gesetz und Lichtquantenhypothese. \emph{Z. Phys.} 26, 178.
\item
  Brannen, C. A. (2006). The lepton masses. Unpublished note, brannenworks.com.
\item
  Bressi, G., Carugno, G., Della Valle, F., Galeazzi, G., Ruoso, G. and Sartori, G. (2011). Testing the neutrality of matter by acoustic means in a spherical resonator. \emph{Phys. Rev.~A} 83, 052101.
\item
  Buchdahl, H. A. (1959). General relativistic fluid spheres. \emph{Phys. Rev.} 116, 1027.
\item
  Buchmüller, W., Di Bari, P. and Plümacher, M. (2005). Leptogenesis for pedestrians. \emph{Ann. Phys.} 315, 305.
\item
  Burkert, V. D., Elouadrhiri, L. and Girod, F. X. (2018). The pressure distribution inside the proton. \emph{Nature} 557, 396.
\item
  Buttazzo, D., Degrassi, G., Giardino, P. P., Giudice, G. F., Sala, F., Salvio, A. and Strumia, A. (2013). Investigating the near-criticality of the Higgs boson. \emph{JHEP} 12, 089.
\item
  Bär, O. and Wiese, U.-J. (2001). Can one see the number of colors? \emph{Nucl. Phys. B} 609, 225.
\item
  Cartan, É. (1894). \emph{Sur la structure des groupes de transformations finis et continus.} Thesis, Paris, Nony.
\item
  Chamseddine, A. H. and Connes, A. (1996). Universal formula for noncommutative geometry actions: unification of gravity and the standard model. \emph{Phys. Rev.~Lett.} 77, 4868.
\item
  Chamseddine, A. H., Connes, A. and Marcolli, M. (2007). Gravity and the standard model with neutrino mixing. \emph{Adv. Theor. Math. Phys.} 11, 991.
\item
  Chikashige, Y., Mohapatra, R. N. and Peccei, R. D. (1981). Are there real Goldstone bosons associated with broken lepton number? \emph{Phys. Lett. B} 98, 265.
\item
  Chiribella, G., D'Ariano, G. M. and Perinotti, P. (2011). Informational derivation of quantum theory. \emph{Phys. Rev.~A} 84, 012311.
\item
  Chodos, A., Jaffe, R. L., Johnson, K., Thorn, C. B. and Weisskopf, V. F. (1974). New extended model of hadrons. \emph{Phys. Rev.~D} 9, 3471.
\item
  Chung, D. J. H., Kolb, E. W. and Riotto, A. (1999). Production of massive particles during reheating. \emph{Phys. Rev.~D} 60, 063504.
\item
  Clauser, J. F. and Freedman, S. J. (1972). Experimental test of local hidden-variable theories. \emph{Phys. Rev.~Lett.} 28, 938.
\item
  Clowe, D., Bradač, M., Gonzalez, A. H., Markevitch, M., Randall, S. W., Jones, C. and Zaritsky, D. (2006). A direct empirical proof of the existence of dark matter. \emph{Astrophys. J.} 648, L109.
\item
  Cohen, A. G., Kaplan, D. B. and Nelson, A. E. (1999). Effective field theory, black holes, and the cosmological constant. \emph{Phys. Rev.~Lett.} 82, 4971.
\item
  Colangelo, G., Lanz, S., Leutwyler, H. and Passemar, E. (2018). Dispersive analysis of \ensuremath{\eta} \ensuremath{\to} 3\ensuremath{\pi}. \emph{Eur. Phys. J. C} 78, 947.
\item
  Craig, N. (2026). A gauge-invariant measure for lattice chiral fermions. arXiv:2610.09093 {[}hep-lat{]}.
\item
  D'Onofrio, M. and Rummukainen, K. (2016). Standard model cross-over on the lattice. \emph{Phys. Rev.~D} 93, 025003.
\item
  Dalla Brida, M., Höllwieser, R., Knechtli, F., Korzec, T., Ramos, A., Sint, S. and Sommer, R. (ALPHA) (2025). The strength of the interaction between quarks and gluons. arXiv:2501.06633.
\item
  Davidson, A. (1979). B \ensuremath{-} L as the fourth color within an SU(2)\textsubscript{L} \ensuremath{\times} U(1)\textsubscript{R} \ensuremath{\times} U(1) model. \emph{Phys. Rev.~D} 20, 776.
\item
  Davidson, S. and Ibarra, A. (2002). A lower bound on the right-handed neutrino mass from leptogenesis. \emph{Phys. Lett. B} 535, 25.
\item
  De Rújula, A., Georgi, H. and Glashow, S. L. (1975). Hadron masses in a gauge theory. \emph{Phys. Rev.~D} 12, 147.
\item
  DeGrand, T., Jaffe, R. L., Johnson, K. and Kiskis, J. (1975). Masses and other parameters of the light hadrons. \emph{Phys. Rev.~D} 12, 2060.
\item
  Deser, S. (1970). Self-interaction and gauge invariance. \emph{Gen.~Rel. Grav.} 1, 9.
\item
  DESI Collaboration (2025). DESI DR2 results II: measurements of baryon acoustic oscillations and cosmological constraints. arXiv:2503.14738.
\item
  Dimopoulos, S. and Georgi, H. (1981). Softly broken supersymmetry and SU(5). \emph{Nucl. Phys. B} 193, 150.
\item
  Dimopoulos, S. and Wilczek, F. (1981). Incomplete multiplets in supersymmetric unified models. NSF-ITP-82-07 (unpublished).
\item
  Dine, M. and Fischler, W. (1983). The not-so-harmless axion. \emph{Phys. Lett. B} 120, 137.
\item
  Dirac, P. A. M. (1937). The cosmological constants. \emph{Nature} 139, 323.
\item
  Dixon, G. M. (1994). \emph{Division Algebras: Octonions, Quaternions, Complex Numbers and the Algebraic Design of Physics.} Kluwer.
\item
  Donoghue, J. F. and Johnson, K. (1980). Pion and nucleon form factors in the bag model. \emph{Phys. Rev.~D} 21, 1975.
\item
  Drago, A., Lavagno, A. and Pagliara, G. (2014). Can very compact and very massive neutron stars both exist? \emph{Phys. Rev.~D} 89, 043014.
\item
  Dubois-Violette, M. and Todorov, I. (2019). Exceptional quantum geometry and particle physics II. \emph{Nucl. Phys. B} 938, 751.
\item
  Dyson, F. W., Eddington, A. S. and Davidson, C. (1920). A determination of the deflection of light by the Sun's gravitational field. \emph{Phil. Trans. R. Soc. A} 220, 291.
\item
  Einstein, A. (1915). Die Feldgleichungen der Gravitation. \emph{Sitzungsber. Preuss. Akad. Wiss.}, 844.
\item
  Einstein, A. (1919). Spielen Gravitationsfelder im Aufbau der materiellen Elementarteilchen eine wesentliche Rolle? \emph{Sitzungsber. Preuss. Akad. Wiss.} 349.
\item
  Einstein, A. (1925). Quantentheorie des einatomigen idealen Gases. Zweite Abhandlung. \emph{Sitzungsber. Preuss. Akad. Wiss.}, 3.
\item
  Einstein, A., Podolsky, B. and Rosen, N. (1935). Can quantum-mechanical description of physical reality be considered complete? \emph{Phys. Rev.} 47, 777.
\item
  Espinosa, J. R., Giudice, G. F., Morgante, E., Riotto, A., Senatore, L., Strumia, A. and Tetradis, N. (2015). The cosmological Higgstory of the vacuum instability. \emph{JHEP} 09, 174.
\item
  Esteban, I., Gonzalez-Garcia, M. C., Maltoni, M., Martinez-Soler, I., Pinheiro, J. P. and Schwetz, T. (2024). NuFit-6.0: updated global analysis of three-flavor neutrino oscillations. \emph{JHEP} 12, 216.
\item
  Farhi, E. and Jaffe, R. L. (1984). Strange matter. \emph{Phys. Rev.~D} 30, 2379.
\item
  Fermi, E. (1934). Versuch einer Theorie der \ensuremath{\beta}-Strahlen. I. \emph{Z. Phys.} 88, 161.
\item
  Feshbach, H. (1958). Unified theory of nuclear reactions. \emph{Ann. Phys.} 5, 357.
\item
  Feynman, R. P., Morinigo, F. B. and Wagner, W. G. (1995). \emph{Feynman Lectures on Gravitation.} Addison-Wesley.
\item
  Fonseca, E. et al.~(2021). Refined mass and geometric measurements of the high-mass PSR J0740+6620. \emph{Astrophys. J. Lett.} 915, L12.
\item
  Foot, R. (1994). A note on Koide's lepton mass relation. arXiv:hep-ph/9402242.
\item
  Fox, P., Pierce, A. and Thomas, S. (2004). Probing a QCD string axion with precision cosmological measurements. arXiv:hep-th/0409059.
\item
  Fraga, E. S., Pisarski, R. D. and Schaffner-Bielich, J. (2001). Small, dense quark stars from perturbative QCD. \emph{Phys. Rev.~D} 63, 121702.
\item
  Freedman, B. A. and McLerran, L. D. (1977). Fermions and gauge vector mesons at finite temperature and density. \emph{Phys. Rev.~D} 16, 1169.
\item
  Freudenthal, H. (1954). Beziehungen der E7 und E8 zur Oktavenebene I. \emph{Proc. Kon. Ned. Akad. Wet. A} 57, 218.
\item
  Friedberg, R. and Lee, T. D. (1977). Fermion-field nontopological solitons. \emph{Phys. Rev.~D} 15, 1694.
\item
  Friedmann, A. (1922). Über die Krümmung des Raumes. \emph{Z. Phys.} 10, 377.
\item
  Fritzsch, H. (1977). Calculating the Cabibbo angle. \emph{Phys. Lett. B} 70, 436.
\item
  Fritzsch, H. and Minkowski, P. (1975). Unified interactions of leptons and hadrons. \emph{Ann. Phys.} 93, 193.
\item
  Froggatt, C. D. and Nielsen, H. B. (1979). Hierarchy of quark masses, Cabibbo angles and CP violation. \emph{Nucl. Phys. B} 147, 277.
\item
  Fröhlich, J. (1982). On the triviality of \ensuremath{\lambda}\ensuremath{\phi}\ensuremath{^{4}}\_d theories and the approach to the critical point in d \ensuremath{\ge} 4 dimensions. \emph{Nucl. Phys. B} 200, 281.
\item
  Fukugita, M. and Yanagida, T. (1986). Baryogenesis without grand unification. \emph{Phys. Lett. B} 174, 45.
\item
  Furey, C. (2016). Standard model physics from an algebra? PhD thesis, University of Waterloo. arXiv:1611.09182.
\item
  Furry, W. H. (1939). On transition probabilities in double beta-disintegration. \emph{Phys. Rev.} 56, 1184.
\item
  Garcia-Cely, C. and Heeck, J. (2017). Neutrino lines from majoron dark matter. \emph{JHEP} 05, 102.
\item
  Gatto, R., Sartori, G. and Tonin, M. (1968). Weak self-masses, Cabibbo angle, and broken SU(2) \ensuremath{\times} SU(2). \emph{Phys. Lett. B} 28, 128.
\item
  Gell-Mann, M. and Low, F. E. (1954). Quantum electrodynamics at small distances. \emph{Phys. Rev.} 95, 1300.
\item
  Gell-Mann, M., Ramond, P. and Slansky, R. (1979). Complex spinors and unified theories. In \emph{Supergravity}, North-Holland, 315.
\item
  Geng, C. Q. and Marshak, R. E. (1989). Uniqueness of quark and lepton representations in the standard model from the anomalies viewpoint. \emph{Phys. Rev.~D} 39, 693.
\item
  Georgi, H. (1975). The state of the art: gauge theories. \emph{AIP Conf. Proc.} 23, 575.
\item
  Georgi, H. and Glashow, S. L. (1974). Unity of all elementary-particle forces. \emph{Phys. Rev.~Lett.} 32, 438.
\item
  Georgi, H. and Jarlskog, C. (1979). A new lepton--quark mass relation in a unified theory. \emph{Phys. Lett. B} 86, 297.
\item
  Georgi, H., Quinn, H. R. and Weinberg, S. (1974). Hierarchy of interactions in unified gauge theories. \emph{Phys. Rev.~Lett.} 33, 451.
\item
  Ginsparg, P. H. and Wilson, K. G. (1982). A remnant of chiral symmetry on the lattice. \emph{Phys. Rev.~D} 25, 2649.
\item
  Giudice, G. F., Kolb, E. W. and Riotto, A. (2001). Largest temperature of the radiation era and its cosmological implications. \emph{Phys. Rev.~D} 64, 023508.
\item
  Glashow, S. L. (1961). Partial-symmetries of weak interactions. \emph{Nucl. Phys.} 22, 579.
\item
  Glashow, S. L. (1980). The future of elementary particle physics. \emph{NATO Sci. Ser. B} 61, 687.
\item
  Gleason, A. M. (1957). Measures on the closed subspaces of a Hilbert space. \emph{J. Math. Mech.} 6, 885.
\item
  Gordon, W. (1923). Zur Lichtfortpflanzung nach der Relativitätstheorie. \emph{Ann. Phys.} 377, 421.
\item
  Gross, D. J. and Wilczek, F. (1973). Ultraviolet behavior of non-abelian gauge theories. \emph{Phys. Rev.~Lett.} 30, 1343.
\item
  Günaydin, M. and Gürsey, F. (1973). Quark structure and octonions. \emph{J. Math. Phys.} 14, 1651.
\item
  Gürsey, F., Ramond, P. and Sikivie, P. (1976). A universal gauge theory model based on E6. \emph{Phys. Lett. B} 60, 177.
\item
  Harari, H. (1979). A schematic model of quarks and leptons. \emph{Phys. Lett. B} 86, 83.
\item
  Hardy, L. (2001). Quantum theory from five reasonable axioms. arXiv:quant-ph/0101012.
\item
  Harlow, D. and Ooguri, H. (2019). Constraints on symmetries from holography. \emph{Phys. Rev.~Lett.} 122, 191601.
\item
  Harrison, P. F., Perkins, D. H. and Scott, W. G. (2002). Tri-bimaximal mixing and the neutrino oscillation data. \emph{Phys. Lett. B} 530, 167.
\item
  Hasenfratz, P., Laliena, V. and Niedermayer, F. (1998). The index theorem in QCD with a finite cut-off. \emph{Phys. Lett. B} 427, 125.
\item
  Hawking, S. W. (1975). Particle creation by black holes. \emph{Commun. Math. Phys.} 43, 199.
\item
  Heisenberg, W. (1927). Über den anschaulichen Inhalt der quantentheoretischen Kinematik und Mechanik. \emph{Z. Phys.} 43, 172.
\item
  Henneaux, M. and Teitelboim, C. (1989). The cosmological constant and general covariance. \emph{Phys. Lett. B} 222, 195.
\item
  Hernández, P., Jansen, K. and Lüscher, M. (1999). Locality properties of Neuberger's lattice Dirac operator. \emph{Nucl. Phys. B} 552, 363.
\item
  Herranen, M., Markkanen, T., Nurmi, S. and Rajantie, A. (2014). Spacetime curvature and the Higgs stability during inflation. \emph{Phys. Rev.~Lett.} 113, 211102.
\item
  Hertzberg, M. P., Tegmark, M. and Wilczek, F. (2008). Axion cosmology and the energy scale of inflation. \emph{Phys. Rev.~D} 78, 083507.
\item
  Hill, C. T. (1981). Quark and lepton masses from renormalization-group fixed points. \emph{Phys. Rev.~D} 24, 691.
\item
  Hinderer, T. (2008). Tidal Love numbers of neutron stars. \emph{Astrophys. J.} 677, 1216.
\item
  Hoffman, Y., Pomarède, D., Tully, R. B. and Courtois, H. M. (2017). The dipole repeller. \emph{Nature Astron.} 1, 0036.
\item
  Hoof, S. and Marsh, D. J. E. (2024). Black hole superradiance constraints on ultralight bosons, revisited. arXiv:2406.10337.
\item
  't Hooft, G. (1976). Symmetry breaking through Bell--Jackiw anomalies. \emph{Phys. Rev.~Lett.} 37, 8.
\item
  Hu, W., Barkana, R. and Gruzinov, A. (2000). Fuzzy cold dark matter: the wave properties of ultralight particles. \emph{Phys. Rev.~Lett.} 85, 1158.
\item
  Hui, L., Ostriker, J. P., Tremaine, S. and Witten, E. (2017). Ultralight scalars as cosmological dark matter. \emph{Phys. Rev.~D} 95, 043541.
\item
  Hurwitz, A. (1898). Über die Composition der quadratischen Formen von beliebig vielen Variablen. \emph{Nachr. Ges. Wiss. Göttingen}, 309.
\item
  Ibáñez, L. E. and Ross, G. G. (1991). Discrete gauge symmetry anomalies. \emph{Phys. Lett. B} 260, 291.
\item
  Jacobson, T. (1995). Thermodynamics of spacetime: the Einstein equation of state. \emph{Phys. Rev.~Lett.} 75, 1260.
\item
  Ji, X. (1995). QCD analysis of the mass structure of the nucleon. \emph{Phys. Rev.~Lett.} 74, 1071.
\item
  Jordan, P. (1933). Über Verallgemeinerungsmöglichkeiten des Formalismus der Quantenmechanik. \emph{Nachr. Ges. Wiss. Göttingen} 1933, 209.
\item
  Jordan, P., von Neumann, J. and Wigner, E. (1934). On an algebraic generalization of the quantum mechanical formalism. \emph{Ann. Math.} 35, 29.
\item
  JUNO Collaboration (2024). Potential to identify neutrino mass ordering with reactor antineutrinos at JUNO. arXiv:2405.18008.
\item
  JUNO Collaboration (2025). First measurement of reactor neutrino oscillations at JUNO. arXiv:2511.14593.
\item
  JUNO Collaboration (2026). Measurement of neutrino oscillation parameters at JUNO and indication of the neutrino mass ordering. arXiv:2609.37895.
\item
  Kadastik, M., Kannike, K. and Raidal, M. (2010). Matter parity as the origin of scalar dark matter. \emph{Phys. Rev.~D} 81, 015002.
\item
  Kadoh, D. and Kikukawa, Y. (2008). A construction of the Glashow--Weinberg--Salam model on the lattice with exact gauge invariance. \emph{JHEP} 05, 095.
\item
  KamLAND-Zen Collaboration (2024). Search for Majorana neutrinos with the complete KamLAND-Zen dataset. arXiv:2406.11438.
\item
  Kaplunovsky, V. S. (1988). One-loop threshold effects in string unification. \emph{Nucl. Phys. B} 307, 145.
\item
  Karachentsev, I. D., Kashibadze, O. G., Makarov, D. I. and Tully, R. B. (2009). The Hubble flow around the Local Group. \emph{Mon. Not. R. Astron. Soc.} 393, 1265.
\item
  Kartavtsev, A. (2011). A remark on the Koide relation for quarks. arXiv preprint.
\item
  Killing, W. (1888). Die Zusammensetzung der stetigen endlichen Transformationsgruppen. \emph{Math. Ann.} 31, 252.
\item
  Kobayashi, M. and Maskawa, T. (1973). CP-violation in the renormalizable theory of weak interaction. \emph{Prog. Theor. Phys.} 49, 652.
\item
  Koide, Y. (1983). A fermion--boson composite model of quarks and leptons. \emph{Phys. Lett. B} 120, 161.
\item
  Krasnov, K. (2019). SO(9) characterisation of the Standard Model gauge group. arXiv:1912.11282; \emph{J. Math. Phys.} 62, 021703 (2021).
\item
  Krauss, L. M. and Wilczek, F. (1989). Discrete gauge symmetry in continuum theories. \emph{Phys. Rev.~Lett.} 62, 1221.
\item
  Lazarides, G. and Shafi, Q. (1991). Origin of matter in the inflationary cosmology. \emph{Phys. Lett. B} 258, 305.
\item
  Lee, T. D. (1973). A theory of spontaneous T violation. \emph{Phys. Rev.~D} 8, 1226.
\item
  Lee, T. D. and Yang, C. N. (1956). Question of parity conservation in weak interactions. \emph{Phys. Rev.} 104, 254.
\item
  Lemaître, G. (1927). Un univers homogène de masse constante et de rayon croissant. \emph{Ann. Soc. Sci. Bruxelles} A47, 49.
\item
  Leutwyler, H. (1996). The ratios of the light quark masses. \emph{Phys. Lett. B} 378, 313.
\item
  Li, M. (2004). A model of holographic dark energy. \emph{Phys. Lett. B} 603, 1.
\item
  Loos, O. (1977). \emph{Bounded Symmetric Domains and Jordan Pairs.} University of California, Irvine.
\item
  Luo, M., Wang, H. and Xiao, Y. (2003). Two-loop renormalization group equations in general gauge field theories. \emph{Phys. Rev.~D} 67, 065019.
\item
  Lüscher, M. (1998). Exact chiral symmetry on the lattice and the Ginsparg--Wilson relation. \emph{Phys. Lett. B} 428, 342.
\item
  Lüscher, M. (1999). Abelian chiral gauge theories on the lattice with exact gauge invariance. \emph{Nucl. Phys. B} 549, 295.
\item
  Machacek, M. E. and Vaughn, M. T. (1984). Two-loop renormalization group equations in a general quantum field theory. II. Yukawa couplings. \emph{Nucl. Phys. B} 236, 221.
\item
  Majorana, E. (1937). Teoria simmetrica dell'elettrone e del positrone. \emph{Nuovo Cimento} 14, 171.
\item
  Maldacena, J. and Susskind, L. (2013). Cool horizons for entangled black holes. \emph{Fortschr. Phys.} 61, 781.
\item
  Mannarelli, M., Rajagopal, K. and Sharma, R. (2007). The rigidity of crystalline color superconducting quark matter. \emph{Phys. Rev.~D} 76, 074026.
\item
  Marciano, W. J. and Sirlin, A. (1993). Radiative corrections to \ensuremath{\pi}\textsubscript{\ensuremath{\ell}2} decays. \emph{Phys. Rev.~Lett.} 71, 3629.
\item
  Marshak, R. E. and Mohapatra, R. N. (1980). Quark--lepton symmetry and B \ensuremath{-} L as the U(1) generator of the electroweak symmetry group. \emph{Phys. Lett. B} 91, 222.
\item
  Masiero, A., Nanopoulos, D. V., Tamvakis, K. and Yanagida, T. (1982). Naturally massless Higgs doublets in supersymmetric SU(5). \emph{Phys. Lett. B} 115, 380.
\item
  McDonough, E. and Ferreira, E. G. M. (2025). The spectrum of n\_s constraints from DESI and CMB data. arXiv:2512.05108.
\item
  McGaugh, S. S., Lelli, F. and Schombert, J. M. (2016). Radial acceleration relation in rotationally supported galaxies. \emph{Phys. Rev.~Lett.} 117, 201101.
\item
  Meitner, L. and Frisch, O. R. (1939). Disintegration of uranium by neutrons: a new type of nuclear reaction. \emph{Nature} 143, 239.
\item
  Michel, L. and Radicati, L. A. (1971). Properties of the breaking of hadronic internal symmetry. \emph{Ann. Phys.} 66, 758.
\item
  MICROSCOPE Collaboration (2022). MICROSCOPE mission: final results of the test of the equivalence principle. \emph{Phys. Rev.~Lett.} 129, 121102.
\item
  Milgrom, M. (1983). A modification of the Newtonian dynamics as a possible alternative to the hidden mass hypothesis. \emph{Astrophys. J.} 270, 365.
\item
  Minahan, J. A., Ramond, P. and Warner, R. C. (1990). Comment on anomaly cancellation in the standard model. \emph{Phys. Rev.~D} 41, 715.
\item
  Minkowski, P. (1977). \ensuremath{\mu} \ensuremath{\to} e\ensuremath{\gamma} at a rate of one out of 10\textsuperscript{9} muon decays? \emph{Phys. Lett. B} 67, 421.
\item
  Mohapatra, R. N. and Senjanović, G. (1980). Neutrino mass and spontaneous parity nonconservation. \emph{Phys. Rev.~Lett.} 44, 912.
\item
  Morningstar, C. J. and Peardon, M. (1999). The glueball spectrum from an anisotropic lattice study. \emph{Phys. Rev.~D} 60, 034509.
\item
  Mukhanov, V. F. and Chibisov, G. V. (1981). Quantum fluctuations and a nonsingular universe. \emph{JETP Lett.} 33, 532.
\item
  Narayanan, R. and Neuberger, H. (1994). Chiral determinant as an overlap of two vacua. \emph{Nucl. Phys. B} 412, 574.
\item
  Nardi, E., Nir, Y., Roulet, E. and Racker, J. (2006). The importance of flavor in leptogenesis. \emph{JHEP} 0601, 164.
\item
  Navas, S. et al.~(Particle Data Group) (2024). Review of particle physics. \emph{Phys. Rev.~D} 110, 030001.
\item
  Neuberger, H. (1998). Exactly massless quarks on the lattice. \emph{Phys. Lett. B} 417, 141.
\item
  Nielsen, H. B. and Ninomiya, M. (1981). Absence of neutrinos on a lattice. \emph{Nucl. Phys. B} 185, 20.
\item
  Noether, E. (1918). Invariante Variationsprobleme. \emph{Nachr. Ges. Wiss. Göttingen}, 235.
\item
  OpenAI (2026). A step toward defining the Standard Model on a lattice. openai.com/index/step-toward-defining-the-standard-model (15 November 2026).
\item
  Oppenheimer, J. R. and Volkoff, G. M. (1939). On massive neutron cores. \emph{Phys. Rev.} 55, 374.
\item
  Osterwalder, K. and Seiler, E. (1978). Gauge field theories on a lattice. \emph{Ann. Phys.} 110, 440.
\item
  Page, D. N. and Wootters, W. K. (1983). Evolution without evolution: dynamics described by stationary observables. \emph{Phys. Rev.~D} 27, 2885.
\item
  Pati, J. C. and Salam, A. (1974). Lepton number as the fourth ``color''. \emph{Phys. Rev.~D} 10, 275.
\item
  Pauli, W. (1925). Über den Zusammenhang des Abschlusses der Elektronengruppen im Atom mit der Komplexstruktur der Spektren. \emph{Z. Phys.} 31, 765.
\item
  Pendleton, B. and Ross, G. G. (1981). Mass and mixing angle predictions from infrared fixed points. \emph{Phys. Lett. B} 98, 291.
\item
  Penrose, R. (1979). Singularities and time-asymmetry. In \emph{General Relativity: An Einstein Centenary Survey}, ed.~Hawking and Israel, Cambridge University Press.
\item
  Penrose, R. and Rindler, W. (1984). \emph{Spinors and Space-Time}, Vol. 1. Cambridge University Press.
\item
  Planck Collaboration (2020). Planck 2018 results. VI. Cosmological parameters. \emph{Astron. Astrophys.} 641, A6.
\item
  Politzer, H. D. (1973). Reliable perturbative results for strong interactions? \emph{Phys. Rev.~Lett.} 30, 1346.
\item
  Preskill, J. P. (1979). Cosmological production of superheavy magnetic monopoles. \emph{Phys. Rev.~Lett.} 43, 1365.
\item
  Preskill, J., Wise, M. B. and Wilczek, F. (1983). Cosmology of the invisible axion. \emph{Phys. Lett. B} 120, 127.
\item
  Rapp, R., Schäfer, T., Shuryak, E. and Velkovsky, M. (1998). Diquark Bose condensates in high density matter and instantons. \emph{Phys. Rev.~Lett.} 81, 53.
\item
  Renou, M.-O. et al.~(2021). Quantum theory based on real numbers can be experimentally falsified. \emph{Nature} 600, 625.
\item
  Reuter, M. (1998). Nonperturbative evolution equation for quantum gravity. \emph{Phys. Rev.~D} 57, 971.
\item
  Rivero, A. (2011). A new Koide tuple: strange--charm--bottom. arXiv:1111.7232.
\item
  Rogers, K. K. and Peiris, H. V. (2021). Strong bound on canonical ultralight axion dark matter from the Lyman-alpha forest. \emph{Phys. Rev.~Lett.} 126, 071302.
\item
  Romani, R. W., Kandel, D., Filippenko, A. V., Brink, T. G. and Zheng, W. (2022). PSR J0952\ensuremath{-}0607: the fastest and heaviest known Galactic neutron star. \emph{Astrophys. J. Lett.} 934, L17.
\item
  Rothstein, I. Z., Babu, K. S. and Seckel, D. (1993). Planck scale symmetry breaking and majoron physics. \emph{Nucl. Phys. B} 403, 725.
\item
  Ryu, S. and Takayanagi, T. (2006). Holographic derivation of entanglement entropy from AdS/CFT. \emph{Phys. Rev.~Lett.} 96, 181602.
\item
  Sakharov, A. D. (1967). Vacuum quantum fluctuations in curved space and the theory of gravitation. \emph{Dokl. Akad. Nauk SSSR} 177, 70.
\item
  Sakharov, A. D. and Zeldovich, Ya. B. (1966). Quark structure and masses of strongly interacting particles. \emph{Yad. Fiz.} 4, 395 (\emph{Sov. J. Nucl. Phys.} 4, 283, 1967).
\item
  Salam, A. (1968). Weak and electromagnetic interactions. In \emph{Elementary Particle Theory} (Nobel Symposium 8), ed.~N. Svartholm, Almqvist and Wiksell, 367.
\item
  Schrödinger, E. (1935). Discussion of probability relations between separated systems. \emph{Proc. Camb. Phil. Soc.} 31, 555.
\item
  Schur, I. (1917). Über Potenzreihen, die im Innern des Einheitskreises beschränkt sind. \emph{J. Reine Angew. Math.} 147, 205.
\item
  Shannon, C. E. (1948). A mathematical theory of communication. \emph{Bell Syst. Tech. J.} 27, 379.
\item
  Shapiro, I. I. (1964). Fourth test of general relativity. \emph{Phys. Rev.~Lett.} 13, 789.
\item
  Shaposhnikov, M. and Wetterich, C. (2010). Asymptotic safety of gravity and the Higgs boson mass. \emph{Phys. Lett. B} 683, 196.
\item
  Shifman, M. A., Vainshtein, A. I. and Zakharov, V. I. (1978). Remarks on Higgs-boson interactions with nucleons. \emph{Phys. Lett. B} 78, 443.
\item
  Shupe, M. A. (1979). A composite model of leptons and quarks. \emph{Phys. Lett. B} 86, 87.
\item
  Slansky, R. (1981). Group theory for unified model building. \emph{Phys. Rep.} 79, 1.
\item
  Srednicki, M. (1993). Entropy and area. \emph{Phys. Rev.~Lett.} 71, 666.
\item
  Starobinsky, A. A. (1980). A new type of isotropic cosmological models without singularity. \emph{Phys. Lett. B} 91, 99.
\item
  Stueckelberg, E. C. G. (1960). Quantum theory in real Hilbert space. \emph{Helv. Phys. Acta} 33, 727.
\item
  Sumino, Y. (2009). Family gauge symmetry as an origin of Koide's mass formula and charged lepton spectrum. \emph{Phys. Lett. B} 671, 477.
\item
  Sutherland, D. G. (1966). Current algebra and the decay \ensuremath{\eta} \ensuremath{\to} 3\ensuremath{\pi}. \emph{Phys. Lett.} 23, 384.
\item
  Suzuki, H. (2000). Real representation in chiral gauge theories on the lattice. \emph{JHEP} 10, 039.
\item
  Takenaka, A. et al.~(Super-Kamiokande Collaboration) (2020). Search for proton decay via p \ensuremath{\to} e\textsuperscript{+}\ensuremath{\pi}\textsuperscript{0} and p \ensuremath{\to} \ensuremath{\mu}\textsuperscript{+}\ensuremath{\pi}\textsuperscript{0} with an enlarged fiducial volume in Super-Kamiokande I--IV. \emph{Phys. Rev.~D} 102, 112011.
\item
  Théberge, S., Thomas, A. W. and Miller, G. A. (1980). Pionic corrections to the MIT bag model: the (3,3) resonance. \emph{Phys. Rev.~D} 22, 2838.
\item
  Tiesinga, E., Mohr, P. J., Newell, D. B. and Taylor, B. N. (2021). CODATA recommended values of the fundamental physical constants: 2018. \emph{Rev.~Mod. Phys.} 93, 025010.
\item
  Todorov, I. and Dubois-Violette, M. (2018). Deducing the symmetry of the Standard Model from the automorphism and structure groups of the exceptional Jordan algebra. \emph{Int. J. Mod. Phys. A} 33, 1850118.
\item
  Tolman, R. C. (1939). Static solutions of Einstein's field equations for spheres of fluid. \emph{Phys. Rev.} 55, 364.
\item
  Vafa, C. and Witten, E. (1984). Parity conservation in quantum chromodynamics. \emph{Phys. Rev.~Lett.} 53, 535.
\item
  van der Bij, J. J., van Dam, H. and Ng, Y. J. (1982). The exchange of massless spin-two particles. \emph{Physica A} 116, 307.
\item
  Van Raamsdonk, M. (2010). Building up spacetime with quantum entanglement. \emph{Gen.~Rel. Grav.} 42, 2323.
\item
  Veltman, M. (1977). Limit on mass differences in the Weinberg model. \emph{Nucl. Phys. B} 123, 89.
\item
  Veltman, M. (1981). The infrared--ultraviolet connection. \emph{Acta Phys. Polon. B} 12, 437.
\item
  Verhulst, P.-F. (1838). Notice sur la loi que la population suit dans son accroissement. \emph{Corresp. Math. Phys.} 10, 113.
\item
  Verlinde, E. (2017). Emergent gravity and the dark universe. \emph{SciPost Phys.} 2, 016.
\item
  van Ritbergen, T., Schellekens, A. N. and Vermaseren, J. A. M. (1999). Group theory factors for Feynman diagrams. \emph{Int. J. Mod. Phys. A} 14, 41.
\item
  Vilenkin, A. (1985). Classical and quantum cosmology of the Starobinsky inflationary model. \emph{Phys. Rev.~D} 32, 2511.
\item
  Wagner, T. A., Schlamminger, S., Gundlach, J. H. and Adelberger, E. G. (2012). Torsion-balance tests of the weak equivalence principle. \emph{Class. Quantum Grav.} 29, 184002.
\item
  Wang, M., Huang, W. J., Kondev, F. G., Audi, G. and Naimi, S. (2021). The AME 2020 atomic mass evaluation (II). Tables, graphs and references. \emph{Chin. Phys. C} 45, 030003.
\item
  Weinberg, S. (1964). Photons and gravitons in S-matrix theory. \emph{Phys. Rev.} 135, B1049.
\item
  Weinberg, S. (1967). A model of leptons. \emph{Phys. Rev.~Lett.} 19, 1264.
\item
  Weinberg, S. (1977). The problem of mass. \emph{Trans. N. Y. Acad. Sci.} 38, 185.
\item
  Weinberg, S. (1979). Ultraviolet divergences in quantum theories of gravitation. In \emph{General Relativity: An Einstein Centenary Survey}, Cambridge University Press, 790.
\item
  Weinberg, S. (1979b). Baryon- and lepton-nonconserving processes. \emph{Phys. Rev.~Lett.} 43, 1566.
\item
  Weinberg, S. (1989). The cosmological constant problem. \emph{Rev.~Mod. Phys.} 61, 1.
\item
  Weizsäcker, C. F. von (1935). Zur Theorie der Kernmassen. \emph{Z. Phys.} 96, 431.
\item
  Weyl, H. (1929). Elektron und Gravitation. I. \emph{Z. Phys.} 56, 330.
\item
  Wilczek, F. (2001). Scaling Mount Planck I: a view from the bottom. \emph{Phys. Today} 54(6), 12.
\item
  Wilczek, F. and Zee, A. (1979). Operator analysis of nucleon decay. \emph{Phys. Rev.~Lett.} 43, 1571.
\item
  Wilson, K. G. (1974). Confinement of quarks. \emph{Phys. Rev.~D} 10, 2445.
\item
  Witten, E. (1983). Fermion quantum numbers in Kaluza--Klein theory. In \emph{Shelter Island II}, MIT Press (1985), 227.
\item
  Witten, E. (1984). Cosmic separation of phases. \emph{Phys. Rev.~D} 30, 272.
\item
  Wu, C. S., Ambler, E., Hayward, R. W., Hoppes, D. D. and Hudson, R. P. (1957). Experimental test of parity conservation in beta decay. \emph{Phys. Rev.} 105, 1413.
\item
  Xing, Z.-z. and Zhou, S. (2007). A partial \ensuremath{\mu}--\ensuremath{\tau} symmetry and its prediction for leptonic CP violation. \emph{Phys. Lett. B} 653, 278.
\item
  Yanagida, T. (1979). Horizontal gauge symmetry and masses of neutrinos. \emph{Proc. Workshop on Unified Theory and Baryon Number in the Universe}, KEK, 95.
\item
  Yang, C. N. and Mills, R. L. (1954). Conservation of isotopic spin and isotopic gauge invariance. \emph{Phys. Rev.} 96, 191.
\item
  Yang, Y.-B., Liang, J., Bi, Y.-J., Chen, Y., Draper, T., Liu, K.-F. and Liu, Z. (2018). Proton mass decomposition from the QCD energy momentum tensor. \emph{Phys. Rev.~Lett.} 121, 212001.
\item
  Yukawa, H. (1935). On the interaction of elementary particles. I. \emph{Proc. Phys.-Math. Soc. Japan} 17, 48.
\item
  Zee, A. (1979). Broken-symmetric theory of gravity. \emph{Phys. Rev.~Lett.} 42, 417.
\item
  Zel'dovich, Ya. B. (1967). Cosmological constant and elementary particles. \emph{JETP Lett.} 6, 316.
\item
  Zeldovich, Ya. B. and Khlopov, M. Yu. (1978). On the concentration of relic magnetic monopoles in the universe. \emph{Phys. Lett. B} 79, 239.
\item
  Zu, L., Feng, L., Yuan, Q. and Fan, Y.-Z. (2020). Stringent constraints on the light boson model with supermassive black hole spin measurements. arXiv:2007.03222.
\end{itemize}

\end{document}
